% Leçon 10 — Grammaire : Phrases comparatives avec…不如… ; A比B 形容词+一点ou 多了 ; A比B更/ 还+形容词 ; 为了 ; 不但…而且… — Chinois 1ère A — Cours signé OnBuch+
% Programme officiel MINESEC. Compiler avec : tectonic ZHA10-grammaire-phrases-comparatives-avec-ab-ou-ab.tex
\newcommand{\DOCMATIERE}{Chinois}
\newcommand{\DOCNIVEAU}{1ère A}
\newcommand{\DOCMODULE}{Module 2 : Environnement et santé}
\newcommand{\DOCLECON}{ZHA10}
\newcommand{\DOCTITRE}{Leçon 10 — Grammaire : Phrases comparatives avec…不如… ; A比B 形容词+一点ou 多了 ; A比B更/ 还+形容词 ; 为了 ; 不但…而且…}
% OnBuch+ — préambule « cours signés » ENRICHI (compilé avec tectonic / XeLaTeX)
% Système visuel « Pop » d'OnBuch+ : fond crème (jamais blanc), bordures encre
% épaisses, ombres « dures » décalées, titres Archivo Black en capitales.
% Version MATHÉMATIQUES : graphes (pgfplots/tikz), boîtes pédagogiques
% (définition, propriété, méthode, attention, le savais-tu, exercice résolu,
% exercice, TP, bilan), page de garde et signature OnBuch+ ancrée en bas.
%
% Macros attendues dans main.tex AVANT \input{preamble} :
%   \DOCMATIERE  \DOCNIVEAU  \DOCMODULE  \DOCLECON (numéro)  \DOCTITRE
\documentclass[11pt]{article}

\usepackage{fontspec}
\usepackage[french]{babel} % français = langue principale ; le chinois via \zh{..} / textebox (xeCJK)
\usepackage{xeCJK}
\usepackage{pifont} % \ding{51} \ding{55} utilisés par les modèles
\usepackage[a4paper,margin=2cm,top=3.4cm,bottom=2.8cm,headheight=1.4cm,footskip=1.3cm]{geometry}
\usepackage{amsmath,amssymb,amsfonts}
\usepackage{mathtools} % \xmapsto, \coloneqq... souvent utilisés par les modèles
\usepackage{cancel} % simplifications barrées (bilans de forces)
% \abs{z} (module, valeur absolue) et \norm{u} (norme), utilisés par les modèles
\providecommand{\abs}{}\let\abs\relax\DeclarePairedDelimiter{\abs}{\lvert}{\rvert}
\providecommand{\norm}{}\let\norm\relax\DeclarePairedDelimiter{\norm}{\lVert}{\rVert}
\usepackage[table]{xcolor} % [table] : fournit aussi \rowcolors (alternance de lignes)
\usepackage{pagecolor}
\usepackage{tikz}
\usetikzlibrary{arrows.meta,positioning,calc,shapes.geometric,shapes.misc,shapes.arrows,decorations.pathmorphing,decorations.pathreplacing,decorations.markings,patterns,shadows,mindmap,trees,fit,backgrounds,matrix,intersections,angles,quotes}
\usepackage{pgfplots}
\usepackage[version=4]{mhchem} % équations nucléaires \ce{^{235}_{92}U}
\usepackage[european,siunitx]{circuitikz} % schémas électriques
\pgfplotsset{compat=1.18}
\usepgfplotslibrary{fillbetween,groupplots} % aires entre courbes (calcul intégral)
\usepackage{enumitem}
\usepackage{fancyhdr}
% [most] charge toutes les bibliothèques tcolorbox (listings, minted, raster,
% documentation...) et gonfle inutilement la mémoire TeX sur un document long
% (~300 boîtes) : on ne charge que ce qui est réellement utilisé.
\usepackage[skins,breakable]{tcolorbox}
\usepackage{graphicx}
\usepackage{colortbl}
\usepackage{booktabs}
\usepackage{tabularx}
\usepackage{diagbox} % case d'en-tête diagonale des tableaux à double entrée
% Colonnes extensibles centrée / gauche / droite (C, L, R), souvent utilisées par les modèles
\newcolumntype{C}{>{\centering\arraybackslash}X}
\newcolumntype{L}{>{\raggedright\arraybackslash}X}
\newcolumntype{R}{>{\raggedleft\arraybackslash}X}
\usepackage{array}
\usepackage{multicol}
\usepackage{multirow}
\usepackage{siunitx}
% parse-numbers=false : accepte aussi \SI{2,5 \times 10^{-3}}{...}
\sisetup{locale=FR,output-decimal-marker={,},group-minimum-digits=4,parse-numbers=false}
\DeclareSIUnit\ohm{\ensuremath{\symup{\Omega}}} % U+2126 absent de la police
\DeclareSIUnit\division{div} % réglages d'oscilloscope (ms/div, V/div)
\DeclareSIUnit\tr{tr} % tours (vitesse de rotation, ex. \SI{1500}{\tr\per\minute})
\DeclareSIUnit\molar{M} % concentration molaire (ex. \SI{3}{\molar}, \SI{1}{\micro\molar})
\DeclareSIUnit\an{an} % année (le modèle confond parfois avec \year, un compteur TeX) — ex. \SI{5}{\kilo\meter\per\an}
\DeclareSIUnit\FCFA{FCFA} % franc CFA (ex. \SI{500}{\FCFA\per\kilo\gram})
\DeclareSIUnit\degreeBrix{\textdegree Brix} % teneur en sucre d'un sirop/jus (réfractométrie)
\DeclareSIUnit\month{mois} % le modèle utilise parfois \month, un compteur TeX, comme unité de durée
\usepackage{qrcode}
% \degree : commande fréquemment utilisée par les modèles pour un angle ou une
% température en dehors de \SI{}{} (non fournie par les paquets chargés).
\providecommand{\degree}{^{\circ}}
% \qed (fin de démonstration), utilisé par les modèles sans amsthm
\providecommand{\qed}{\hfill$\square$}
% \barre : double barre des valeurs interdites dans un tableau de variations (tkz-tab)
\providecommand{\barre}{\ensuremath{\|}}
% environnement proof (amsthm non chargé) utilisé par les modèles
\ifdefined\proof\else\newenvironment{proof}[1][Démonstration]{\par\noindent\textit{#1.}\ }{\qed\par}\fi
\usepackage{lastpage}
\usepackage{hyperref}
\hypersetup{hidelinks}

% ============================================================
% Palette « Pop » OnBuch+ — mêmes hex que la classe P (plus_ui.dart)
% ============================================================
\definecolor{poporange}{HTML}{F59321}
\definecolor{popdark}{HTML}{E07A0C}
\definecolor{popcream}{HTML}{FFF6E8}
\definecolor{popink}{HTML}{1C1714}
\definecolor{poppurple}{HTML}{7A5AE0}
\definecolor{popgreen}{HTML}{2E9E6A}
\definecolor{poppink}{HTML}{E0457B}
\definecolor{popblue}{HTML}{2F7DE1}
\definecolor{popgold}{HTML}{C98A12}
\definecolor{popmuted}{HTML}{8A7F76}
\definecolor{popline}{HTML}{EADFD2}
\definecolor{popgray}{HTML}{8A7F76}
\colorlet{popgrey}{popgray}
% Alias tolérants (le modèle nomme parfois les couleurs différemment)
\colorlet{popwhite}{white}
\colorlet{popblack}{popink}
\colorlet{popred}{poppink}
\colorlet{popyellow}{popgold}
% Teintes claires (fonds de tableaux/encadrés) — le modèle nomme parfois
% les couleurs avec un suffixe L (ex. popblueL) sans les avoir définies.
\colorlet{poporangeL}{poporange!20}
\colorlet{popdarkL}{popdark!20}
\colorlet{poppurpleL}{poppurple!20}
\colorlet{popgreenL}{popgreen!20}
\colorlet{poppinkL}{poppink!20}
\colorlet{popblueL}{popblue!20}
\colorlet{popgoldL}{popgold!20}
\colorlet{popredL}{popred!20}
\colorlet{popgrayL}{popgray!20}
% Teintes claires pour fonds de tableaux / zones
\colorlet{popblueL}{popblue!10!white}
\colorlet{poporangeL}{poporange!14!white}
\colorlet{popgreenL}{popgreen!12!white}
\colorlet{poppurpleL}{poppurple!12!white}
\colorlet{poppinkL}{poppink!12!white}
\colorlet{popgoldL}{popgold!14!white}

\pagecolor{popcream}

% ============================================================
% Typographie
% ============================================================
\defaultfontfeatures{Ligatures=TeX}
\setmainfont{PlusJakartaSans-Medium.ttf}[Path=../../../fonts/,BoldFont=PlusJakartaSans-Bold.ttf]
% Symboles, API et emojis absents de la police principale : repli sur DejaVu Sans (caractères actifs)
\newfontface\symfont{DejaVuSans.ttf}[Path=../../../fonts/]
\makeatletter
\newcommand\symactive[1]{\catcode#1=13 \begingroup\lccode`\~=#1 \lowercase{\endgroup\def~}{{\symfont\char#1}}}
\newcommand\symblank[1]{\catcode#1=13 \begingroup\lccode`\~=#1 \lowercase{\endgroup\def~}{}}
\newcommand\symhyph[1]{\catcode#1=13 \begingroup\lccode`\~=#1 \lowercase{\endgroup\def~}{\mbox{-}}}
\newcount\symc
\newcommand\symrange[2]{\symc=#1 \@whilenum\symc<#2 \do{\expandafter\symactive\expandafter{\the\symc}\advance\symc by 1 }}
\newcommand\symblankrange[2]{\symc=#1 \@whilenum\symc<#2 \do{\expandafter\symblank\expandafter{\the\symc}\advance\symc by 1 }}
\makeatother
\symrange{"0250}{"02FF}\symactive{"2713}\symactive{"2714}\symactive{"2717}\symactive{"2718}\symactive{"2610}\symactive{"2611}\symactive{"2612}
\symrange{"2600}{"27BF}
\catcode"2705=13 \begingroup\lccode`\~="2705 \lowercase{\endgroup\def~}{{\symfont\char"2713}}\symblank{"26BD}\symblank{"26A1}
\symrange{"2460}{"257F}\symactive{"223C}\symactive{"2248}\symactive{"2260}
\symhyph{"2011}\symhyph{"2E1A}\symblankrange{"1F300}{"1FAFF}\symblank{"FE0F}
% Caractères chinois : WenQuanYi Zen Hei (sous-ensemble GB2312 + pinyin), gras/italique simulés
\setCJKmainfont{WenQuanYiZenHei-subset.ttf}[Path=../../../fonts/,AutoFakeBold=3,AutoFakeSlant=0.2]
\xeCJKsetup{CJKspace=false}
% Pinyin : voyelles à caron (ǎ ǐ ǒ ǔ ǖ ǘ ǚ ǜ) absentes de la police principale -> caractères actifs
\newfontface\pyfont{WenQuanYiZenHei-subset.ttf}[Path=../../../fonts/]
\newcommand\pyactive[1]{\catcode#1=13 \begingroup\lccode`\~=#1 \lowercase{\endgroup\def~}{{\pyfont\char#1}}}
\pyactive{"01CD}\pyactive{"01CE}\pyactive{"01CF}\pyactive{"01D0}\pyactive{"01D1}\pyactive{"01D2}\pyactive{"01D3}\pyactive{"01D4}\pyactive{"01D5}\pyactive{"01D6}\pyactive{"01D7}\pyactive{"01D8}\pyactive{"01D9}\pyactive{"01DA}\pyactive{"01DB}\pyactive{"01DC}
% Maths en sans-serif (Fira Math) avec chiffres et lettres droites de Plus
% Jakarta, pour que les formules se fondent dans le texte.
\usepackage[math-style=ISO,bold-style=ISO]{unicode-math}
\setmathfont{FiraMath-Regular.otf}[Path=../../../fonts/]
\setmathfont{PlusJakartaSans-Medium.ttf}[Path=../../../fonts/,range={up/num,up/{Latin,latin},\mathup}]
\usepackage{icomma} % virgule décimale sans espace en mode math
% Caractères mathématiques Unicode absents des polices de texte (Plus Jakarta,
% Archivo Black) : rendus par leur équivalent mathématique, en texte comme en
% formule — sinon ils s'affichent en carré vide (ex. « Arithmétique dans ℤ »).
\usepackage{newunicodechar}
\newunicodechar{ℝ}{\ensuremath{\mathbb{R}}}
\newunicodechar{ℤ}{\ensuremath{\mathbb{Z}}}
\newunicodechar{ℂ}{\ensuremath{\mathbb{C}}}
\newunicodechar{ℕ}{\ensuremath{\mathbb{N}}}
\newunicodechar{ℚ}{\ensuremath{\mathbb{Q}}}
\newunicodechar{𝔻}{\ensuremath{\mathbb{D}}}
\newunicodechar{α}{\ensuremath{\alpha}}
\newunicodechar{β}{\ensuremath{\beta}}
\newunicodechar{γ}{\ensuremath{\gamma}}
\newunicodechar{δ}{\ensuremath{\delta}}
\newunicodechar{ε}{\ensuremath{\varepsilon}}
\newunicodechar{θ}{\ensuremath{\theta}}
\newunicodechar{λ}{\ensuremath{\lambda}}
\newunicodechar{μ}{\ensuremath{\mu}}
\newunicodechar{σ}{\ensuremath{\sigma}}
\newunicodechar{τ}{\ensuremath{\tau}}
\newunicodechar{φ}{\ensuremath{\varphi}}
\newunicodechar{ψ}{\ensuremath{\psi}}
\newunicodechar{ω}{\ensuremath{\omega}}
\newunicodechar{Σ}{\ensuremath{\Sigma}}
% majuscules produites par \MakeUppercase dans les titres (α-aminés -> Α)
\newunicodechar{Α}{\ensuremath{\alpha}}
\newunicodechar{Β}{\ensuremath{\beta}}
\newunicodechar{Γ}{\ensuremath{\Gamma}}
\newunicodechar{Δ}{\ensuremath{\Delta}}
\newunicodechar{Ε}{\ensuremath{\varepsilon}}
\newunicodechar{Θ}{\ensuremath{\Theta}}
\newunicodechar{Λ}{\ensuremath{\Lambda}}
\newunicodechar{Μ}{\ensuremath{\mu}}
\newunicodechar{Τ}{\ensuremath{\tau}}
\newunicodechar{Φ}{\ensuremath{\Phi}}
\newunicodechar{Ψ}{\ensuremath{\Psi}}
\newunicodechar{Ω}{\ensuremath{\Omega}}
\newunicodechar{↦}{\ensuremath{\mapsto}}
\newunicodechar{∈}{\ensuremath{\in}}
\newunicodechar{∉}{\ensuremath{\notin}}
\newunicodechar{∧}{\ensuremath{\wedge}}
\newunicodechar{⇔}{\ensuremath{\Leftrightarrow}}
\newunicodechar{⟺}{\ensuremath{\Longleftrightarrow}}
\newunicodechar{⇒}{\ensuremath{\Rightarrow}}
\newunicodechar{⟹}{\ensuremath{\Longrightarrow}}
\newunicodechar{∘}{\ensuremath{\circ}}
\newunicodechar{∪}{\ensuremath{\cup}}
\newunicodechar{∩}{\ensuremath{\cap}}
\newunicodechar{≡}{\ensuremath{\equiv}}
\newunicodechar{⊂}{\ensuremath{\subset}}
\newunicodechar{⊥}{\ensuremath{\perp}}
\newunicodechar{∃}{\ensuremath{\exists}}
\newunicodechar{∀}{\ensuremath{\forall}}
\newunicodechar{∛}{\ensuremath{\sqrt[3]{\,}}}
\newunicodechar{∜}{\ensuremath{\sqrt[4]{\,}}}
\newunicodechar{⃗}{\ensuremath{{}^{\to}}}
\newunicodechar{ⁿ}{\ifmmode^{n}\else\textsuperscript{\ensuremath{n}}\fi}
\newunicodechar{ˣ}{\ifmmode^{x}\else\textsuperscript{\ensuremath{x}}\fi}
\newunicodechar{ᵇ}{\ifmmode^{b}\else\textsuperscript{\ensuremath{b}}\fi}
\newunicodechar{ʳ}{\ifmmode^{r}\else\textsuperscript{\ensuremath{r}}\fi}
\newunicodechar{ᵘ}{\ifmmode^{u}\else\textsuperscript{\ensuremath{u}}\fi}
\newunicodechar{ʸ}{\ifmmode^{y}\else\textsuperscript{\ensuremath{y}}\fi}
\newunicodechar{ᵃ}{\ifmmode^{a}\else\textsuperscript{\ensuremath{a}}\fi}
\newunicodechar{ᵉ}{\ifmmode^{e}\else\textsuperscript{\ensuremath{e}}\fi}
\newunicodechar{ᵏ}{\ifmmode^{k}\else\textsuperscript{\ensuremath{k}}\fi}
\newunicodechar{ᵖ}{\ifmmode^{p}\else\textsuperscript{\ensuremath{p}}\fi}
\newunicodechar{ᴺ}{\ifmmode^{N}\else\textsuperscript{\ensuremath{N}}\fi}
\newunicodechar{ᶜ}{\ifmmode^{c}\else\textsuperscript{\ensuremath{c}}\fi}
\newunicodechar{ʰ}{\ifmmode^{h}\else\textsuperscript{\ensuremath{h}}\fi}
\newunicodechar{ᵠ}{\ifmmode^{q}\else\textsuperscript{\ensuremath{q}}\fi}
\newunicodechar{ₙ}{\ifmmode_{n}\else\textsubscript{\ensuremath{n}}\fi}
\newunicodechar{ₐ}{\ifmmode_{a}\else\textsubscript{\ensuremath{a}}\fi}
\newunicodechar{ᵢ}{\ifmmode_{i}\else\textsubscript{\ensuremath{i}}\fi}
\newunicodechar{ᵣ}{\ifmmode_{r}\else\textsubscript{\ensuremath{r}}\fi}
\newunicodechar{ₖ}{\ifmmode_{k}\else\textsubscript{\ensuremath{k}}\fi}
\newunicodechar{ᵦ}{\ifmmode_{\beta}\else\textsubscript{\ensuremath{\beta}}\fi}
\newunicodechar{ₓ}{\ifmmode_{x}\else\textsubscript{\ensuremath{x}}\fi}
\newfontfamily\popdisplay{ArchivoBlack-Regular.ttf}[Path=../../../fonts/]
\newfontfamily\poplogo{SpaceGrotesk-Bold.ttf}[Path=../../../fonts/]
\newfontfamily\popsemi{PlusJakartaSans-SemiBold.ttf}[Path=../../../fonts/]
\newfontfamily\popxbold{PlusJakartaSans-ExtraBold.ttf}[Path=../../../fonts/]
\newfontfamily\popxboldls{PlusJakartaSans-ExtraBold.ttf}[Path=../../../fonts/,LetterSpace=3.0]

\setlist[enumerate]{leftmargin=1.7em,itemsep=5pt,topsep=4pt}
\setlist[itemize]{leftmargin=1.7em,itemsep=4pt,topsep=4pt,label={\color{poporange}\textbullet}}
\setlength{\parindent}{0pt}
\setlength{\parskip}{5pt}
\renewcommand{\arraystretch}{1.25}

% pgfplots : style Pop par défaut
\pgfplotsset{
  every axis/.append style={axis line style={popink,line width=1pt},tick style={popink},
    grid style={popline,line width=0.5pt},label style={font=\small},tick label style={font=\footnotesize}},
  popcurve/.style={poporange,line width=1.8pt,no markers,smooth},
}
% Même style disponible dans un \draw TikZ simple (hors pgfplots)
\tikzset{popcurve/.style={poporange,line width=1.8pt}}
\tikzset{popgrid/.style={popline,very thin}}
\colorlet{popcurve}{poporange} % le nom du style est parfois utilisé comme couleur

% ============================================================
% Pill / squiggle
% ============================================================
\newcommand{\poppill}[3]{%
  \tikz[baseline=-0.55ex]\node[draw=popink,line width=1.2pt,rounded corners=2.6pt,%
    fill=#2,inner xsep=6.5pt,inner ysep=3pt]{%
    \popxboldls\fontsize{7}{7}\selectfont\color{#3}\MakeUppercase{#1}};%
}
\newcommand{\popsquiggle}[1]{%
  \tikz[baseline=-0.4ex]{
    \draw[#1,line width=2.6pt,line cap=round]
      plot [smooth] coordinates {(0,0.09) (0.34,0.01) (0.68,0.21) (1.02,0.01) (1.36,0.10)};
  }%
}
% Mot-clé surligné (marqueur orange sous le texte)
\newcommand{\cle}[1]{{\popxbold\color{popdark}#1}}

% ============================================================
% En-tête / pied
% ============================================================
\pagestyle{fancy}
\fancyhf{}
\renewcommand{\headrulewidth}{0pt}
\renewcommand{\footrulewidth}{0pt}
\lhead{\raisebox{-0.15ex}{\poplogo\fontsize{13}{13}\selectfont\color{popink}OnBuch\color{poporange}+}}
\rhead{\poppill{\DOCMATIERE\ \textbullet\ \DOCNIVEAU\ \textbullet\ Leçon \DOCLECON}{white}{popink}}
\lfoot{\popxbold\fontsize{7.6}{7.6}\selectfont\color{popmuted}\MakeUppercase{Cours signé OnBuch+}\ \textbullet\ {\popsemi\DOCTITRE}}
\rfoot{\tikz[baseline=-0.6ex]\node[rounded corners=8.5pt,fill=popink,minimum height=17pt,minimum width=17pt,inner xsep=5pt,inner sep=0]{\popxbold\fontsize{8}{8}\selectfont\color{popcream}\thepage\,/\,\pageref*{LastPage}};}

% ============================================================
% Titres
% ============================================================
\newcommand{\chaptitre}[2]{%
  \begin{tcolorbox}[enhanced,colback=poporange,colframe=popink,boxrule=2.6pt,arc=11pt,
      drop shadow=popink,left=15pt,right=15pt,top=12pt,bottom=11pt,before skip=3pt,after skip=18pt]
    \poppill{#2}{popink}{popcream}\par\vspace{9pt}
    {\raggedright\hyphenpenalty=10000\exhyphenpenalty=10000\popdisplay\fontsize{22}{24}\selectfont\color{popcream}\MakeUppercase{#1}\par}
  \end{tcolorbox}
}
% Section numérotée : pastille orange avec le numéro + titre Archivo
\newcounter{popsec}
\newcommand{\coursec}[1]{%
  \stepcounter{popsec}\setcounter{popsub}{0}%
  \par\vspace{16pt}\noindent
  \begin{minipage}{\linewidth}
  \tikz[baseline=-0.7ex]\node[draw=popink,line width=1.6pt,fill=poporange,rounded corners=4pt,minimum size=22pt,inner sep=0]{\popdisplay\fontsize{12}{12}\selectfont\color{popcream}\thepopsec};%
  \hspace{8pt}{\popdisplay\fontsize{15}{17}\selectfont\color{popink}\MakeUppercase{#1}}\par
  \vspace{4pt}\noindent{\color{poporange}\rule{36pt}{3.6pt}}
  \end{minipage}\par\nopagebreak\vspace{8pt}
}
\newcounter{popsub}
\newcommand{\courssub}[1]{%
  \stepcounter{popsub}%
  \par\vspace{9pt}\noindent{\popxbold\fontsize{11.5}{13}\selectfont\color{popink}{\color{poporange}\thepopsec.\thepopsub}\hspace{6pt}#1}\par\nopagebreak\vspace{4pt}
}

% ============================================================
% Boîtes pédagogiques — même recette : fond blanc, bordure épaisse colorée,
% ombre dure encre, badge plein. Titre optionnel [#1].
% ============================================================
\tcbset{popbase/.style={breakable,enhanced,colback=white,boxrule=2.3pt,arc=9pt,
  shadow={3pt}{-3pt}{0pt}{popink},left=14pt,right=14pt,top=11pt,bottom=11pt,before skip=11pt,after skip=15pt,
  fontupper=\popsemi\fontsize{10.6}{14.5}\selectfont\color{popink},
  fontlower=\popsemi\fontsize{10.6}{14.5}\selectfont\color{popink}}}
% Coche dessinée (le glyphe ✓ n'existe pas dans les polices du document)
\newcommand{\popcheck}[1]{\tikz[baseline=-0.5ex]\draw[#1,line width=1.6pt,line cap=round,line join=round](-0.13,0.0)--(-0.04,-0.09)--(0.13,0.1);}
\providecommand{\coche}{\popcheck{popgreen}}
\providecommand{\resultat}[1]{\par\smallskip\noindent\fcolorbox{popgreen}{popgreenL}{\parbox{\dimexpr\linewidth-2\fboxsep-2\fboxrule\relax}{\textbf{Résultat.} #1}}\par\smallskip}
\newcommand{\popboxhead}[3]{\poppill{#1}{#2}{white}\ifx\relax#3\relax\else\hspace{6pt}{\popxbold\fontsize{10.6}{12}\selectfont\color{popink}#3}\fi\par\vspace{7pt}}

\newtcolorbox{aretenir}[1][]{popbase,colframe=popgreen,before upper={\popboxhead{À retenir}{popgreen}{#1}}}
% Texte en chinois (document de lecture, dialogue, extrait) : cadre
% d'étude. Un titre optionnel (source, auteur) s'affiche dans l'en-tête.
\newtcolorbox{textebox}[1][]{popbase,colframe=popblue,before upper={\popboxhead{文本 · Texte}{popblue}{#1}},after upper={}}
% Marques ✓ / ✗ (forme correcte / fautive) souvent utilisées par les modèles
\providecommand{\cross}{\textcolor{poppink}{\ensuremath{\times}}}
\providecommand{\xmark}{\cross}\providecommand{\crossmark}{\cross}\providecommand{\cmark}{\ensuremath{\checkmark}}
% Ligne de réponse / justification à compléter (inventée par les modèles)
\providecommand{\justification}[1]{\par\noindent\textit{Justification~:} \dotfill\par}
% \textipa (package tipa non chargé) : la police ne couvre pas l'API -> simple passage
\providecommand{\textipa}[1]{#1}
% Soulignement (ulem non chargé) : \uline, \ul
\providecommand{\uline}[1]{\underline{#1}}\providecommand{\ul}[1]{\underline{#1}}
% \glossaire{\item ...} inventée par les modèles : liste de glossaire
\providecommand{\glossaire}[1]{\begin{itemize}#1\end{itemize}}
% \alert (beamer) inventée par les modèles : texte mis en évidence
\providecommand{\alert}[1]{\textbf{\textcolor{poppink}{#1}}}
% variantes ASCII d'ordinaux écrites en macro
\providecommand{\iere}{\textsuperscript{re}}\providecommand{\ieme}{\textsuperscript{e}}\providecommand{\ere}{\textsuperscript{re}}\providecommand{\eme}{\textsuperscript{e}}\providecommand{\er}{\textsuperscript{er}}\providecommand{\ie}{\textsuperscript{e}}
% Ordinaux français écrits en macro par les modèles : 1\ière, 3\ième
\newcommand{\ière}{\textsuperscript{re}}\newcommand{\ième}{\textsuperscript{e}}
% \exemple (inventée par les modèles) : étiquette « Exemple : »
\providecommand{\exemple}{\textbf{Exemple~:}\ }
% \square utilisable aussi hors mode math (cases à cocher)
\AtBeginDocument{\let\squareMath\square\renewcommand{\square}{\ensuremath{\squareMath}}}
% Mot ou phrase en chinois à l'intérieur d'un paragraphe français
\newcommand{\zh}[1]{\ifmmode\text{#1}\else{#1}\fi}
\let\eng\zh \let\esp\zh \let\all\zh % alias tolérés
% flèches d'intonation inventées par les modèles
\providecommand{\textsearrow}{}\renewcommand{\textsearrow}{\ensuremath{\searrow}}\providecommand{\textnearrow}{}\renewcommand{\textnearrow}{\ensuremath{\nearrow}}
% \fr{..} : mot français (inventé par les modèles)
\providecommand{\fr}[1]{\textit{#1}}
% zhbox : encadré de texte chinois inventé par les modèles
\newenvironment{zhbox}{\begin{quote}}{\end{quote}}
% \vs : « versus » inventé par les modèles
\providecommand{\vs}{\textit{vs}\ }
% \blank : case à compléter inventée par les modèles
\providecommand{\blank}[1][]{\underline{\hspace{1.6em}}}
\newtcolorbox{exemplebox}[1][]{popbase,colframe=poporange,before upper={\popboxhead{Exemple}{poporange}{#1}}}
\newtcolorbox{definition}[1][]{popbase,colframe=popblue,before upper={\popboxhead{Définition}{popblue}{#1}}}
\newtcolorbox{propriete}[1][]{popbase,colframe=poppurple,before upper={\popboxhead{Propriété}{poppurple}{#1}}}
\newtcolorbox{methode}[1][]{popbase,colframe=popgold,before upper={\popboxhead{Méthode}{popgold}{#1}}}
\newtcolorbox{attention}[1][]{popbase,colframe=poppink,before upper={\popboxhead{Attention}{poppink}{#1}}}
\newtcolorbox{experience}[1][]{popbase,colframe=popblue,colback=popblueL,before upper={\popboxhead{Expérience}{popblue}{#1}}}
% « Le savais-tu ? » : carte violette pleine (contexte, culture, Cameroun)
\newtcolorbox{savaistu}[1][]{popbase,colback=poppurple,colframe=popink,
  fontupper=\popsemi\fontsize{10.4}{14.2}\selectfont\color{white},
  before upper={\poppill{Le savais-tu\,?}{white}{poppurple}\ifx\relax#1\relax\else\hspace{6pt}{\popxbold\color{white}#1}\fi\par\vspace{7pt}}}
% Exercice résolu : énoncé en haut, \tcblower puis la solution sur fond crème
\newcounter{popexo}\newcounter{popexr}
\newtcolorbox{exoresolu}[1][]{popbase,colframe=poporange,colbacklower=popcream,
  segmentation style={popink,line width=1.2pt,dashed},
  before upper={\stepcounter{popexr}\popboxhead{Exercice résolu \thepopexr}{poporange}{#1}},
  before lower={\poppill{Solution}{popink}{popcream}\par\vspace{6pt}}}
% Exercice (non résolu en place) — la correction est dans la partie Corrigés
\newtcolorbox{exercice}[1][]{popbase,colframe=popink,
  before upper={\stepcounter{popexo}\popboxhead{Exercice \thepopexo}{popink}{#1}}}
% Corrigé
\newtcolorbox{corrige}[1][]{popbase,colframe=popgreen,colback=popgreenL,
  before upper={\popboxhead{Corrigé}{popgreen}{#1}}}
% Objectifs / prérequis / bilan
\newtcolorbox{objectifs}{popbase,colframe=popink,colback=poporangeL,
  before upper={\popboxhead{Objectifs de la leçon}{popink}{}}}
\newtcolorbox{prerequis}{popbase,colframe=popink,colback=popblueL,
  before upper={\popboxhead{Prérequis}{popblue}{}}}
\newtcolorbox{bilan}{popbase,colframe=popink,colback=white,boxrule=2.8pt,
  before upper={\popboxhead{Fiche bilan}{poporange}{L'essentiel à connaître pour l'examen}}}
% Figure encadrée avec légende
\newtcolorbox{popfigure}[1][]{enhanced,breakable=false,colback=white,colframe=popline,boxrule=1.4pt,arc=8pt,
  left=10pt,right=10pt,top=10pt,bottom=8pt,before skip=10pt,after skip=12pt,halign=center,
  lower separated=false,halign lower=center,
  fontlower=\popsemi\fontsize{9}{11.5}\selectfont\color{popmuted},
  #1}
\newcommand{\legende}[1]{\par\vspace{6pt}{\popsemi\fontsize{9}{11.5}\selectfont\color{popmuted}#1}}

% ============================================================
% Signature OnBuch+
% ============================================================
\newcommand{\poplogoinline}[1]{{\poplogo\fontsize{#1}{#1}\selectfont\color{popink}OnBuch\color{poporange}+}}
\newcommand{\signatureonbuch}{%
  \begin{tcolorbox}[enhanced,colback=popink,colframe=popink,boxrule=2.2pt,arc=10pt,
      drop shadow=poporange,left=14pt,right=14pt,top=10pt,bottom=10pt,before skip=0pt,after skip=0pt]
    \tikz[baseline=-0.7ex]\node[circle,fill=popgreen,draw=popcream,line width=1.3pt,minimum size=22pt,inner sep=0]{\popcheck{white}};%
    \hspace{9pt}\begin{minipage}[c]{0.62\linewidth}
      {\popxbold\fontsize{12}{13}\selectfont\color{popcream}Signé\ \textbullet\ {\poplogo OnBuch\color{poporange}+}}\par\vspace{2pt}
      {\raggedright\popsemi\fontsize{8}{10.5}\selectfont\color{popline}\MakeUppercase{Équipe pédagogique OnBuch+}\\\MakeUppercase{Conforme au programme officiel MINESEC}\par}
    \end{minipage}\hfill
    {\poplogo\fontsize{22}{22}\selectfont\color{popcream}OnBuch\color{poporange}+}
  \end{tcolorbox}%
}

% ============================================================
% Page de garde
% #1 titre · #2 matière · #3 niveau · #4 module · #5 n° leçon · #6 durée ·
% #7 description / objectifs (texte ou itemize)
% ============================================================
\newcommand{\pagegarde}[7]{%
  \thispagestyle{empty}
  \newgeometry{a4paper,margin=1.7cm,top=1.6cm,bottom=1.4cm}%
  \begin{tikzpicture}[remember picture,overlay]
    \fill[poporange!18] ([xshift=-3.2cm,yshift=-3.6cm]current page.north east) circle (4.6cm);
    \fill[poppurple!14] ([xshift=2.4cm,yshift=7.6cm]current page.south west) circle (3.8cm);
  \end{tikzpicture}%
  {\poplogo\fontsize{30}{30}\selectfont\color{popink}OnBuch\color{poporange}+}\hfill
  \raisebox{0.6ex}{\poppill{Cours signé \textbullet\ Premium}{popink}{popcream}}\par
  \vspace{1pt}\popsquiggle{poppurple}\par
  \vspace{8pt}
  \begin{tcolorbox}[enhanced,colback=poporange,colframe=popink,boxrule=2.8pt,arc=13pt,
      drop shadow=popink,left=18pt,right=18pt,top=12pt,bottom=13pt,after skip=10pt]
    \poppill{#2 \textbullet\ #3}{popink}{popcream}\hspace{5pt}\poppill{#4}{white}{popink}\par\vspace{8pt}
    {\popdisplay\fontsize{14}{14}\selectfont\color{popink}LEÇON #5}\par\vspace{4pt}
    {\raggedright\hyphenpenalty=10000\exhyphenpenalty=10000\popdisplay\fontsize{25}{27}\selectfont\color{popcream}\MakeUppercase{#1}\par}
    \vspace{8pt}
    {\popxbold\fontsize{9.5}{9.5}\selectfont\color{popink}\MakeUppercase{Durée conseillée : #6}}
  \end{tcolorbox}
  \begin{tcolorbox}[enhanced,colback=white,colframe=popink,boxrule=2.2pt,arc=10pt,
      drop shadow=popink,left=16pt,right=16pt,top=10pt,bottom=10pt,after skip=8pt]
    \poppill{Dans cette leçon}{poppurple}{white}\par\vspace{5pt}
    \popsemi\fontsize{9.3}{12.6}\selectfont\color{popink}#7
  \end{tcolorbox}
  \noindent\begin{minipage}[t]{0.487\linewidth}
    \begin{tcolorbox}[enhanced,colback=popgreen,colframe=popink,boxrule=2.2pt,arc=10pt,
        drop shadow=popink,left=11pt,right=11pt,top=10pt,bottom=10pt,height=3.1cm,valign=center]
      \tcbox[colback=white,colframe=popink,boxrule=1.3pt,arc=5pt,boxsep=0pt,nobeforeafter,
        left=3pt,right=3pt,top=3pt,bottom=3pt,valign=center]{\qrcode[height=1.9cm]{https://play.google.com/store/apps/details?id=com.luvvix.onbuch&hl=fr}}%
      \hspace{8pt}\begin{minipage}[c]{0.5\linewidth}
        {\popdisplay\fontsize{11}{12.5}\selectfont\color{white}\MakeUppercase{Télécharge OnBuch}}\par\vspace{4pt}
        {\raggedright\popsemi\fontsize{8.4}{11}\selectfont\color{white}Gratuit \textbullet\ Google Play\\Tuteur IA, annales, concours, résultats\par}
      \end{minipage}
    \end{tcolorbox}
  \end{minipage}\hfill
  \begin{minipage}[t]{0.487\linewidth}
    \begin{tcolorbox}[enhanced,colback=poppurple,colframe=popink,boxrule=2.2pt,arc=10pt,
        drop shadow=popink,left=11pt,right=11pt,top=10pt,bottom=10pt,height=3.1cm,valign=center]
      \tikz[baseline=-0.7ex]\node[circle,fill=white,draw=popink,line width=1.3pt,minimum size=1.6cm,inner sep=0]{\popdisplay\fontsize{24}{24}\selectfont\color{poppurple}+};%
      \hspace{8pt}\begin{minipage}[c]{0.6\linewidth}
        {\popdisplay\fontsize{11}{12.5}\selectfont\color{white}\MakeUppercase{Passe à OnBuch+}}\par\vspace{4pt}
        {\raggedright\popsemi\fontsize{8.4}{11}\selectfont\color{white}Tous les cours signés, Tuteur IA illimité, fiches PDF — onglet OnBuch+ de l'app\par}
      \end{minipage}
    \end{tcolorbox}
  \end{minipage}
  \vspace{8pt}
  \popcontenu
  \vfill
  \signatureonbuch
  \restoregeometry
  \newpage
}

% Rangée « Ce cours contient » (page de garde)
\newcommand{\poptile}[3]{% #1 couleur #2 gros texte #3 légende
  \begin{tcolorbox}[enhanced,colback=#1,colframe=popink,boxrule=2pt,arc=9pt,shadow={2.5pt}{-2.5pt}{0pt}{popink},
      left=6pt,right=6pt,top=9pt,bottom=9pt,halign=center,valign=center,height=1.75cm,width=\linewidth,nobeforeafter]
    {\popdisplay\fontsize{12.5}{13}\selectfont\color{white}\MakeUppercase{#2}}\par\vspace{3pt}
    {\centering\popsemi\fontsize{7.8}{9.5}\selectfont\color{white}#3\par}
  \end{tcolorbox}}
\newcommand{\popcontenu}{%
  \par\noindent{\popxboldls\fontsize{7.5}{7.5}\selectfont\color{popmuted}CE COURS SIGNÉ CONTIENT}\par\vspace{7pt}
  \noindent\begin{minipage}[t]{0.235\linewidth}\poptile{popblue}{Cours}{complet, illustré, pas à pas}\end{minipage}\hfill
  \begin{minipage}[t]{0.235\linewidth}\poptile{poporange}{Méthodes}{et exercices résolus}\end{minipage}\hfill
  \begin{minipage}[t]{0.235\linewidth}\poptile{poppink}{Exercices}{type examen + corrigés}\end{minipage}\hfill
  \begin{minipage}[t]{0.235\linewidth}\poptile{popgreen}{Fiche bilan}{l'essentiel à retenir}\end{minipage}\par}

% Sommaire
\newtcolorbox{sommaire}{enhanced,colback=white,colframe=popink,boxrule=2.2pt,arc=10pt,
  drop shadow=popink,left=18pt,right=18pt,top=13pt,bottom=13pt,before skip=6pt,after skip=20pt,
  before upper={\poppill{Sommaire}{poppurple}{white}\par\vspace{9pt}\popsemi\fontsize{10.6}{14.5}\selectfont\color{popink}}}

% Signature de fin de cours — poussée tout en bas de la dernière page
\newcommand{\findecours}{%
  \par\vfill\nopagebreak
  \begin{center}\popsquiggle{poporange}\end{center}\vspace{4pt}
  \signatureonbuch
}
\AtBeginDocument{\def\llbracket{\mathopen{[\mkern-3.5mu[}}\def\rrbracket{\mathclose{]\mkern-3.5mu]}}} % intervalles d'entiers (glyphe ⟦ absent de la police)
% Glyphes absents de Fira Math : redéfinis à partir de glyphes disponibles
\AtBeginDocument{%
  \def\setminus{\mathbin{\backslash}}\let\smallsetminus\setminus
  \def\vdots{\mathord{\vbox{\baselineskip3.2pt\lineskiplimit0pt\kern4pt\hbox{.}\hbox{.}\hbox{.}}}}%
  \def\ddots{\mathinner{\mkern1mu\raise6.5pt\hbox{.}\mkern2mu\raise3.6pt\hbox{.}\mkern2mu\raise0.7pt\hbox{.}\mkern1mu}}%
  \def\triangle{\mathord{\scalebox{0.9}{$\symup{\Delta}$}}}%
  \def\nabla{\mathord{\raisebox{\depth}{\scalebox{1}[-1]{$\symup{\Delta}$}}}}%
  \def\checkmark{\popcheck{popdark}}%
  \def\star{\mathbin{\ast}}\def\bigstar{\mathord{\ast}}%
  \def\diamond{\mathbin{\scalebox{0.6}{$\Diamond$}}}%
  \def\bigcup{\mathop{\mathchoice{\raisebox{-0.3em}{\scalebox{1.7}{$\cup$}}}{\scalebox{1.25}{$\cup$}}{\cup}{\cup}}\displaylimits}%
  \def\bigcap{\mathop{\mathchoice{\raisebox{-0.3em}{\scalebox{1.7}{$\cap$}}}{\scalebox{1.25}{$\cap$}}{\cap}{\cap}}\displaylimits}%
  \def\lesseqqgtr{\mathrel{\lessgtr}}\def\bigoplus{\mathop{\oplus}\displaylimits}%
}
\providecommand{\ssi}{si et seulement si} % abréviation fréquente
\AtBeginDocument{\providecommand{\wideparen}{\overparen}} % arc de cercle

\begin{document}

\pagegarde{Leçon 10 — Grammaire : Phrases comparatives avec…不如… ; A比B 形容词+一点ou 多了 ; A比B更/ 还+形容词 ; 为了 ; 不但…而且…}{Chinois}{1ère A}{Module 2 : Environnement et santé}{ZHA10}{2 h}{Cette leçon de grammaire présente cinq structures essentielles du chinois niveau 1ère : les comparatifs avec 不如 et 比 (nuancés par 一点/多了 et 更/还), la construction de but avec 为了 et la corrélation additive 不但…而且…. Chaque structure est observée, schématisée, comparée au français et entraînée par des exercices de transformation.
\vspace{4pt}
\begin{multicols}{2}\small
\begin{itemize}
\item À la fin de cette leçon, je sais construire une phrase comparative négative avec A不如B (+ adjectif)
\item À la fin de cette leçon, je sais exprimer un écart faible ou grand avec A比B + adjectif + 一点/一些/得多/多了
\item À la fin de cette leçon, je sais renforcer une comparaison avec A比B更/还 + adjectif
\item À la fin de cette leçon, je sais exprimer le but avec 为了 + objectif, + action
\item À la fin de cette leçon, je sais relier deux idées additives avec 不但…而且…
\item À la fin de cette leçon, je sais placer correctement les mots-outils (比, 不如, 更, 还) dans la phrase et éviter les calques du français
\end{itemize}\end{multicols}}

\begin{sommaire}
\begin{multicols}{2}
\begin{enumerate}
\item Observation et rappel : la comparaison avec 比
\item Nuancer l'écart : A比B + adjectif + 一点/多了
\item Renforcer la comparaison : A比B更/还 + adjectif
\item La comparaison négative : A不如B (+ adjectif)
\item Exprimer le but avec 为了
\item Additionner les idées : 不但…而且… et synthèse
\item Activité d'intégration
\item Méthodes et exercices résolus
\item Exercices
\item Corrigés des exercices
\item Fiche bilan
\end{enumerate}
\end{multicols}
\end{sommaire}

\begin{objectifs}
\begin{itemize}
\item À la fin de cette leçon, je sais construire une phrase comparative négative avec A不如B (+ adjectif)
\item À la fin de cette leçon, je sais exprimer un écart faible ou grand avec A比B + adjectif + 一点/一些/得多/多了
\item À la fin de cette leçon, je sais renforcer une comparaison avec A比B更/还 + adjectif
\item À la fin de cette leçon, je sais exprimer le but avec 为了 + objectif, + action
\item À la fin de cette leçon, je sais relier deux idées additives avec 不但…而且…
\item À la fin de cette leçon, je sais placer correctement les mots-outils (比, 不如, 更, 还) dans la phrase et éviter les calques du français
\item À la fin de cette leçon, je sais transformer une phrase affirmative en phrase comparative et inversement
\item À la fin de cette leçon, je sais employer ces cinq structures dans un court texte comparatif (ville, école, santé)
\end{itemize}
\end{objectifs}

\begin{prerequis}
\begin{itemize}
\item La structure comparative de base A比B + adjectif (vue en début d'année)
\item Les adjectifs qualificatifs courants (好, 大, 热, 贵, 多, 快…) et leur emploi sans 是
\item La négation avec 不 et 没有
\item L'ordre de base de la phrase chinoise : sujet – verbe – complément
\item Le vocabulaire du Module 2 : environnement et santé (身体, 环境, 运动, 医院…)
\end{itemize}
\end{prerequis}

\newpage

\coursec{Observation et rappel : la comparaison avec 比}

\courssub{Situation d'accroche}

Pendant les vacances, deux élèves de Yaoundé comparent leurs villes : « Douala est plus chaude que Yaoundé, mais Yaoundé est moins bruyante que Douala. » Un autre ajoute : « Pour réussir au Bac, je travaille non seulement le chinois, mais aussi l'anglais. » En chinois, ces comparaisons et ces justifications obéissent à des structures précises : \zh{不如}, \zh{A比B…}, \zh{为了}, \zh{不但…而且…}. Maîtriser ces cinq structures, c'est pouvoir \cle{comparer}, \cle{nuancer} et \cle{argumenter} comme un vrai locuteur chinois.

\textbf{Problématique :} comment dire en chinois « plus… que », « moins… que », « non seulement… mais aussi », sans jamais calquer le français ? Cette première section pose le fondement de toute la leçon : la structure de base de la comparaison avec \zh{比}.

\courssub{Lecture guidée : observer avant de comprendre}

Lisons attentivement les phrases suivantes. Ne cherchez pas encore la règle : observez simplement \emph{l'ordre des mots} et ce qui est présent ou absent.

\begin{textebox}[Quatre exemples à observer]
今天比昨天热。\\
\emph{Jīntiān bǐ zuótiān rè.}\\
« Aujourd'hui, il fait plus chaud qu'hier. »\\[0.4em]
汉语比法语难。\\
\emph{Hànyǔ bǐ Fǎyǔ nán.}\\
« Le chinois est plus difficile que le français. »\\[0.4em]
杜阿拉比雅温得热。\\
\emph{Dù'ālā bǐ Yǎwēndé rè.}\\
« Douala est plus chaude que Yaoundé. »\\[0.4em]
足球比篮球受欢迎。\\
\emph{Zúqiú bǐ lánqiú shòu huānyíng.}\\
« Le football est plus populaire que le basket-ball. »
\end{textebox}

\textbf{Questions d'observation :}
\begin{enumerate}
\item Quel mot revient dans les quatre phrases ? Où est-il placé ?
\item Y a-t-il un mot comme \zh{很} (très) ou \zh{是} (être) avant l'adjectif ?
\item En français, quel mot correspond à \zh{比} ?
\end{enumerate}

\textbf{Réponses attendues :} le mot \zh{比} (\emph{bǐ}) revient toujours, placé \emph{entre} les deux éléments comparés ; il n'y a ni \zh{很} ni \zh{是} avant l'adjectif ; \zh{比} correspond au « que » français de la comparaison, mais il porte \emph{aussi} l'idée de « plus ».

\courssub{Le schéma de base : A + 比 + B + adjectif}

\begin{definition}[La comparaison avec 比]
\zh{比} (\emph{bǐ}) est une préposition qui introduit le second terme de la comparaison. La structure \cle{A + 比 + B + adjectif} signifie « A est plus (adjectif) que B ». L'adjectif seul, sans aucun autre mot, exprime déjà la supériorité.
\end{definition}

\begin{propriete}[Schéma et emploi de la structure 比]
\textbf{Schéma :} A + \zh{比} + B + adjectif\\
\textbf{Emploi :} comparer deux personnes, deux lieux, deux moments, deux choses sur une qualité.\\
\textbf{Point capital :} dans cette structure, l'adjectif ne prend \textbf{jamais} \zh{很} ni \zh{非常} ni \zh{是} : le degré (« plus ») est déjà exprimé par la comparaison elle-même.
\end{propriete}

\begin{popfigure}
\begin{tikzpicture}[node distance=0.8cm, >=Stealth]
\node[draw=popblue, fill=popblueL, rounded corners, minimum width=1.8cm, minimum height=1cm] (A) {A};
\node[draw=poporange, fill=poporangeL, rounded corners, minimum width=1.8cm, minimum height=1cm, right=of A] (bi) {\zh{比}};
\node[draw=popblue, fill=popblueL, rounded corners, minimum width=1.8cm, minimum height=1cm, right=of bi] (B) {B};
\node[draw=popgreen, fill=popgreenL, rounded corners, minimum width=2.5cm, minimum height=1cm, right=of B] (adj) {adjectif};
\draw[thick, popdark, ->] (A) -- (bi);
\draw[thick, popdark, ->] (bi) -- (B);
\draw[thick, popdark, ->] (B) -- (adj);
\node[draw=poppink, fill=poppinkL, rounded corners, below=1.2cm of bi, align=center, text width=3cm] (hen) {\textbf{Interdits ici :} \\ \zh{很} \quad \zh{非常} \quad \zh{是}};
\draw[thick, poppink] (hen.north west) -- (hen.south east);
\draw[thick, poppink] (hen.north east) -- (hen.south west);
\end{tikzpicture}
\legende{Schéma de la structure A + 比 + B + adjectif : aucun mot de degré ne peut s'intercaler entre 比 et l'adjectif.}
\end{popfigure}

\begin{exemplebox}[Exemples traduits]
\zh{杜阿拉比雅温得热。} \emph{Dù'ālā bǐ Yǎwēndé rè.} « Douala est plus chaude que Yaoundé. »\\
\zh{我哥哥比我高。} \emph{Wǒ gēge bǐ wǒ gāo.} « Mon grand frère est plus grand que moi. »\\
\zh{这本书比那本书有意思。} \emph{Zhè běn shū bǐ nà běn shū yǒu yìsi.} « Ce livre est plus intéressant que celui-là. »\\
\zh{今天比昨天冷。} \emph{Jīntiān bǐ zuótiān lěng.} « Aujourd'hui, il fait plus froid qu'hier. »\\
\zh{喀麦隆的足球比篮球有名。} \emph{Kāmàilóng de zúqiú bǐ lánqiú yǒumíng.} « Au Cameroun, le football est plus célèbre que le basket-ball. »
\end{exemplebox}

\courssub{Boîte à outils : adjectifs fréquents pour la comparaison}

\begin{center}
\begin{tabularx}{\linewidth}{|X|X|X|X|}
\hline
\rowcolor{popblueL} Caractères & Pinyin & Nature & Traduction \\
\hline
热 & rè & adj. & chaud \\
\hline
冷 & lěng & adj. & froid \\
\hline
大 & dà & adj. & grand, gros \\
\hline
小 & xiǎo & adj. & petit, jeune (âge) \\
\hline
高 & gāo & adj. & grand (taille), haut \\
\hline
矮 & ǎi & adj. & petit (taille), bas \\
\hline
贵 & guì & adj. & cher \\
\hline
便宜 & piányi & adj. & bon marché \\
\hline
远 & yuǎn & adj. & loin \\
\hline
近 & jìn & adj. & près \\
\hline
漂亮 & piàoliang & adj. & beau, joli \\
\hline
有意思 & yǒu yìsi & adj. & intéressant \\
\hline
受欢迎 & shòu huānyíng & adj. & populaire, apprécié \\
\hline
有名 & yǒumíng & adj. & célèbre \\
\hline
难 & nán & adj. & difficile \\
\hline
容易 & róngyì & adj. & facile \\
\hline
\end{tabularx}
\end{center}

\courssub{Deux interdictions absolues : pas de 很, pas de 是}

\begin{attention}[Erreurs fréquentes des francophones]
En français, on dit « il est \emph{plus} grand que moi » : le mot « plus » est séparé de « que ». Le francophone traduit donc souvent « plus » par \zh{很} et obtient la phrase fautive *\zh{他比我很高}. C'est \textbf{impossible} : en chinois, \zh{比} seul porte déjà l'idée de « plus ». Ajouter \zh{很} ou \zh{非常} devant l'adjectif est une faute grave.\\
De même, le verbe « être » français pousse à écrire *\zh{他比我是高} : c'est faux aussi. L'adjectif chinois fonctionne comme un verbe d'état : il n'a pas besoin de \zh{是}.
\end{attention}

\begin{center}
\begin{tabularx}{\linewidth}{|X|X|X|}
\hline
\rowcolor{popblueL} Phrase & Statut & Explication \\
\hline
\zh{他比我高。} Tā bǐ wǒ gāo. & Correcte & \zh{比} + adjectif seul : « Il est plus grand que moi. » \\
\hline
\zh{他比我很高。} & Fausse & \zh{很} interdit entre \zh{比} et l'adjectif. \\
\hline
\zh{他比我是高。} & Fausse & \zh{是} interdit : l'adjectif suffit. \\
\hline
\zh{他非常高，比我高。} & Correcte & \zh{非常} possible dans une phrase \emph{séparée}, sans \zh{比}. \\
\hline
\end{tabularx}
\end{center}

\courssub{La négation : 不比 ne veut pas dire « moins »}

\begin{propriete}[Forme négative de la comparaison avec 比]
\textbf{Schéma :} A + \zh{不比} + B + adjectif\\
\textbf{Sens :} « A n'est pas plus (adjectif) que B », c'est-à-dire que A et B sont à peu près égaux, ou que A est même un peu moins.\\
\textbf{Attention :} \zh{A不比B} + adj n'est \emph{pas} la même chose que \zh{A比B不} + adj, qui est une structure très rare et marquée (« A est plus non-(adjectif) que B »). Pour dire simplement « moins… que », le chinois utilise \zh{不如} (section 4) ou \zh{没有}.
\end{propriete}

\begin{exemplebox}[Comparer les deux négations]
\zh{今天不比昨天热。} \emph{Jīntiān bù bǐ zuótiān rè.} « Aujourd'hui, il ne fait pas plus chaud qu'hier. » (les deux jours se valent)\\
\zh{雅温得不比杜阿拉热。} \emph{Yǎwēndé bù bǐ Dù'ālā rè.} « Yaoundé n'est pas plus chaude que Douala. »\\
\zh{我不比他矮。} \emph{Wǒ bù bǐ tā ǎi.} « Je ne suis pas plus petit que lui. » (nous faisons à peu près la même taille)
\end{exemplebox}

\begin{center}
\begin{tabularx}{\linewidth}{|X|X|X|}
\hline
\rowcolor{popblueL} Structure & Exemple (caractères + pinyin) & Traduction \\
\hline
A \zh{比} B + adj & \zh{他比我高。} Tā bǐ wǒ gāo. & Il est plus grand que moi. \\
\hline
A \zh{不比} B + adj & \zh{他不比我高。} Tā bù bǐ wǒ gāo. & Il n'est pas plus grand que moi. \\
\hline
A \zh{比} B + \zh{不} + adj & \zh{他比我不好。} (rare) & Il est plus « pas bien » que moi (il est pire que moi). \\
\hline
\end{tabularx}
\end{center}

\courssub{Comparaison avec le français}

\begin{center}
\begin{tabularx}{\linewidth}{|X|X|X|}
\hline
\rowcolor{popblueL} Notion & Français & Chinois \\
\hline
Marque de la supériorité & « plus… que » (deux mots séparés) & \zh{比} seul, placé entre A et B \\
\hline
Verbe « être » & obligatoire : « il \emph{est} plus grand » & absent : l'adjectif suffit \\
\hline
Mot de degré & « plus », « beaucoup plus » & pas de \zh{很}/\zh{非常} ; on verra \zh{多了}, \zh{更}, \zh{还} \\
\hline
Négation & « il n'est pas plus grand que » & \zh{不比} devant B \\
\hline
\end{tabularx}
\end{center}

\courssub{Carte mentale de synthèse}

\begin{popfigure}
\begin{tikzpicture}[mindmap, grow cyclic, every node/.style={concept, align=center, font=\small}, root concept/.append style={fill=poporange, minimum size=2.2cm}, level 1/.append style={level distance=3.2cm, sibling angle=120}, level 1 concept/.append style={minimum size=1.9cm}]
\node[root concept]{\zh{比} bǐ}
  child[concept color=popblueL] { node[align=center]{Schéma \\ A \zh{比} B + adj} }
  child[concept color=poppinkL] { node[align=center]{Interdits \\ pas de \zh{很} \\ pas de \zh{是}} }
  child[concept color=popgreenL] { node[align=center]{Négation \\ A \zh{不比} B + adj} };
\end{tikzpicture}
\legende{Carte mentale : les trois choses à retenir sur la comparaison avec 比.}
\end{popfigure}

\courssub{Exercice résolu et premier entraînement oral}

\begin{exoresolu}[Traduire en chinois]
Énoncé : traduisez en chinois les phrases suivantes. 1) « Yaoundé est plus grande que Bafoussam. » 2) « Le chinois n'est pas plus difficile que l'anglais. » 3) « Ma sœur est plus jeune que moi. »
\tcblower
Solution détaillée :
\begin{enumerate}
\item \zh{雅温得比巴富萨姆大。} \emph{Yǎwēndé bǐ Bāfùsàmǔ dà.} — Schéma A + \zh{比} + B + adjectif ; aucun \zh{很} devant \zh{大}.
\item \zh{汉语不比英语难。} \emph{Hànyǔ bù bǐ Yīngyǔ nán.} — Négation : \zh{不} se place \emph{devant} \zh{比}, jamais devant l'adjectif.
\item \zh{我妹妹比我小。} \emph{Wǒ mèimei bǐ wǒ xiǎo.} — Pour l'âge, le chinois dit \zh{大} (grand/âgé) et \zh{小} (petit/jeune), pas « vieux/jeune » littéralement.
\end{enumerate}
\end{exoresolu}

\begin{experience}[Premier entraînement oral : Yaoundé contre Douala]
Par groupes de deux, comparez Yaoundé et Douala en utilisant uniquement le schéma A + \zh{比} + B + adjectif, puis sa négation \zh{不比}. Adjectifs utiles : \zh{热} (chaud), \zh{大} (grand), \zh{贵} (cher), \zh{远} (loin), \zh{漂亮} (beau).\\
Modèles : \zh{杜阿拉比雅温得热。} \emph{Dù'ālā bǐ Yǎwēndé rè.} « Douala est plus chaude que Yaoundé. » ; \zh{雅温得不比杜阿拉大。} \emph{Yǎwēndé bù bǐ Dù'ālā dà.} « Yaoundé n'est pas plus grande que Douala. »\\
Chaque binôme produit quatre phrases : deux affirmatives, deux négatives, puis les présente à la classe.
\end{experience}

\begin{aretenir}[L'essentiel de la section]
\begin{itemize}
\item Schéma de base : \cle{A + 比 + B + adjectif} = « A est plus (adjectif) que B ».
\item \zh{比} porte à lui seul l'idée de « plus » : jamais de \zh{很}, \zh{非常} ou \zh{是} devant l'adjectif.
\item Négation : \zh{不} se place devant \zh{比} (\zh{A不比B} + adj = « A n'est pas plus… que B »).
\item Pour dire « moins… que », on utilisera \zh{不如} (section 4) ; pour nuancer l'écart, \zh{一点} / \zh{多了} (section 2).
\end{itemize}
\end{aretenir}

\coursec{Nuancer l'écart : A比B + adjectif + 一点/多了}

Dans la section précédente, nous avons revu la structure de base de la comparaison avec \zh{比} : \textbf{A + 比 + B + adjectif}, comme dans \zh{他比我高。} \emph{Tā bǐ wǒ gāo.} « Il est plus grand que moi. » Mais cette phrase reste vague : plus grand \emph{de combien} ? D'un centimètre ou de vingt ? En français, on dirait « un peu plus grand », « beaucoup plus grand », « deux centimètres de plus ». Le chinois possède des marqueurs précis pour cela : \cle{一点}, \cle{一些}, \cle{多了} et \cle{得多}. C'est l'objet de cette section : apprendre à \textbf{nuancer l'écart} d'une comparaison.

\courssub{Observation : deux phrases à comparer}

\begin{textebox}[Dialogue au marché de Mokolo (Yaoundé)]
\zh{甲：这件衬衫怎么样？}\\
\emph{Jiǎ : Zhè jiàn chènshān zěnmeyàng ?}\\
« A : Comment trouves-tu cette chemise ? »\\[4pt]
\zh{乙：不错，可是比那件贵一点。}\\
\emph{Yǐ : Búcuò, kěshì bǐ nà jiàn guì yìdiǎn.}\\
« B : Pas mal, mais elle est un peu plus chère que l'autre. »\\[4pt]
\zh{甲：那边的店便宜多了！}\\
\emph{Jiǎ : Nàbiān de diàn piányi duō le !}\\
« A : La boutique là-bas est bien moins chère ! »
\end{textebox}

\textbf{Observons.} Dans les deux répliques de B et de A, on retrouve la structure \zh{比} déjà connue. Mais quelque chose a été \emph{ajouté} :

\begin{itemize}
\item \zh{贵一点} \emph{guì yìdiǎn} : l'adjectif \zh{贵} « cher » est suivi de \zh{一点} « un peu » → écart \textbf{faible}.
\item \zh{便宜多了} \emph{piányi duō le} : l'adjectif \zh{便宜} « bon marché » est suivi de \zh{多了} → écart \textbf{grand}.
\end{itemize}

Regardons encore deux exemples isolés :

\zh{他比我高一点。} \emph{Tā bǐ wǒ gāo yìdiǎn.} « Il est un peu plus grand que moi. »

\zh{这本书比那本贵多了。} \emph{Zhè běn shū bǐ nà běn guì duō le.} « Ce livre est beaucoup plus cher que celui-là. »

La logique est limpide : le marqueur qui mesure l'écart vient \textbf{après} l'adjectif, comme une précision ajoutée à la fin.

\courssub{La règle : écart faible et écart grand}

\begin{propriete}[Marqueurs d'écart dans la comparaison avec 比]
Pour préciser l'ampleur de la différence, on place un marqueur d'écart \textbf{après l'adjectif} :
\begin{itemize}
\item \textbf{Écart faible} : A + 比 + B + adjectif + \cle{一点} (ou \cle{一些}, légèrement plus soutenu)\\
\zh{今天比昨天冷一点。} \emph{Jīntiān bǐ zuótiān lěng yìdiǎn.} « Il fait un peu plus froid aujourd'hui qu'hier. »
\item \textbf{Écart grand} : A + 比 + B + adjectif + \cle{多了} (ou \cle{得多}, registre écrit)\\
\zh{杜阿拉比雅温得热多了。} \emph{Dùlā'ā bǐ Yǎwēndé rè duō le.} « Douala est beaucoup plus chaude que Yaoundé. »
\end{itemize}
\textbf{Point capital :} le marqueur d'écart se place \textbf{toujours après l'adjectif}, jamais avant.
\end{propriete}

\begin{popfigure}
\begin{tikzpicture}[scale=1]
\draw[-{Stealth}, very thick, popline] (0,0) -- (10.5,0);
\fill[popmuted] (1,0) circle (0.09);
\node[align=center, font=\small, text=popmuted] at (1,0.95) {adjectif seul\\ \zh{高}};
\fill[popblue] (3.8,0) circle (0.09);
\node[align=center, font=\small, text=popblue] at (3.8,0.95) {+ \zh{一点}\\ écart faible};
\fill[popgreen] (6.4,0) circle (0.09);
\node[align=center, font=\small, text=popgreen] at (6.4,0.95) {+ \zh{一些}\\ un peu plus soutenu};
\fill[poporange] (9.2,0) circle (0.09);
\node[align=center, font=\small, text=poporange] at (9.2,0.95) {+ \zh{多了} / \zh{得多}\\ écart fort};
\node[font=\small\itshape, text=popmuted] at (5.25,-0.55) {intensité croissante de l'écart};
\end{tikzpicture}
\legende{Frise d'intensité : de l'adjectif seul à l'écart fort.}
\end{popfigure}

\begin{popfigure}
\begin{tikzpicture}[node distance=0.25cm]
\node[draw=popblue, fill=popblueL, rounded corners, minimum width=1.1cm, minimum height=0.8cm] (A) {\zh{A}};
\node[draw=popdark, fill=popcream, rounded corners, right=of A, minimum width=1.1cm, minimum height=0.8cm] (bi) {\zh{比}};
\node[draw=popblue, fill=popblueL, rounded corners, right=of bi, minimum width=1.1cm, minimum height=0.8cm] (B) {\zh{B}};
\node[draw=poppurple, fill=poppurpleL, rounded corners, right=of B, minimum width=1.8cm, minimum height=0.8cm] (adj) {adjectif};
\node[draw=poporange, fill=poporangeL, rounded corners, right=of adj, minimum width=2.4cm, minimum height=0.8cm] (ec) {\zh{一点} / \zh{多了}};
\draw[-{Stealth}, thick, poporange] (adj.south) .. controls +(0,-0.7) and +(0,-0.7) .. node[below, font=\small, text=poporange]{l'écart se mesure ici, après l'adjectif} (ec.south);
\end{tikzpicture}
\legende{Schéma de la structure : A 比 B + adjectif + marqueur d'écart.}
\end{popfigure}

\courssub{Les quatre marqueurs en détail}

\begin{definition}[一点, 一些, 多了, 得多]
\begin{itemize}
\item \zh{一点} \emph{yìdiǎn} : « un peu », écart faible, registre courant et oral. C'est le marqueur le plus fréquent.
\item \zh{一些} \emph{yìxiē} : « un peu, quelque peu », même valeur que \zh{一点} mais légèrement plus soutenu, fréquent à l'écrit.
\item \zh{多了} \emph{duō le} : « beaucoup plus », écart grand, registre oral et courant. Le \zh{了} final est obligatoire.
\item \zh{得多} \emph{de duō} : « beaucoup plus », équivalent de \zh{多了}, registre plutôt écrit. Attention : ici \zh{得} s'écrit avec la clé du pas (彳), pas \zh{的} ni \zh{地}.
\end{itemize}
\end{definition}

\begin{savaistu}[La monnaie et les classifieurs de prix]
En Chine, l'unité officielle est le \zh{元} \emph{yuán} (symbole ¥), mais à l'oral on utilise presque toujours \zh{块} \emph{kuài} (classifieur pour l'argent, comme « balles » ou « euros » en familier).
\begin{itemize}
\item \zh{十块} \emph{shí kuài} = 10 yuans (≈ 600 FCFA).
\item \zh{五块钱} \emph{wǔ kuài qián} = 5 yuans (le mot \zh{钱} « argent » est souvent ajouté).
\end{itemize}
Pour les prix « à la pièce », on n'utilise pas de classifieur entre le nombre et \zh{块} : \zh{贵十块} (plus cher de 10 yuans).
\end{savaistu}

\begin{exemplebox}[Exemples traduits]
\zh{坐公共汽车比坐出租车便宜多了。}\\
\emph{Zuò gōnggòng qìchē bǐ zuò chūzūchē piányi duō le.}\\
« Prendre le bus est bien moins cher que prendre le taxi. »\\[4pt]
\zh{汉语比英语难一点。}\\
\emph{Hànyǔ bǐ Yīngyǔ nán yìdiǎn.}\\
« Le chinois est un peu plus difficile que l'anglais. »\\[4pt]
\zh{这个城市比那个城市干净一些。}\\
\emph{Zhè ge chéngshì bǐ nà ge chéngshì gānjìng yìxiē.}\\
« Cette ville est un peu plus propre que celle-là. »\\[4pt]
\zh{坐火车比坐汽车舒服得多。}\\
\emph{Zuò huǒchē bǐ zuò qìchē shūfu de duō.}\\
« Le train est beaucoup plus confortable que le car. »
\end{exemplebox}

\courssub{Exprimer un écart chiffré précis}

Au lieu d'un marqueur vague, on peut donner la \textbf{mesure exacte} de l'écart : nombre + mot de mesure (classifieur), toujours \emph{après} l'adjectif.

\begin{textebox}[Écart chiffré]
\zh{哥哥比我大两岁。}\\
\emph{Gēge bǐ wǒ dà liǎng suì.}\\
« Mon grand frère a deux ans de plus que moi. »\\[4pt]
\zh{这条鱼比那条贵十块。}\\
\emph{Zhè tiáo yú bǐ nà tiáo guì shí kuài.}\\
« Ce poisson coûte dix yuans de plus que celui-là. »
\end{textebox}

\begin{methode}[Construire une comparaison avec écart chiffré]
\begin{enumerate}
\item Poser la structure de base : A + 比 + B + adjectif (\zh{哥哥比我大}).
\item Choisir la mesure : âge → \zh{岁}, prix → \zh{块} / \zh{块钱}, taille → \zh{厘米}, poids → \zh{公斤}.
\item Placer « nombre + mesure » \textbf{après} l'adjectif : \zh{大两岁}, \zh{贵十块}, \zh{高三厘米}.
\item Ne jamais ajouter \zh{了} après un écart chiffré : on dit \zh{大两岁}, pas \zh{大两岁了} dans ce sens.
\end{enumerate}
\end{methode}

\courssub{Emploi avec les verbes psychologiques}

La nuance d'écart fonctionne aussi après les \textbf{verbes de sentiment et de préférence} (\zh{喜欢}, \zh{爱}, \zh{想}, \zh{怕}…) :

\zh{我比他喜欢运动。} \emph{Wǒ bǐ tā xǐhuan yùndòng.} « J'aime le sport plus que lui. »

\zh{他比我喜欢运动一点。} \emph{Tā bǐ wǒ xǐhuan yùndòng yìdiǎn.} « Il aime le sport un peu plus que moi. »

\zh{妹妹比我怕狗多了。} \emph{Mèimei bǐ wǒ pà gǒu duō le.} « Ma petite sœur a beaucoup plus peur des chiens que moi. »

Le marqueur \zh{一点} / \zh{多了} se place alors \textbf{après le verbe} (ou après son objet court), exactement comme après un adjectif. \textbf{Note :} l'emploi de \zh{更} ou \zh{还} (voir section suivante) renforce encore le degré.

\courssub{Comparaison avec le français}

\begin{center}
\begin{tabularx}{\linewidth}{|X|X|X|}
\hline
\rowcolor{popblueL} Structure & Exemple en caractères + pinyin & Traduction \\
\hline
A + 比 + B + adj. + 一点 & \zh{他比我高一点。} Tā bǐ wǒ gāo yìdiǎn. & Il est un peu plus grand que moi. \\
\hline
A + 比 + B + adj. + 一些 & \zh{今天比昨天冷一些。} Jīntiān bǐ zuótiān lěng yìxiē. & Il fait un peu plus froid qu'hier. \\
\hline
A + 比 + B + adj. + 多了 & \zh{这本书比那本贵多了。} Zhè běn shū bǐ nà běn guì duō le. & Ce livre est beaucoup plus cher que celui-là. \\
\hline
A + 比 + B + adj. + 得多 & \zh{火车比汽车快得多。} Huǒchē bǐ qìchē kuài de duō. & Le train est bien plus rapide que la voiture. \\
\hline
A + 比 + B + adj. + nombre + mesure & \zh{哥哥比我大两岁。} Gēge bǐ wǒ dà liǎng suì. & Mon frère a deux ans de plus que moi. \\
\hline
A + 比 + B + verbe psych. + 一点/多了 & \zh{我比他喜欢足球多了。} Wǒ bǐ tā xǐhuan zúqiú duō le. & J'aime le football bien plus que lui. \\
\hline
\end{tabularx}
\end{center}

\textbf{Différence majeure avec le français :} en français, « un peu » et « beaucoup » se placent \emph{avant} l'adjectif (« un peu plus grand », « beaucoup plus cher »). En chinois, c'est l'inverse : le marqueur vient \emph{après} l'adjectif. L'ordre chinois est : « moi comparé à lui, grand — d'un peu ». Il faut donc penser la phrase dans cet ordre et résister au calque du français.

\begin{attention}[Erreurs fréquentes des francophones]
\begin{itemize}
\item Ne jamais dire \zh{*他比我一点高} : le marqueur d'écart ne précède \textbf{jamais} l'adjectif. Correct : \zh{他比我高一点。}
\item Ne pas mettre \zh{很} avant l'adjectif dans une comparaison : \zh{*他比我很高} est fautif ; \zh{比} suffit à marquer le degré.
\item Ne pas oublier le \zh{了} de \zh{多了} : \zh{贵多} seul n'existe pas ; on dit \zh{贵多了}.
\item Ne pas confondre \zh{得多} (comparaison, après l'adjectif) avec \zh{的} et \zh{地} : trois caractères différents, trois emplois différents.
\item Avec un écart chiffré, pas de \zh{了} : \zh{大两岁}, non \zh{*大两岁了}.
\end{itemize}
\end{attention}

\courssub{Exercice résolu de transformation}

\begin{exoresolu}[Ajouter 一点 puis 多了]
Énoncé : transformez la phrase \zh{苹果比香蕉贵。} \emph{Píngguǒ bǐ xiāngjiāo guì.} « Les pommes sont plus chères que les bananes » en ajoutant (a) \zh{一点}, puis (b) \zh{多了}.
\tcblower
Solution détaillée :
\begin{enumerate}
\item On repère l'adjectif : \zh{贵} « cher ».
\item (a) Écart faible : on place \zh{一点} après l'adjectif → \zh{苹果比香蕉贵一点。} \emph{Píngguǒ bǐ xiāngjiāo guì yìdiǎn.} « Les pommes sont un peu plus chères que les bananes. »
\item (b) Écart grand : on place \zh{多了} après l'adjectif → \zh{苹果比香蕉贵多了。} \emph{Píngguǒ bǐ xiāngjiāo guì duō le.} « Les pommes sont beaucoup plus chères que les bananes. »
\end{enumerate}
Vérification : dans les deux cas, le marqueur suit l'adjectif ; la structure A + 比 + B reste inchangée.
\end{exoresolu}

\begin{exercice}[À vous de jouer]
Transformez les phrases suivantes en ajoutant le marqueur indiqué, puis traduisez :
\begin{enumerate}
\item \zh{雅温得比贝尔图阿大。}(一点)
\item \zh{地铁比公共汽车快。}(多了)
\item \zh{这个房间比那个房间干净。}(一些)
\item \zh{弟弟比我高。}(écart chiffré : trois centimètres, \zh{三厘米})
\end{enumerate}
\end{exercice}

\begin{corrige}[Exercice « À vous de jouer »]
\begin{enumerate}
\item \zh{雅温得比贝尔图阿大一点。} \emph{Yǎwēndé bǐ Bèi'ěrtú'ā dà yìdiǎn.} « Yaoundé est un peu plus grande que Bertoua. »
\item \zh{地铁比公共汽车快多了。} \emph{Dìtiě bǐ gōnggòng qìchē kuài duō le.} « Le métro est beaucoup plus rapide que le bus. »
\item \zh{这个房间比那个房间干净一些。} \emph{Zhè ge fángjiān bǐ nà ge fángjiān gānjìng yìxiē.} « Cette chambre est un peu plus propre que celle-là. »
\item \zh{弟弟比我高三厘米。} \emph{Dìdi bǐ wǒ gāo sān límǐ.} « Mon petit frère est plus grand que moi de trois centimètres. »
\end{enumerate}
\end{corrige}

\begin{aretenir}[L'essentiel de la section]
\begin{itemize}
\item Structure : \textbf{A + 比 + B + adjectif + marqueur d'écart}.
\item Écart faible : \zh{一点} (courant), \zh{一些} (plus soutenu). Écart grand : \zh{多了} (oral), \zh{得多} (écrit).
\item Le marqueur se place \textbf{toujours après} l'adjectif — contrairement au français.
\item Écart chiffré : adjectif + nombre + mesure (\zh{大两岁}, \zh{贵十块}), sans \zh{了}.
\item La nuance fonctionne aussi après les verbes psychologiques : \zh{我比他喜欢运动一点。}
\end{itemize}
\end{aretenir}

\coursec{Renforcer la comparaison : A比B更/还 + adjectif}

Dans la section précédente, nous avons appris à \cle{nuancer l'écart} entre deux éléments : grâce à \zh{一点} (yīdiǎn, « un peu ») et \zh{多了} (duō le, « beaucoup ») placés \emph{après} l'adjectif, nous pouvons dire si la différence est petite ou grande. Mais il existe une autre façon d'enrichir la comparaison : placer un adverbe \emph{avant} l'adjectif pour \cle{renforcer} le constat. C'est ce que nous allons découvrir avec \zh{更} (gèng) et \zh{还} (hái), qui correspondent au français « encore plus », « encore mieux », « encore moins cher ».

\courssub{Observation : deux phrases à comparer}

Lisons attentivement les deux phrases suivantes et comparons-les avec la structure de base \zh{A + 比 + B + adjectif}.

\begin{exemplebox}[Observation guidée]
\zh{这个比那个更好。} \emph{Zhège bǐ nàge gèng hǎo.} « Celui-ci est encore meilleur que celui-là. »\\[4pt]
\zh{他的汉语比我的还好。} \emph{Tā de Hànyǔ bǐ wǒ de hái hǎo.} « Son chinois est encore meilleur que le mien. »
\end{exemplebox}

Questions d'observation :
\begin{enumerate}
\item Où sont placés \zh{更} et \zh{还} par rapport à l'adjectif \zh{好} ? (Réponse : juste \emph{avant} l'adjectif, et \emph{après} B.)
\item Que devient le sens de la phrase ? La différence entre A et B est-elle simplement constatée, ou est-elle \emph{soulignée, accentuée} ?
\item Peut-on remplacer \zh{更} ou \zh{还} par \zh{很} ? Essayons : \zh{*这个比那个很好。} — cette phrase est \textbf{incorrecte}.
\end{enumerate}

Nous observons donc que \zh{更} et \zh{还} occupent exactement la place où un francophone serait tenté de mettre \zh{很} (« très »). C'est un point capital que nous allons étudier en détail.

\courssub{La règle : schéma, forme et emploi}

\begin{propriete}[Schéma de la comparaison renforcée]
\textbf{A + 比 + B + 更/还 + adjectif}\\[4pt]
\zh{更} et \zh{还} renforcent l'adjectif dans une comparaison avec \zh{比} : ils signifient « encore plus », « encore mieux ». \zh{更} est un renforcement \textbf{neutre}, utilisable à l'oral comme à l'écrit ; \zh{还} ajoute souvent une nuance d'\textbf{étonnement} ou d'inattendu et appartient plutôt à la langue orale.
\end{propriete}

\begin{definition}[更 gèng et 还 hái]
\zh{更} (gèng) : adverbe de degré, « encore plus, davantage ». Il insiste sur l'intensité de la différence, sans jugement particulier. Registre neutre, écrit et oral.\\[4pt]
\zh{还} (hái) : adverbe, « encore, encore plus ». Dans une comparaison, il exprime souvent que le résultat \emph{surprend}, qu'on ne s'y attendait pas. Registre plutôt oral.
\end{definition}

\begin{popfigure}
\begin{center}
\begin{tikzpicture}[node distance=0.4cm, every node/.style={font=\small}]
\node[draw=popblue, fill=popblueL, rounded corners, minimum width=1.5cm, minimum height=0.8cm] (A) {A};
\node[draw=popblue, fill=popblueL, rounded corners, minimum width=1.5cm, minimum height=0.8cm, right=of A] (bi) {比};
\node[draw=popblue, fill=popblueL, rounded corners, minimum width=1.5cm, minimum height=0.8cm, right=of bi] (B) {B};
\node[draw=popgreen, fill=popgreenL, rounded corners, minimum width=2.2cm, minimum height=0.8cm, right=of B, align=center] (adv) {更 / 还};
\node[draw=poppurple, fill=poppurpleL, rounded corners, minimum width=2.2cm, minimum height=0.8cm, right=of adv, align=center] (adj) {adjectif};
\draw[-{Stealth}, thick, popdark] (A) -- (bi);
\draw[-{Stealth}, thick, popdark] (bi) -- (B);
\draw[-{Stealth}, thick, popdark] (B) -- (adv);
\draw[-{Stealth}, thick, popdark] (adv) -- (adj);
\node[draw=popgreen, fill=popgreenL, rounded corners, below=1.1cm of adv, align=center, inner sep=5pt] (ok) {\textbf{更, 还 = AUTORISÉS}\\ (ils renforcent)};
\node[draw=poporange, fill=poporangeL, rounded corners, left=0.6cm of ok, align=center, inner sep=5pt] (no) {\textbf{很 = INTERDIT}\\ (avec 比)};
\draw[-{Stealth}, thick, popgreen] (adv) -- (ok);
\end{tikzpicture}
\legende{Schéma de la structure A + 比 + B + 更/还 + adjectif : avant l'adjectif, seuls 更 et 还 sont possibles ; 很 est interdit.}
\end{center}
\end{popfigure}

\courssub{更 ou 还 ? Le choix de la nuance}

Les deux adverbes renforcent la comparaison, mais ils ne sont pas toujours interchangeables. Le tableau suivant résume la différence.

\begin{center}
\begin{tabularx}{\linewidth}{|X|X|X|}
\hline
\rowcolor{popblueL} \textbf{Critère} & \textbf{更 gèng} & \textbf{还 hái} \\
\hline
Sens & encore plus, davantage & encore plus (étonnamment) \\
\hline
Nuance & neutre, constat objectif & surprise, inattendu \\
\hline
Registre & oral et écrit & plutôt oral \\
\hline
Exemple & \zh{今天比昨天更热。} \emph{Jīntiān bǐ zuótiān gèng rè.} & \zh{今天比昨天还热。} \emph{Jīntiān bǐ zuótiān hái rè.} \\
\hline
Traduction & Aujourd'hui, il fait encore plus chaud qu'hier. & Aujourd'hui, il fait encore plus chaud qu'hier (incroyable !) \\
\hline
\end{tabularx}
\end{center}

\begin{exemplebox}[Exemples variés avec 更]
\zh{雅温得的交通比杜阿拉的更方便。} \emph{Yǎwēndé de jiāotōng bǐ Dù'ālā de gèng fāngbiàn.} « La circulation à Yaoundé est encore plus pratique qu'à Douala. »\\[4pt]
\zh{医生说我应该比过去更注意健康。} \emph{Yīshēng shuō wǒ yīnggāi bǐ guòqù gèng zhùyì jiànkāng.} « Le médecin dit que je dois faire encore plus attention à ma santé qu'avant. »\\[4pt]
\zh{这个问题比那个问题更重要。} \emph{Zhège wèntí bǐ nàge wèntí gèng zhòngyào.} « Cette question est encore plus importante que celle-là. »
\end{exemplebox}

\begin{exemplebox}[Exemples variés avec 还]
\zh{弟弟比哥哥还高。} \emph{Dìdi bǐ gēge hái gāo.} « Le petit frère est encore plus grand que le grand frère ! » (surprenant)\\[4pt]
\zh{这个市场比那个市场还热闹。} \emph{Zhège shìchǎng bǐ nàge shìchǎng hái rènao.} « Ce marché est encore plus animé que celui-là. »\\[4pt]
\zh{他跑得比我还快。} \emph{Tā pǎo de bǐ wǒ hái kuài.} « Il court encore plus vite que moi ! »
\end{exemplebox}

\courssub{Texte support : au marché}

\begin{textebox}[Au marché de Mokolo — 在莫科洛市场]
\zh{南方的水果比北方的还便宜。} \emph{Nánfāng de shuǐguǒ bǐ běifāng de hái piányi.} « Les fruits du Sud sont encore moins chers que ceux du Nord. »\\[4pt]
\zh{这里的芒果比超市的更新鲜，也更便宜。} \emph{Zhèlǐ de mángguǒ bǐ chāoshì de gèng xīnxiān, yě gèng piányi.} « Les mangues d'ici sont encore plus fraîches que celles du supermarché, et encore moins chères aussi. »\\[4pt]
\zh{妈妈说，为了健康，我们应该吃更多水果。} \emph{Māma shuō, wèile jiànkāng, wǒmen yīnggāi chī gèng duō shuǐguǒ.} « Maman dit que pour la santé, nous devrions manger encore plus de fruits. »
\end{textebox}

\begin{center}
\begin{tabularx}{\linewidth}{|X|X|X|X|}
\hline
\rowcolor{popblueL} \textbf{Caractères} & \textbf{Pinyin} & \textbf{Nature} & \textbf{Traduction} \\
\hline
南方 & nánfāng & nom & le Sud \\
\hline
北方 & běifāng & nom & le Nord \\
\hline
新鲜 & xīnxiān & adjectif & frais \\
\hline
超市 & chāoshì & nom & supermarché \\
\hline
芒果 & mángguǒ & nom & mangue \\
\hline
\end{tabularx}
\end{center}

Questions de compréhension :
\begin{enumerate}
\item Que compare la première phrase ? Quel adverbe de renforcement est utilisé ?
\item Dans la deuxième phrase, combien y a-t-il de comparaisons renforcées ?
\item Retrouvez dans la troisième phrase un emploi de \zh{更} \emph{sans} \zh{比} : que signifie-t-il ?
\end{enumerate}

\courssub{Combiner les deux nuances : 更/还 + adjectif + 一点/多了}

Rien n'empêche de combiner ce que nous avons appris dans la section précédente avec la structure d'aujourd'hui. On obtient une \cle{double nuance} : le renforcement avant l'adjectif (\zh{更}/\zh{还}) et la mesure de l'écart après l'adjectif (\zh{一点}/\zh{多了}).

\begin{propriete}[Combinaison possible]
\textbf{A + 比 + B + 更/还 + adjectif + 一点/多了}\\[4pt]
\zh{这本书比那本更厚一点。} \emph{Zhè běn shū bǐ nà běn gèng hòu yīdiǎn.} « Ce livre est encore un peu plus épais que celui-là. »\\[4pt]
\zh{他的成绩比我的还好多了。} \emph{Tā de chéngjì bǐ wǒ de hái hǎo duō le.} « Ses résultats sont encore bien meilleurs que les miens. »
\end{propriete}

Attention à l'ordre : \zh{更}/\zh{还} se placent \textbf{avant} l'adjectif, \zh{一点}/\zh{多了} \textbf{après}. Ne jamais inverser : \zh{*这本书比那本一点更厚} est incorrect.

\courssub{更 sans comparaison explicite : un emploi utile pour l'examen}

\zh{更} peut aussi s'employer \textbf{seul}, sans \zh{比}, quand la comparaison est implicite (comparaison avec le passé, avec une situation antérieure). C'est un emploi fréquent dans les sujets d'examen.

\begin{exemplebox}[更 hors comparaison explicite]
\zh{我想更健康。} \emph{Wǒ xiǎng gèng jiànkāng.} « Je veux être en meilleure santé (qu'avant). »\\[4pt]
\zh{我们应该更努力学习。} \emph{Wǒmen yīnggāi gèng nǔlì xuéxí.} « Nous devons étudier encore plus sérieusement. »\\[4pt]
\zh{吃了药以后，他更舒服了。} \emph{Chī le yào yǐhòu, tā gèng shūfu le.} « Après avoir pris le médicament, il se sent mieux. »
\end{exemplebox}

Dans ces phrases, le français traduit naturellement par « mieux », « davantage » : la comparaison est sous-entendue. Notez bien qu'ici aussi, \zh{很} ne conviendrait pas, car le sens n'est pas « très » mais « plus qu'avant ».

\courssub{Comparaison avec le français et pièges des francophones}

En français, on dit « encore \emph{plus} grand », « encore \emph{mieux} » : l'adverbe de renforcement se place avant l'adjectif, comme en chinois. La structure est donc facile à comprendre. Le vrai piège est ailleurs : le francophone apprend très tôt que l'adjectif chinois s'accompagne presque toujours de \zh{很}, et il a tendance à le mettre partout.

\begin{attention}[很 ou 更 ? Le piège classique]
Les francophones hésitent entre \zh{很} et \zh{更}. Retenez la règle absolue : \textbf{dans une phrase avec 比, seul 更/还 est possible, jamais 很}.\\[4pt]
\zh{*他比我很高。} \textbf{FAUX} — \zh{他比我更高。} \textbf{JUSTE} « Il est encore plus grand que moi. »\\[4pt]
Pourquoi ? Parce que \zh{很} exprime un degré absolu (« très »), alors que \zh{比} exprime un degré relatif (« plus... que »). Les deux logiques sont incompatibles. \zh{更} et \zh{还}, eux, renforcent le degré relatif : ils sont donc parfaitement compatibles avec \zh{比}.
\end{attention}

Autres erreurs fréquentes :
\begin{itemize}
\item Placer \zh{更} avant \zh{比} : \zh{*他更比我高。} est faux ; \zh{更} se place après B, juste avant l'adjectif.
\item Oublier que \zh{还} se lit ici \emph{hái} et non \emph{huán}.
\item Traduire « meilleur » par \zh{更好} hors comparaison : c'est correct, mais avec \zh{比} on dira \zh{比……更好} et non \zh{*比……很好}.
\end{itemize}

\begin{savaistu}[还 : deux prononciations, deux sens]
Le caractère \zh{还} est un caractère \cle{polyphonique} : il possède deux lectures. Dans notre leçon, il se prononce \textbf{hái} et signifie « encore » (\zh{还好}, \zh{还有} « il y a encore »). Mais il se prononce aussi \textbf{huán} avec le sens de « rendre, restituer » : \zh{还钱} (huán qián) signifie « rendre de l'argent », \zh{还书} (huán shū) « rendre un livre » à la bibliothèque. Le contexte et le vocabulaire permettent de choisir la bonne lecture : devant un verbe d'action concrète comme « rendre », lisez huán ; dans le sens temporel ou comparatif, lisez hái.
\end{savaistu}

\courssub{Tableau récapitulatif}

\begin{center}
\begin{tabularx}{\linewidth}{|X|X|X|}
\hline
\rowcolor{popblueL} \textbf{Structure} & \textbf{Exemple (caractères + pinyin)} & \textbf{Traduction} \\
\hline
A + 比 + B + 更 + adj & \zh{今天比昨天更冷。} \emph{Jīntiān bǐ zuótiān gèng lěng.} & Il fait encore plus froid aujourd'hui qu'hier. \\
\hline
A + 比 + B + 还 + adj & \zh{他比我还忙。} \emph{Tā bǐ wǒ hái máng.} & Il est encore plus occupé que moi ! \\
\hline
A + 比 + B + 更/还 + adj + 一点 & \zh{这条路比那条更宽一点。} \emph{Zhè tiáo lù bǐ nà tiáo gèng kuān yīdiǎn.} & Cette route est encore un peu plus large que celle-là. \\
\hline
A + 比 + B + 更/还 + adj + 多了 & \zh{汉语比英语难多了，但更有意思。} \emph{Hànyǔ bǐ Yīngyǔ nán duō le, dàn gèng yǒu yìsi.} & Le chinois est bien plus difficile que l'anglais, mais encore plus intéressant. \\
\hline
更 + adj (sans 比) & \zh{我想更健康。} \emph{Wǒ xiǎng gèng jiànkāng.} & Je veux être en meilleure santé. \\
\hline
\zh{*A + 比 + B + 很 + adj} & \zh{*他比我很好。} & PHRASE INTERDITE \\
\hline
\end{tabularx}
\end{center}

\courssub{Exercice d'application}

\begin{exoresolu}[Réécrire avec 更 ou 还]
Réécrivez les phrases suivantes en ajoutant \zh{更} ou \zh{还} selon le contexte indiqué entre parenthèses.\\[4pt]
1. \zh{杜阿拉比雅温得热。}(constat neutre)\\[2pt]
2. \zh{妹妹比我高。}(quel étonnement !)\\[2pt]
3. \zh{这家医院比那家好。}(constat neutre, et insister sur l'écart : beaucoup plus)\\[2pt]
4. \zh{他的发音比我的好。}(surprise à l'oral)
\tcblower
\textbf{Solution détaillée.}\\[4pt]
1. Constat neutre, à l'écrit : on choisit \zh{更}. \zh{杜阿拉比雅温得更热。} \emph{Dù'ālā bǐ Yǎwēndé gèng rè.} « Douala est encore plus chaude que Yaoundé. »\\[4pt]
2. Nuance d'étonnement : on choisit \zh{还}. \zh{妹妹比我还高。} \emph{Mèimei bǐ wǒ hái gāo.} « Ma petite sœur est encore plus grande que moi ! »\\[4pt]
3. Renforcement neutre + écart grand : combinaison \zh{更……多了}. \zh{这家医院比那家更好多了。} \emph{Zhè jiā yīyuàn bǐ nà jiā gèng hǎo duō le.} « Cet hôpital est encore bien meilleur que celui-là. »\\[4pt]
4. Surprise à l'oral : \zh{还}. \zh{他的发音比我的还好。} \emph{Tā de fāyīn bǐ wǒ de hái hǎo.} « Sa prononciation est encore meilleure que la mienne ! »
\end{exoresolu}

\begin{exercice}[À vous de jouer]
Traduisez en chinois avec \zh{更} ou \zh{还} :
\begin{enumerate}
\item Le ndolé de ce restaurant est encore meilleur que celui de l'autre restaurant. (neutre)
\item Le bus est encore moins cher que le taxi ! (surprise)
\item Après le sport, je me sens encore mieux qu'avant. (neutre)
\item Ce match de football est encore plus intéressant que le précédent, et de beaucoup. (combiner avec \zh{多了})
\end{enumerate}
\end{exercice}

\begin{corrige}[Exercice « À vous de jouer »]
1. \zh{这家餐馆的恩多莱比那家的更好吃。} \emph{Zhè jiā cānguǎn de ndolé bǐ nà jiā de gèng hǎochī.}\\[2pt]
2. \zh{公共汽车比出租车还便宜。} \emph{Gōnggòng qìchē bǐ chūzūchē hái piányi.}\\[2pt]
3. \zh{运动以后，我感觉比以前更舒服。} \emph{Yùndòng yǐhòu, wǒ gǎnjué bǐ yǐqián gèng shūfu.}\\[2pt]
4. \zh{这场足球比赛比上一场的更有意思多了。} \emph{Zhè chǎng zúqiú bǐsài bǐ shàng yī chǎng de gèng yǒu yìsi duō le.}
\end{corrige}

\begin{experience}[Expression orale : Comparer pour choisir]
\textbf{Consigne :} En binôme, imaginez un dialogue entre deux lycéens qui hésitent entre deux options pour le week-end (ex: aller au stade vs aller au marché Mokolo ; manger du ndolé vs manger du poulet DG ; réviser le chinois vs réviser les maths).\\[4pt]
Utilisez au moins 3 structures vues aujourd'hui :\\[2pt]
- \zh{A比B更/还+adj} (choix neutre ou surpris)\\[2pt]
- \zh{A比B更/还+adj+一点/多了} (double nuance)\\[2pt]
- \zh{更+adj} sans \zh{比} (projet, résolution)\\[4pt]
\textbf{Exemple de début :}\\[2pt]
\zh{A: 这个周末去踢球比去市场更有意思。}\\[2pt]
\zh{B: 可是市场比球场还热闹，而且东西更便宜多了。}\\[2pt]
\zh{A: 那我们要更早起床。}
\end{experience}

\begin{aretenir}[L'essentiel de la section]
\begin{itemize}
\item Schéma : \textbf{A + 比 + B + 更/还 + adjectif} — « A est encore plus... que B ».
\item \zh{更} (gèng) : renforcement neutre, oral et écrit ; \zh{还} (hái) : renforcement avec surprise, plutôt oral.
\item Avec \zh{比}, \zh{很} est \textbf{interdit} ; seuls \zh{更} et \zh{还} peuvent précéder l'adjectif.
\item Combinaison possible : \zh{更/还} avant l'adjectif + \zh{一点/多了} après.
\item \zh{更} peut s'employer sans \zh{比} quand la comparaison est implicite : \zh{我想更健康。}
\item \zh{还} a deux lectures : \emph{hái} (encore) et \emph{huán} (rendre).
\end{itemize}
\end{aretenir}

\coursec{La comparaison négative : A不如B (+ adjectif)}

Dans la section précédente, nous avons appris à \emph{renforcer} une comparaison : \zh{A比B更/还 + adjectif} permet de dire « A est encore plus… que B ». Mais la vie n'est pas faite que de supériorités ! Comment dire, en chinois, « je suis \emph{moins} grand que lui », « ma santé est \emph{moins bonne} qu'avant », « la moto est \emph{moins sûre} que la voiture » ? Le français utilise « moins… que » ; le chinois, lui, change de verbe : il abandonne \zh{比} et emploie \zh{不如} (bùrú), littéralement « ne pas égaler, ne pas atteindre ». C'est l'objet de cette section : la comparaison négative.

\courssub{Observation : découvrir 不如 en contexte}

\begin{textebox}[Au dispensaire de Yaoundé — dialogue]
\zh{医生：你最近身体怎么样？}\\
\emph{Yīshēng : Nǐ zuìjìn shēntǐ zěnmeyàng ?}\\
« Docteur : Comment va ta santé ces derniers temps ? »\\[4pt]
\zh{病人：说实话，我的身体不如以前好。}\\
\emph{Bìngrén : Shuō shíhuà, wǒ de shēntǐ bùrú yǐqián hǎo.}\\
« Patient : Pour être honnête, ma santé n'est plus aussi bonne qu'avant. »\\[4pt]
\zh{医生：你以前每天运动吗？}\\
\emph{Yīshēng : Nǐ yǐqián měitiān yùndòng ma ?}\\
« Docteur : Faisais-tu du sport tous les jours avant ? »\\[4pt]
\zh{病人：以前我常常跑步，现在工作太忙，运动的时间不如以前多。}\\
\emph{Bìngrén : Yǐqián wǒ chángcháng pǎobù, xiànzài gōngzuò tài máng, yùndòng de shíjiān bùrú yǐqián duō.}\\
« Patient : Avant je courais souvent ; maintenant le travail me prend trop de temps, et je fais moins de sport qu'avant. »\\[4pt]
\zh{医生：坐摩托车上班不如坐汽车安全，你也要注意。}\\
\emph{Yīshēng : Zuò mótuōchē shàngbān bùrú zuò qìchē ānquán, nǐ yě yào zhùyì.}\\
« Docteur : Aller au travail en moto est moins sûr qu'en voiture, fais aussi attention à cela. »\\[4pt]
\zh{病人：您说得对。为了健康，我要多运动。}\\
\emph{Bìngrén : Nín shuō de duì. Wèile jiànkāng, wǒ yào duō yùndòng.}\\
« Patient : Vous avez raison. Pour ma santé, je dois faire plus de sport. »
\end{textebox}

\textbf{Questions de compréhension.}
\begin{enumerate}
\item Que dit le patient de sa santé actuelle par rapport au passé ?
\item Pourquoi fait-il moins de sport qu'avant ?
\item Selon le médecin, quel moyen de transport est le plus sûr ?
\end{enumerate}

\textbf{Observons.} Relevons les trois phrases clés du dialogue :

\begin{exemplebox}[Trois exemples à observer]
\zh{我的身体不如以前好。} \emph{Wǒ de shēntǐ bùrú yǐqián hǎo.} « Ma santé n'est plus aussi bonne qu'avant. »\\[4pt]
\zh{我不如他高。} \emph{Wǒ bùrú tā gāo.} « Je ne suis pas aussi grand que lui. » (= il est plus grand que moi)\\[4pt]
\zh{今天不如昨天热。} \emph{Jīntiān bùrú zuótiān rè.} « Aujourd'hui, il fait moins chaud qu'hier. »
\end{exemplebox}

Dans les trois cas, le premier élément (A) est \cle{inférieur} au second (B) sur le point comparé : ma santé $<$ avant ; moi $<$ lui (en taille) ; aujourd'hui $<$ hier (en chaleur). Le mot qui porte cette relation est \zh{不如}.

\begin{definition}[不如 bùrú]
\zh{不如} est un verbe qui signifie « ne pas égaler, ne pas valoir, être inférieur à ». Il sert à construire la comparaison d'\emph{infériorité} : « A n'atteint pas B », « A est moins… que B ». Attention : en pinyin standard, on écrit \textbf{bùrú} (ton 4 + ton 2), mais à l'oral, le \zh{不} (bù, ton 4) change en ton montant (bú, ton 2) devant le ton 2 de \zh{如} (rú) — on prononce donc \emph{búrú}.
\end{definition}

\courssub{Schéma 1 : A + 不如 + B (comparaison globale)}

Sans adjectif, \zh{不如} compare deux éléments \emph{globalement} : « A ne vaut pas B », « A n'est pas à la hauteur de B ».

\begin{propriete}[Structure : A + 不如 + B]
\textbf{Schéma :} \zh{A + 不如 + B}\\[4pt]
Emploi : A est jugé globalement inférieur à B (qualité, niveau, valeur).\\[4pt]
\zh{我的法语不如他的。} \emph{Wǒ de Fǎyǔ bùrú tā de.} « Mon français ne vaut pas le sien. »\\[4pt]
\zh{这个手机不如那个。} \emph{Zhège shǒujī bùrú nàge.} « Ce téléphone-ci ne vaut pas celui-là. »\\[4pt]
\zh{我不如你。} \emph{Wǒ bùrú nǐ.} « Je ne te vaux pas. » (je ne suis pas aussi bon que toi)
\end{propriete}

Remarquez le \zh{的} final dans \zh{他的} : il remplace le nom déjà connu (« son français »), exactement comme le pronom possessif français « le sien ».

\courssub{Schéma 2 : A + 不如 + B + adjectif (comparaison précise)}

Pour préciser \emph{sur quel point} A est inférieur à B, on ajoute un adjectif après B.

\begin{propriete}[Structure : A + 不如 + B + adjectif]
\textbf{Schéma :} \zh{A + 不如 + B + Adj.}\\[4pt]
Sens : « A est moins… que B », « A n'est pas aussi… que B ».\\[4pt]
\zh{坐摩托车不如坐汽车安全。} \emph{Zuò mótuōchē bùrú zuò qìchē ānquán.} « La moto est moins sûre que la voiture. »\\[4pt]
\zh{弟弟不如哥哥高。} \emph{Dìdi bùrú gēge gāo.} « Le petit frère est moins grand que le grand frère. »\\[4pt]
\zh{乡下的空气不如以前干净。} \emph{Xiāngxià de kōngqì bùrú yǐqián gānjìng.} « L'air de la campagne est moins pur qu'avant. »
\end{propriete}

\begin{attention}[Pas de 一点 / 多了 avec 不如 !]
Avec \zh{比}, on peut nuancer l'écart : \zh{他比我高一点} (un peu plus grand), \zh{他比我高多了} (beaucoup plus grand). Avec \zh{不如}, on \textbf{ne peut pas} ajouter \zh{一点} ou \zh{多了} après l'adjectif : on ne dit jamais \zh{我不如他高一点}. En revanche — et c'est l'inverse de \zh{比} — l'adjectif \textbf{peut être précédé} de \zh{这么} ou \zh{那么} :\\[4pt]
\zh{我不如他那么努力。} \emph{Wǒ bùrú tā nàme nǔlì.} « Je ne suis pas aussi travailleur que lui. »\\[4pt]
Rappel : avec \zh{比}, \zh{这么/那么} est interdit (on ne dit pas \zh{他比我那么高}).
\end{attention}

\courssub{Équivalence de sens : 不如 et 比, deux points de vue sur une même réalité}

Voici le point capital de cette leçon. Les deux phrases suivantes décrivent \emph{la même réalité} :

\begin{exemplebox}[Même réalité, deux formulations]
\zh{我不如他高。} \emph{Wǒ bùrú tā gāo.} « Je suis moins grand que lui. »\\[4pt]
\zh{他比我高。} \emph{Tā bǐ wǒ gāo.} « Il est plus grand que moi. »
\end{exemplebox}

La différence n'est pas dans les faits, mais dans la \cle{perspective} : avec \zh{比}, on part de l'élément \emph{fort} (celui qui gagne la comparaison) ; avec \zh{不如}, on part de l'élément \emph{faible} (celui qui perd).

\begin{popfigure}
\begin{tikzpicture}[>=Stealth, node distance=1.2cm]
\node[draw=popblue, fill=popblueL, rounded corners, align=center, text width=4.6cm, minimum height=1.6cm] (buru) at (0,0) {\zh{我不如他高。}\\ \emph{point de vue du « faible »}\\ (moi, moins grand)};
\node[draw=poporange, fill=poporangeL, rounded corners, align=center, text width=4.6cm, minimum height=1.6cm] (bi) at (8.2,0) {\zh{他比我高。}\\ \emph{point de vue du « fort »}\\ (lui, plus grand)};
\draw[<->, very thick, poppurple] (buru) -- (bi) node[midway, above, align=center, text=popdark]{\textbf{même réalité}\\ taille : lui $>$ moi};
\node[align=center, text=popmuted] at (4.1,-1.6) {\zh{A不如B + adj.} \; $\Leftrightarrow$ \; \zh{B比A + adj.}};
\end{tikzpicture}
\legende{Double flèche d'équivalence : A不如B+adj. et B比A+adj. décrivent le même fait, avec des perspectives inversées.}
\end{popfigure}

\begin{propriete}[Règle d'équivalence]
\zh{A + 不如 + B + Adj.} $\Leftrightarrow$ \zh{B + 比 + A + Adj.}\\[4pt]
\zh{不如} est l'\emph{inverse logique} de \zh{比} : il signifie « A n'atteint pas B ». Pour transformer une phrase, il suffit d'\textbf{échanger A et B} et de remplacer \zh{比} par \zh{不如} (ou l'inverse), \textbf{sans toucher à l'adjectif}.\\[4pt]
\zh{今天不如昨天热。} $\Leftrightarrow$ \zh{昨天比今天热。} « Hier il faisait plus chaud qu'aujourd'hui. »
\end{propriete}

Quand choisir l'une ou l'autre ? Suivez l'organigramme :

\begin{popfigure}
\begin{tikzpicture}[>=Stealth, node distance=1.1cm]
\node[draw=popdark, fill=popcream, rounded corners, align=center, text width=6.4cm] (q) at (0,0) {\textbf{De quel élément est-ce que je pars ?}};
\node[draw=poporange, fill=poporangeL, rounded corners, align=center, text width=4.4cm] (fort) at (-3.4,-2.2) {\textbf{Élément fort}\\ (le « gagnant »)\\ \zh{A比B + adj.}\\ \zh{他比我高。}};
\node[draw=popblue, fill=popblueL, rounded corners, align=center, text width=4.4cm] (faible) at (3.4,-2.2) {\textbf{Élément faible}\\ (le « perdant »)\\ \zh{A不如B + adj.}\\ \zh{我不如他高。}};
\draw[->, thick, poporange] (q.south west) -- (fort.north) node[midway, left, text=popdark]{le fort};
\draw[->, thick, popblue] (q.south east) -- (faible.north) node[midway, right, text=popdark]{le faible};
\end{tikzpicture}
\legende{Organigramme de choix : partir de l'élément fort → 比 ; partir de l'élément faible → 不如.}
\end{popfigure}

\courssub{Comparaison avec le français et tableau récapitulatif}

Le français exprime l'infériorité de deux façons : « moins… que » et « pas aussi… que ». Le chinois \zh{不如} couvre les deux : \zh{我不如他高} peut se traduire « je suis moins grand que lui » \emph{ou} « je ne suis pas aussi grand que lui ». Autre différence majeure : en français, « moins » se colle à l'adjectif (« moins grand ») ; en chinois, \zh{不如} se place \emph{avant B}, et l'adjectif reste intact après B.

\begin{center}
\begin{tabularx}{\linewidth}{|X|X|X|}
\hline
\rowcolor{popblueL} Structure & Exemple (caractères + pinyin) & Traduction \\
\hline
A + 不如 + B & \zh{我的法语不如他的。} Wǒ de Fǎyǔ bùrú tā de. & Mon français ne vaut pas le sien. \\
\hline
A + 不如 + B + Adj. & \zh{我不如他高。} Wǒ bùrú tā gāo. & Je suis moins grand que lui. \\
\hline
A + 不如 + B + 这么/那么 + Adj. & \zh{我不如他那么努力。} Wǒ bùrú tā nàme nǔlì. & Je ne suis pas aussi travailleur que lui. \\
\hline
Équivalence : B + 比 + A + Adj. & \zh{他比我高。} Tā bǐ wǒ gāo. & Il est plus grand que moi. \\
\hline
\end{tabularx}
\end{center}

\begin{center}
\begin{tabularx}{\linewidth}{|X|X|X|X|}
\hline
\rowcolor{popblueL} Caractères & Pinyin & Nature & Traduction \\
\hline
不如 & bùrú & verbe & ne pas égaler, être inférieur à \\
\hline
以前 & yǐqián & nom de temps & avant, autrefois \\
\hline
安全 & ānquán & adjectif & sûr, en sécurité \\
\hline
努力 & nǔlì & adjectif/verbe & travailleur, faire des efforts \\
\hline
干净 & gānjìng & adjectif & propre, pur \\
\hline
身体 & shēntǐ & nom & corps, santé \\
\hline
\end{tabularx}
\end{center}

\begin{attention}[不如 ≠ 不比 : le piège n° 1 des francophones]
Ne confondez jamais \zh{不如} et \zh{不比} : le sens change complètement.\\[4pt]
\zh{他不如我高。} \emph{Tā bùrú wǒ gāo.} = « Il est \textbf{moins} grand que moi. » (infériorité certaine)\\[4pt]
\zh{他不比我高。} \emph{Tā bù bǐ wǒ gāo.} = « Il n'est \textbf{pas plus} grand que moi. » (il est peut-être égal, peut-être plus petit — on nie la supériorité, on n'affirme pas l'infériorité)\\[4pt]
En français : « moins grand que » = \zh{不如…高} ; « pas plus grand que » = \zh{不比…高}. Traduire « il n'est pas plus grand que moi » par \zh{他不如我高} est une erreur de sens.
\end{attention}

\begin{savaistu}[La modestie, valeur cardinale en Chine]
En Chine, dire \zh{我不如你} (wǒ bùrú nǐ, « je ne te vaux pas ») n'est pas un signe de faiblesse : c'est une marque de \emph{modestie} très appréciée. Quand on félicite un Chinois pour ses résultats, il répond souvent \zh{哪里哪里，我不如你} (nǎli nǎli, wǒ bùrú nǐ — « mais non, mais non, je ne t'arrive pas à la cheville »). Au Cameroun, on dirait plutôt « tu exagères, merci quand même ! ». La structure \zh{不如} est donc très fréquente dans la politesse quotidienne chinoise : l'utiliser pour parler de soi, c'est parler comme un locuteur cultivé.
\end{savaistu}

\courssub{Exercice résolu : transformer 比 en 不如}

\begin{exoresolu}[Transformation : de 比 à 不如]
Énoncé. Transformez les phrases suivantes en phrases avec \zh{不如}, sans changer le sens :\\
1. \zh{他比我高。} 2. \zh{汽车比摩托车安全。} 3. \zh{昨天比今天热。} 4. \zh{哥哥比我努力。} 5. \zh{城市比乡下吵。}
\tcblower
Solution détaillée. Méthode : on échange A et B, on remplace \zh{比} par \zh{不如}, l'adjectif ne bouge pas.\\
1. \zh{我不如他高。} \emph{Wǒ bùrú tā gāo.} « Je suis moins grand que lui. »\\
2. \zh{坐摩托车不如坐汽车安全。} \emph{Zuò mótuōchē bùrú zuò qìchē ānquán.} « La moto est moins sûre que la voiture. »\\
3. \zh{今天不如昨天热。} \emph{Jīntiān bùrú zuótiān rè.} « Aujourd'hui il fait moins chaud qu'hier. »\\
4. \zh{我不如哥哥努力。} \emph{Wǒ bùrú gēge nǔlì.} « Je suis moins travailleur que mon grand frère. »\\
5. \zh{乡下不如城市吵。} \emph{Xiāngxià bùrú chéngshì chǎo.} « La campagne est moins bruyante que la ville. »
\end{exoresolu}

\begin{exercice}[À vous : transformation inverse]
Transformez ces phrases avec \zh{不如} en phrases avec \zh{比} :\\
1. \zh{我的汉语不如他的好。} 2. \zh{走路不如骑车快。} 3. \zh{这个房间不如那个干净。} 4. \zh{我不如妹妹聪明。} 5. \zh{今年不如去年热。}
\end{exercice}

\begin{corrige}[Exercice : transformation inverse]
1. \zh{他的汉语比我的好。} « Son chinois est meilleur que le mien. »\\
2. \zh{骑车比走路快。} « Le vélo est plus rapide que la marche. »\\
3. \zh{那个房间比这个干净。} « Cette chambre-là est plus propre que celle-ci. »\\
4. \zh{妹妹比我聪明。} « Ma petite sœur est plus intelligente que moi. »\\
5. \zh{去年比今年热。} « L'année dernière, il faisait plus chaud que cette année. »
\end{corrige}

\begin{aretenir}[L'essentiel de la section]
\begin{itemize}
\item \zh{A不如B} = « A ne vaut pas B » (comparaison globale) ; \zh{A不如B + adj.} = « A est moins… que B ».
\item Équivalence : \zh{A不如B + adj.} $\Leftrightarrow$ \zh{B比A + adj.} — on échange A et B, l'adjectif ne change pas.
\item Avec \zh{不如}, on peut dire \zh{那么/这么 + adj.} ; jamais \zh{一点/多了} après l'adjectif.
\item \zh{不如} (moins que) $\neq$ \zh{不比} (pas plus que) : attention au sens !
\item \zh{不如} est poli et atténué : parfait pour parler de soi avec modestie.
\end{itemize}
\end{aretenir}

\coursec{Exprimer le but avec 为了}

Dans la section précédente, nous avons appris à comparer en négatif avec \zh{不如} : dire que \zh{A不如B}, c'est reconnaître que A n'atteint pas le niveau de B. Nous quittons maintenant le domaine de la comparaison pour aborder une question essentielle de la communication : \textbf{exprimer le but}, c'est-à-dire répondre à la question « dans quel but ? », « pour quoi faire ? ». En chinois, la réponse passe presque toujours par un petit mot de deux caractères : \zh{为了} (wèile). Cette structure est omniprésente dans les discours sur l'environnement et la santé — le thème de notre module — car on y explique sans cesse \emph{pourquoi} on agit : pour protéger la nature, pour rester en bonne santé, pour réussir.

\courssub{Observation : lire le but dans des phrases authentiques}

Lisons d'abord deux phrases. Soulignons mentalement la partie qui exprime le but.

\begin{exemplebox}[Phrases d'observation]
\zh{为了健康，他每天运动。}\\
\emph{Wèile jiànkāng, tā měitiān yùndòng.}\\
« Pour sa santé, il fait du sport chaque jour. »

\medskip
\zh{为了通过考试，我们不但复习语法，而且练习口语。}\\
\emph{Wèile tōngguò kǎoshì, wǒmen búdàn fùxí yǔfǎ, érqiě liànxí kǒuyǔ.}\\
« Pour réussir l'examen, non seulement nous révisons la grammaire, mais en plus nous pratiquons l'oral. »
\end{exemplebox}

Que remarque-t-on ? Dans les deux phrases, \zh{为了} se trouve \textbf{au tout début}, suivi d'un groupe qui exprime le but (\zh{健康} « la santé », un nom ; \zh{通过考试} « réussir l'examen », une proposition verbale), puis d'une \textbf{virgule}, puis de l'action principale. La deuxième phrase combine élégamment \zh{为了} avec la structure \zh{不但……而且……} étudiée dans cette même leçon : le chinois enchaîne naturellement les outils grammaticaux dans une seule phrase.

Élargissons l'observation avec un texte court sur notre thème du module.

\begin{textebox}[Pour un environnement propre à Yaoundé]
\zh{为了保护环境，我们应该少用塑料袋。}\\
\emph{Wèile bǎohù huánjìng, wǒmen yīnggāi shǎo yòng sùliàodài.}\\
« Pour protéger l'environnement, nous devrions moins utiliser de sacs plastiques. »

\medskip
\zh{为了让城市更干净，雅温得的学生每个星期六打扫学校。}\\
\emph{Wèile ràng chéngshì gèng gānjìng, Yǎwēndé de xuésheng měi ge xīngqīliù dǎsǎo xuéxiào.}\\
« Pour rendre la ville plus propre, les élèves de Yaoundé nettoient leur école chaque samedi. »

\medskip
\zh{他为了身体健康，每天早上跑步，晚上不吃太多。}\\
\emph{Tā wèile shēntǐ jiànkāng, měitiān zǎoshang pǎobù, wǎnshang bù chī tài duō.}\\
« Pour sa santé, il court chaque matin et ne mange pas trop le soir. »
\end{textebox}

Troisième observation décisive : dans la dernière phrase, \zh{为了} n'est \textbf{plus en tête} : il vient \emph{après le sujet} \zh{他}. Les deux positions sont donc possibles. Précisons maintenant la règle.

\courssub{La règle : schéma, formes et positions}

\begin{propriete}[La structure 为了 + but]
\zh{为了} (wèile, « pour, afin de ») introduit le \cle{but} recherché, c'est-à-dire l'objectif que l'on veut atteindre en réalisant l'action. La proposition de but \textbf{précède toujours l'action principale}.

\medskip
\textbf{Schéma 1 (position la plus fréquente) :}\\
\zh{为了} + but，+ Sujet + Verbe (action)

\textbf{Schéma 2 (après le sujet) :}\\
Sujet + \zh{为了} + but + Verbe (action)

\medskip
Le but peut être :
\begin{itemize}
\item un \textbf{nom ou groupe nominal} : \zh{为了健康} (pour la santé), \zh{为了环境} (pour l'environnement), \zh{为了家人} (pour sa famille) ;
\item un \textbf{verbe ou une proposition} : \zh{为了考上好大学} (pour entrer dans une bonne université), \zh{为了找到好工作} (pour trouver un bon emploi).
\end{itemize}
\end{propriete}

\begin{exemplebox}[为了 + nom et 为了 + proposition]
\textbf{Avec un nom :}\\
\zh{为了孩子，父母努力工作。}\\
\emph{Wèile háizi, fùmǔ nǔlì gōngzuò.}\\
« Pour leurs enfants, les parents travaillent dur. »

\medskip
\textbf{Avec une proposition verbale :}\\
\zh{为了找到好工作，他学习汉语和英语。}\\
\emph{Wèile zhǎodào hǎo gōngzuò, tā xuéxí Hànyǔ hé Yīngyǔ.}\\
« Pour trouver un bon emploi, il apprend le chinois et l'anglais. »

\medskip
\textbf{Après le sujet :}\\
\zh{他为了学好汉语去了中国。}\\
\emph{Tā wèile xuéhǎo Hànyǔ qù le Zhōngguó.}\\
« Il est allé en Chine pour bien apprendre le chinois. »
\end{exemplebox}

\begin{attention}[Pas de verbe conjugué « pour » en fin de phrase]
En français, on dit « il fait du sport \emph{pour être} en bonne santé » : le but vient \textbf{après} l'action. En chinois, c'est impossible : le but exprimé par \zh{为了} doit venir \textbf{avant} l'action. On ne peut pas dire \zh{他每天运动，为了健康} dans un énoncé neutre : il faut inverser → \zh{为了健康，他每天运动。} ou \zh{他为了健康每天运动。} Retenez : en chinois, \textbf{le but d'abord, l'action ensuite}.
\end{attention}

\courssub{为了 (but) contre 因为 (cause) : l'opposition fondamentale}

C'est le point le plus délicat pour les francophones. Le français utilise des mots différents (« pour » / « parce que »), mais dans la rapidité de la production, les élèves les confondent souvent en chinois. Or la logique est inverse :

\begin{itemize}
\item \zh{因为} (yīnwèi, « parce que ») introduit la \cle{cause} : quelque chose qui \textbf{existe déjà} ou qui s'est \textbf{déjà produit}, et qui explique l'action. Il répond à la question « \emph{pourquoi ?} ».
\item \zh{为了} (wèile, « pour ») introduit le \cle{but} : quelque chose que l'on \textbf{veut obtenir}, orienté vers l'avenir. Il répond à la question « \emph{dans quel but ?} ».
\end{itemize}

\begin{exemplebox}[Même situation, deux logiques]
\zh{因为他生病了，所以没来上课。}\\
\emph{Yīnwèi tā shēngbìng le, suǒyǐ méi lái shàngkè.}\\
« Parce qu'il est malade (cause réelle, présente), il n'est pas venu en cours. »

\medskip
\zh{为了不生病，他每天吃水果。}\\
\emph{Wèile bù shēngbìng, tā měitiān chī shuǐguǒ.}\\
« Pour ne pas tomber malade (but visé, futur), il mange des fruits chaque jour. »
\end{exemplebox}

\begin{center}
\begin{tabularx}{\linewidth}{|X|X|X|}
\hline
\rowcolor{popblueL} Critère & \zh{为了} wèile (but) & \zh{因为} yīnwèi (cause) \\
\hline
Question posée & Dans quel but ? & Pourquoi ? \\
\hline
Orientation & Vers le futur (objectif) & Vers le présent/passé (fait) \\
\hline
Équivalent français & pour, afin de & parce que, car \\
\hline
Partenaire fréquent & seul en tête de phrase & souvent avec \zh{所以} (« donc ») \\
\hline
Exemple & \zh{为了健康，他运动。} & \zh{因为生病，他没来。} \\
\hline
Traduction & Pour sa santé, il fait du sport. & Parce qu'il est malade, il n'est pas venu. \\
\hline
\end{tabularx}
\end{center}

\begin{popfigure}
\begin{tikzpicture}[>=Stealth, node distance=1.1cm]
\node[draw=popgreen, fill=popgreenL, rounded corners, align=center, font=\small, minimum width=3.2cm, minimum height=1cm] (but) {为了 + BUT\\ \emph{objectif futur}};
\node[draw=popblue, fill=popblueL, rounded corners, align=center, font=\small, minimum width=3.2cm, minimum height=1cm, right=3.2cm of but] (act) {ACTION\\ \emph{ce qu'on fait}};
\draw[->, very thick, popgreen] (but) -- node[above, font=\small\itshape, text=popdark]{finalité} (act);
\node[draw=poporange, fill=poporangeL, rounded corners, align=center, font=\small, minimum width=3.2cm, minimum height=1cm, below=1.6cm of but] (cause) {因为 + CAUSE\\ \emph{fait présent/passé}};
\node[draw=poppurple, fill=poppurpleL, rounded corners, align=center, font=\small, minimum width=3.2cm, minimum height=1cm, below=1.6cm of act] (res) {RÉSULTAT\\ \emph{ce qui arrive}};
\draw[->, very thick, poporange] (cause) -- node[above, font=\small\itshape, text=popdark]{causalité} (res);
\end{tikzpicture}
\legende{Deux logiques opposées : 为了 regarde vers le but à atteindre ; 因为 regarde vers la cause déjà là.}
\end{popfigure}

\courssub{Carte mentale et vocabulaire utile}

\begin{popfigure}
\begin{tikzpicture}[mindmap, grow cyclic, every node/.style={concept, font=\small}, concept color=popblueL, level 1/.style={level distance=3.4cm, sibling angle=90}, level 1 concept/.append style={font=\small}]
\node[concept color=popblue!50, align=center]{为了\\ wèile\\ « pour »}
child[concept color=popgreenL]{ node[align=center]{+ nom\\ 为了健康\\ 为了环境} }
child[concept color=poporangeL]{ node[align=center]{+ proposition\\ 为了考上大学} }
child[concept color=poppurpleL]{ node[align=center]{position\\ en tête + ，\\ ou après le sujet} }
child[concept color=poppinkL]{ node[align=center]{contraire de\\ 因为 (cause)\\ but ≠ cause} };
\end{tikzpicture}
\legende{Carte mentale : tout ce qu'il faut savoir sur 为了.}
\end{popfigure}

\begin{center}
\begin{tabularx}{\linewidth}{|X|X|X|X|}
\hline
\rowcolor{popblueL} Caractères & Pinyin & Nature & Traduction \\
\hline
为了 & wèile & préposition & pour, afin de \\
\hline
因为 & yīnwèi & conjonction & parce que \\
\hline
目的 & mùdì & nom & but, objectif \\
\hline
健康 & jiànkāng & nom / adjectif & santé ; en bonne santé \\
\hline
保护 & bǎohù & verbe & protéger \\
\hline
环境 & huánjìng & nom & environnement \\
\hline
通过 & tōngguò & verbe & réussir (un examen), passer \\
\hline
考上 & kǎoshàng & verbe + résultat & être admis (à un examen) \\
\hline
努力 & nǔlì & verbe / adverbe & faire des efforts ; avec effort \\
\hline
实现 & shíxiàn & verbe & réaliser (un rêve, un but) \\
\hline
\end{tabularx}
\end{center}

\courssub{Comparaison avec le français et méthode de traduction}

Le français dispose de plusieurs formules : « pour + infinitif » (\emph{pour réussir}), « pour + nom » (\emph{pour sa santé}), « afin de », « dans le but de ». Le chinois est plus simple : \zh{为了} couvre tous ces emplois, suivi d'un nom ou d'une proposition complète. Deux différences à retenir :

\begin{enumerate}
\item \textbf{L'ordre} : le français place souvent le but après l'action (« il travaille \emph{pour réussir} ») ; le chinois place le but \textbf{avant} (\zh{为了成功，他努力工作。}).
\item \textbf{La virgule} : quand \zh{为了} ouvre la phrase, on met une virgule chinoise \zh{，} après le groupe de but ; quand il suit le sujet, la virgule est inutile.
\end{enumerate}

\begin{methode}[Traduire « pour + infinitif/nom » du français vers le chinois]
\begin{enumerate}
\item Repérer dans la phrase française l'expression du but : « pour », « afin de », « dans le but de ».
\item Vérifier qu'il s'agit bien d'un \textbf{but} (objectif visé) et non d'une \textbf{cause} : poser la question « dans quel but ? ». Si la réponse est « parce que », utiliser \zh{因为} à la place.
\item Traduire le but : nom seul (\zh{健康}) ou proposition verbale (\zh{考上好大学}).
\item Placer \zh{为了} + but \textbf{en tête de phrase}, suivi d'une virgule \zh{，}.
\item Traduire ensuite le sujet et l'action principale, dans l'ordre chinois habituel.
\end{enumerate}
\end{methode}

\begin{exoresolu}[Traduction guidée du français vers le chinois]
Énoncé : traduisez « Pour protéger la santé des enfants, l'école ne vend pas de boissons sucrées. »
\tcblower
Solution détaillée :
\begin{enumerate}
\item Le but est « protéger la santé des enfants » : c'est bien un objectif (question « dans quel but ? »), donc \zh{为了}.
\item Traduction du but : \zh{保护孩子们的健康} (bǎohù háizimen de jiànkāng).
\item Placement en tête + virgule : \zh{为了保护孩子们的健康，…}
\item Action principale : « l'école ne vend pas de boissons sucrées » → \zh{学校不卖甜的饮料} (xuéxiào bù mài tián de yǐnliào).
\end{enumerate}
Phrase complète : \zh{为了保护孩子们的健康，学校不卖甜的饮料。}\\
\emph{Wèile bǎohù háizimen de jiànkāng, xuéxiào bù mài tián de yǐnliào.}
\end{exoresolu}

\courssub{Exercices d'application}

\begin{exercice}[为了 ou 因为 ? Complétez puis traduisez]
Complétez chaque phrase avec \zh{为了} ou \zh{因为}, puis traduisez en français.

\begin{enumerate}
\item \zh{……身体好，他每天早上跑步。}
\item \zh{……下雨，所以足球比赛取消了。}
\item \zh{……通过汉语考试，玛丽每天复习生词。}
\item \zh{他……工作太忙，没有时间运动。}
\item \zh{……让家乡更美丽，我们种了很多树。}
\end{enumerate}
\end{exercice}

\begin{corrige}[Exercice : 为了 ou 因为]
\begin{enumerate}
\item \zh{为了身体好，他每天早上跑步。} — Pour être en bonne santé, il court chaque matin. (but futur → \zh{为了})
\item \zh{因为下雨，所以足球比赛取消了。} — Parce qu'il pleut, le match de football est annulé. (cause réelle + \zh{所以} → \zh{因为})
\item \zh{为了通过汉语考试，玛丽每天复习生词。} — Pour réussir l'examen de chinois, Marie révise les mots nouveaux chaque jour. (but → \zh{为了})
\item \zh{他因为工作太忙，没有时间运动。} — Parce qu'il est trop occupé par le travail, il n'a pas le temps de faire du sport. (cause présente → \zh{因为})
\item \zh{为了让家乡更美丽，我们种了很多树。} — Pour rendre notre village natal plus beau, nous avons planté beaucoup d'arbres. (but → \zh{为了} ; notez \zh{让} « rendre, faire que » après \zh{为了})
\end{enumerate}
\end{corrige}

\begin{exercice}[Traduction français → chinois]
Traduisez en chinois avec \zh{为了} :
\begin{enumerate}
\item Pour apprendre le chinois, il est allé en Chine.
\item Pour la santé de sa famille, elle cuisine des légumes chaque jour.
\item Pour trouver un bon emploi, les jeunes de Douala apprennent l'anglais et le chinois.
\end{enumerate}
\end{exercice}

\begin{corrige}[Exercice : traduction]
\begin{enumerate}
\item \zh{为了学习汉语，他去了中国。} \emph{Wèile xuéxí Hànyǔ, tā qù le Zhōngguó.}
\item \zh{为了家人的健康，她每天做蔬菜。} \emph{Wèile jiārén de jiànkāng, tā měitiān zuò shūcài.}
\item \zh{为了找到好工作，杜阿拉的年轻人学习英语和汉语。} \emph{Wèile zhǎodào hǎo gōngzuò, Dù'ālā de niánqīngrén xuéxí Yīngyǔ hé Hànyǔ.}
\end{enumerate}
\end{corrige}

\begin{attention}[Les trois pièges des francophones avec 为了]
\begin{enumerate}
\item \textbf{Confondre but et cause} : avant d'écrire, demandez-vous si l'on répond à « pourquoi ? » (\zh{因为}) ou « dans quel but ? » (\zh{为了}). \zh{因为生病} (parce qu'il est malade — c'est un fait) ≠ \zh{为了不生病} (pour ne pas tomber malade — c'est un objectif).
\item \textbf{Placer le but après l'action} : calque du français « il court pour sa santé ». En chinois : \zh{为了健康，他跑步。} et non l'inverse.
\item \textbf{Oublier la virgule} en tête de phrase : \zh{为了健康，他每天运动。} — la virgule \zh{，} marque la pause entre le but et l'action ; elle est obligatoire quand \zh{为了} ouvre la phrase.
\end{enumerate}
\end{attention}

\begin{aretenir}[L'essentiel sur 为了]
\begin{itemize}
\item \zh{为了} + but (nom ou proposition) exprime l'\textbf{objectif visé} : « pour, afin de ».
\item Schéma : \zh{为了} + but，+ sujet + action ; ou sujet + \zh{为了} + but + action.
\item Le but précède \textbf{toujours} l'action en chinois, contrairement au français.
\item \zh{为了} (but, futur) s'oppose à \zh{因为} (cause, présent/passé), souvent suivi de \zh{所以}.
\item Test infaillible : « dans quel but ? » → \zh{为了} ; « pourquoi ? » → \zh{因为}.
\end{itemize}
\end{aretenir}

\coursec{Additionner les idées : 不但…而且… et synthèse}

Dans la section précédente, nous avons appris à exprimer le \emph{but} avec \zh{为了} (wèile) : \zh{为了健康，我每天运动。} « Pour ma santé, je fais du sport chaque jour. » Mais souvent, une même action a \emph{plusieurs} effets positifs, et nous voulons les énumérer en allant \emph{du plus simple au plus fort}. Le français dit « non seulement… mais aussi… » ; le chinois utilise la paire de corrélatifs \zh{不但…而且…} (bùdàn… érqiě…). C'est la dernière structure de notre leçon, et elle va nous permettre de combiner toutes les autres dans des phrases riches.

\courssub{Observation : découvrir la structure}

\begin{textebox}[Le sport : un double bénéfice]
运动不但对身体好，而且对学习也有帮助。\\
\emph{Yùndòng bùdàn duì shēntǐ hǎo, érqiě duì xuéxí yě yǒu bāngzhù.}\\
« Le sport est non seulement bon pour la santé, mais aussi utile pour les études. »\\[4pt]
他不但会说汉语，而且会说英语。\\
\emph{Tā bùdàn huì shuō Hànyǔ, érqiě huì shuō Yīngyǔ.}\\
« Non seulement il sait parler chinois, mais aussi l'anglais. »\\[4pt]
不但老师来了，而且校长也来了。\\
\emph{Bùdàn lǎoshī lái le, érqiě xiàozhǎng yě lái le.}\\
« Non seulement le professeur est venu, mais le directeur est venu aussi. »
\end{textebox}

Observons attentivement ces trois phrases :

\begin{itemize}
\item La paire \zh{不但…而且…} encadre deux informations : la deuxième est toujours \emph{plus forte}, plus surprenante ou plus importante que la première (gradation ascendante).
\item Dans les phrases 1 et 2, le sujet (\zh{运动}, \zh{他}) est \emph{le même} pour les deux propositions : il est placé \textbf{une seule fois, avant} \zh{不但}.
\item Dans la phrase 3, les sujets sont \emph{différents} (\zh{老师} / \zh{校长}) : chaque sujet se place \emph{après} son corrélatif, et \zh{也} apparaît dans la deuxième proposition.
\end{itemize}

\begin{propriete}[Règle : 不但…而且…]
\zh{不但…而且…} (bùdàn… érqiě…) relie deux informations en \textbf{gradation ascendante} : l'élément introduit par \zh{而且} est toujours « plus fort », plus remarquable ou inattendu.

\textbf{Cas 1 — même sujet :} Sujet + \zh{不但} + élément 1，\zh{而且} + élément 2.\\
\zh{她不但聪明，而且很努力。} \emph{Tā bùdàn cōngming, érqiě hěn nǔlì.} « Elle est non seulement intelligente, mais aussi très travailleuse. »

\textbf{Cas 2 — sujets différents :} \zh{不但} + sujet 1 + …，\zh{而且} + sujet 2 + \zh{也} + …\\
\zh{不但学生喜欢这部电影，而且老师也喜欢。} \emph{Bùdàn xuésheng xǐhuan zhè bù diànyǐng, érqiě lǎoshī yě xǐhuan.} « Non seulement les élèves aiment ce film, mais les professeurs l'aiment aussi. »

\textbf{Variantes :} \zh{不仅…而且…} (bùjǐn… érqiě…) est plus soutenu, fréquent à l'écrit. \zh{而且} peut être renforcé par \zh{还} ou \zh{也} : \zh{他不但会说汉语，而且还会写汉字。} « Non seulement il parle chinois, mais il sait même écrire les caractères. »
\end{propriete}

\begin{popfigure}
\begin{tikzpicture}[font=\small, node distance=0.35cm]
\node[draw=popblue, fill=popblueL, rounded corners, minimum width=2.2cm, minimum height=0.9cm, align=center] (s) {Sujet\\他};
\node[draw=popgreen, fill=popgreenL, rounded corners, minimum width=3.2cm, minimum height=0.9cm, align=center, right=0.6cm of s] (b1) {不但 + idée 1\\会说汉语};
\node[draw=poporange, fill=poporangeL, rounded corners, minimum width=3.6cm, minimum height=0.9cm, align=center, right=0.6cm of b1, yshift=0.9cm] (b2) {而且 + idée 2 \textbf{plus forte}\\还会写汉字};
\draw[-{Stealth}, thick, popblue] (s) -- (b1);
\draw[-{Stealth}, very thick, poporange] (b1.north east) -- (b2.west) node[midway, above, sloped, popdark]{gradation $\uparrow$};
\end{tikzpicture}
\legende{Schéma de la gradation ascendante : le sujet commun précède 不但, l'idée après 而且 est plus forte.}
\end{popfigure}

\courssub{Comparaison avec le français}

La correspondance avec « non seulement… mais aussi… » est presque directe, ce qui est une bonne nouvelle pour les francophones. Deux différences importantes :

\begin{center}
\begin{tabularx}{\linewidth}{|X|X|X|}\hline
\rowcolor{popblueL} \textbf{Point} & \textbf{Français} & \textbf{Chinois} \\ \hline
Sujet identique & souvent répété ou repris par un pronom (« il… mais il aussi… ») & placé \textbf{une seule fois avant} \zh{不但} \\ \hline
Sujets différents & « non seulement le professeur… mais le directeur aussi » & \zh{不但} + sujet 1…，\zh{而且} + sujet 2 + \zh{也}… \\ \hline
Ordre & gradation ascendante & gradation ascendante (identique) \\ \hline
\end{tabularx}
\end{center}

\begin{attention}[Piège des francophones : ne répétez pas le sujet !]
Sous l'influence du français, les élèves écrivent souvent : *\zh{他不但会说汉语，而且他还会说英语。} Cette phrase n'est pas fautive à 100 \%, mais elle est \emph{lourde} et typique d'un apprenant étranger. Quand le sujet est identique, le chinois le place \textbf{avant} \zh{不但} et ne le répète pas : \zh{他不但会说汉语，而且还会说英语。} Autre erreur : oublier \zh{而且} et n'écrire que \zh{不但} — la paire fonctionne toujours \emph{ensemble}, comme « non seulement… mais… ».
\end{attention}

\begin{exemplebox}[Exemples variés (contexte Cameroun)]
\begin{itemize}
\item \zh{杜阿拉不但人口多，而且经济也很发达。} \emph{Dùālā bùdàn rénkǒu duō, érqiě jīngjì yě hěn fādá.} « Douala est non seulement très peuplée, mais son économie est aussi très développée. »
\item \zh{汉语不但有意思，而且很有用。} \emph{Hànyǔ bùdàn yǒu yìsi, érqiě hěn yǒuyòng.} « Le chinois est non seulement intéressant, mais aussi très utile. »
\item \zh{不但我喜欢恩多莱，而且我的中国朋友也喜欢。} \emph{Bùdàn wǒ xǐhuan ēn duō lái, érqiě wǒ de Zhōngguó péngyou yě xǐhuan.} « Non seulement j'aime le ndolé, mais mon ami chinois l'aime aussi. »
\item \zh{为了健康，我不但吃得比以前好，而且每天运动得更多了。} \emph{Wèile jiànkāng, wǒ bùdàn chī de bǐ yǐqián hǎo, érqiě měitiān yùndòng de gèng duō le.} « Pour ma santé, non seulement je mange mieux qu'avant, mais en plus je fais plus de sport chaque jour. » — remarquez la combinaison de \textbf{trois structures de la leçon} : \zh{为了} + \zh{比} + \zh{更} à l'intérieur de \zh{不但…而且…} !
\end{itemize}
\end{exemplebox}

\begin{exoresolu}[Transformation : combiner deux phrases]
Énoncé : combinez avec \zh{不但…而且…} : \zh{他会说汉语。他会写汉字。} (même sujet)
\tcblower
Solution : les deux propositions ont le même sujet \zh{他}, donc le sujet se place une seule fois avant \zh{不但}, et l'élément le plus fort (écrire les caractères, plus difficile que parler) va après \zh{而且}, renforcé par \zh{还} :\\
\zh{他不但会说汉语，而且还会写汉字。}\\
\emph{Tā bùdàn huì shuō Hànyǔ, érqiě hái huì xiě Hànzì.}\\
« Non seulement il sait parler chinois, mais il sait même écrire les caractères. »\\
Méthode : 1) vérifier si les sujets sont identiques ; 2) placer l'idée la plus forte en second ; 3) ne pas répéter le sujet.
\end{exoresolu}

\begin{savaistu}[Pour aller plus loin : 不仅…甚至…]
Il existe une gradation encore plus forte : \zh{那些…甚至…} (bùjǐn… shènzhì…), « non seulement… mais même… ». Exemple : \zh{他那些会开车，那些会开飞机。} \emph{Tā bùjǐn huì kāichē, shènzhì huì kāi fēijī.} « Non seulement il sait conduire une voiture, mais il sait même piloter un avion. » Cette structure est appréciée à l'examen dans les rédactions : elle montre une maîtrise fine de la gradation.
\end{savaistu}

\begin{definition}[Vocabulaire clé de la section]
\begin{center}
\begin{tabularx}{\linewidth}{|X|X|X|X|}\hline
\rowcolor{popblueL} \textbf{Caractères} & \textbf{Pinyin} & \textbf{Nature} & \textbf{Traduction} \\ \hline
不但 & bùdàn & conj. & non seulement \\ \hline
而且 & érqiě & conj. & mais aussi / de plus \\ \hline
那些 & bùjǐn & conj. & non seulement (soutenu) \\ \hline
那些 & shènzhì & adv. & même / jusqu'à / allant jusqu'à \\ \hline
为了 & wèile & prép. & pour / afin de / dans le but de \\ \hline
比 & bǐ & prép. & (marque de comparaison) que \\ \hline
不如 & bùrú & v./prép. & ne pas valoir / être inférieur à \\ \hline
更 & gèng & adv. & encore plus \\ \hline
还 & hái & adv. & encore / aussi / même \\ \hline
一点 & yīdiǎn & quant. & un peu \\ \hline
多了 & duō le & structure & beaucoup plus (après adj.) \\ \hline
空气 & kōngqì & n. & air \\ \hline
风景 & fēngjǐng & n. & paysage / décor \\ \hline
恩多莱 & ēn duō lái & n. & ndolé (plat camerounais) \\ \hline
经济 & jīngjì & n. & économie \\ \hline
发达 & fādá & adj. & développé \\ \hline
有用 & yǒuyòng & adj. & utile \\ \hline
有意思 & yǒu yìsi & adj. & intéressant \\ \hline
花钱 & huā qián & v.o. & coûter de l'argent / dépenser \\ \hline
\end{tabularx}
\end{center}
\end{definition}

\courssub{Synthèse générale des cinq structures de la leçon}

\begin{popfigure}
\begin{tikzpicture}[font=\footnotesize]
\node[draw=popblue, fill=popblueL, rounded corners, align=center, minimum width=4.6cm] (a) at (0,3.2) {\zh{比} : A 比 B + adj.\\今天比昨天热。};
\node[draw=popgreen, fill=popgreenL, rounded corners, align=center, minimum width=4.6cm] (b) at (6.4,3.2) {\zh{比} + 一点/多了\\今天比昨天热多了。};
\node[draw=poporange, fill=poporangeL, rounded corners, align=center, minimum width=4.6cm] (c) at (0,1.6) {\zh{比} + 更/还 + adj.\\今天比昨天更热。};
\node[draw=poppurple, fill=poppurpleL, rounded corners, align=center, minimum width=4.6cm] (d) at (6.4,1.6) {\zh{不如} : A 不如 B (+ adj.)\\我不如他高。};
\node[draw=popgold, fill=popgoldL, rounded corners, align=center, minimum width=4.6cm] (e) at (0,0) {\zh{为了} + but\\为了健康，我每天运动。};
\node[draw=poppink, fill=poppinkL, rounded corners, align=center, minimum width=4.6cm] (f) at (6.4,0) {\zh{不但…而且…}\\他不但聪明，而且努力。};
\draw[-{Stealth}, thick, popmuted] (a) -- (b);
\draw[-{Stealth}, thick, popmuted] (a) -- (c);
\draw[-{Stealth}, thick, popmuted] (c) -- (d);
\end{tikzpicture}
\legende{Carte récapitulative des six schémas étudiés dans la leçon 10.}
\end{popfigure}

\begin{center}
\begin{tabularx}{\linewidth}{|X|X|X|X|}\hline
\rowcolor{popblueL} \textbf{Structure} & \textbf{Emploi} & \textbf{Exemple} & \textbf{Piège} \\ \hline
A \zh{比} B + adj. & comparer deux éléments & \zh{他比我高。} Il est plus grand que moi. & pas de \zh{很} après \zh{比} (*\zh{比我很高}) \\ \hline
A \zh{比} B + adj. + \zh{一点/多了} & préciser l'écart (petit / grand) & \zh{他比我高一点。} Il est un peu plus grand que moi. & \zh{一点} se met \emph{après} l'adjectif \\ \hline
A \zh{比} B \zh{更/还} + adj. & renforcer la comparaison & \zh{他比我更高。} Il est encore plus grand que moi. & \zh{更/还} se place \emph{avant} l'adjectif \\ \hline
A \zh{不如} B (+ adj.) & « A n'est pas aussi… que B » & \zh{我不如他高。} Je ne suis pas aussi grand que lui. & ne pas ajouter \zh{不} devant l'adjectif \\ \hline
\zh{为了} + but + , + action & exprimer le but & \zh{为了健康，我多运动。} Pour ma santé, je fais plus de sport. & \zh{为了} en tête de phrase, pas en fin \\ \hline
\zh{不但…而且…} & additionner en gradation & \zh{他不但聪明，而且努力。} & ne pas répéter le sujet identique \\ \hline
\end{tabularx}
\end{center}

\begin{exercice}[Exercices de transformation finaux]
\begin{enumerate}
\item Combinez : \zh{雅温得的空气好。雅温得的风景美。} (même sujet, avec \zh{不但…而且…})
\item Traduisez : « Pour apprendre le chinois, Marie étudie non seulement le soir, mais elle écoute aussi des chansons chinoises. »
\item Phrase complexe : dites que le sport est meilleur pour la santé que la télévision (\zh{比}), et qu'il est non seulement gratuit mais aussi amusant (\zh{不但…而且…}).
\item Corrigez : *\zh{他不但会说汉语，而且他还会说法语。} (sujet répété)
\end{enumerate}
\end{exercice}

\begin{corrige}[Exercice final]
1. \zh{雅温得不但空气好，而且风景也美。} \emph{Yǎwēndé bùdàn kōngqì hǎo, érqiě fēngjǐng yě měi.} — sujet unique placé avant \zh{不但}, \zh{也} renforce la seconde idée.\\
2. \zh{为了学汉语，Marie 不但晚上学习，而且还听中文歌。} \emph{Wèile xué Hànyǔ, Marie bùdàn wǎnshang xuéxí, érqiě hái tīng Zhōngwén gē.}\\
3. \zh{运动比看电视对健康更好；运动不但不花钱，而且很有意思。} \emph{Yùndòng bǐ kàn diànshì duì jiànkāng gèng hǎo ; yùndòng bùdàn bù huā qián, érqiě hěn yǒu yìsi.}\\
4. \zh{他不但会说汉语，而且还会说法语。} — on supprime le deuxième \zh{他}, car le sujet est identique.
\end{corrige}

\begin{aretenir}[L'essentiel de la section]
\zh{不但…那些…} = « non seulement… mais aussi… », en gradation ascendante. Sujet identique : une seule fois, avant \zh{不但}. Sujets différents : chaque sujet après son corrélatif, avec \zh{也}. Variante soutenue : \zh{那些…那些…}. Gradation forte : \zh{那些…那些…}. Cette structure se combine librement avec \zh{为了}, \zh{比} et \zh{更} pour produire des phrases d'examen riches et naturelles.
\end{aretenir}

\newpage

\coursec{Activité d'intégration}

\coursec{Leçon 10 — Grammaire : Phrases comparatives avec 不如 ; A比B 形容词+一点/多了 ; A比B 更/还+形容词 ; 为了 ; 不但…而且…}

\courssub{Observation et rappel : la comparaison avec 比}

\begin{propriete}[Rappel : la structure de base A 比 B + adjectif]
\textbf{Schéma :} \cle{Sujet A + 比 + Sujet B + Adjectif (sans 很/非常)} \\
\textbf{Emploi :} Exprimer une supériorité de A sur B pour une qualité donnée. En chinois, l'adjectif n'est \textbf{jamais} précédé de \zh{很} (hěn), \zh{非常} (fēicháng) ou \zh{特别} (tèbié) dans cette structure ; le mot \zh{比} marque déjà le degré comparatif.
\textbf{Exemples :}
\begin{itemize}
\item \zh{雅温得比杜阿拉凉快。} \emph{Yǎwēndé bǐ Dù'ālā liángkuai.} « Yaoundé est plus frais que Douala. »
\item \zh{这本书比那本贵。} \emph{Zhè běn shū bǐ nà běn guì.} « Ce livre est plus cher que celui-là. »
\item \zh{哥哥比弟弟高。} \emph{Gēge bǐ dìdi gāo.} « Le grand frère est plus grand que le petit frère. »
\end{itemize}
\textbf{Attention francophones :} Ne dites pas \zh{*雅温得比杜阿拉很凉快}. Le \zh{很} est incompatible avec \zh{比}. Pour accentuer, on utilise les structures vues ci-dessous (\zh{一点}, \zh{多了}, \zh{更}, \zh{还}).
\end{propriete}

\courssub{Nuancer l'écart : A比B + adjectif + 一点 / 多了}

\begin{propriete}[Degrés de l'écart : \zh{一点 / 一些} (petit) vs \zh{多了 / 得多} (grand)]
\begin{center}
\begin{tabularx}{\linewidth}{|X|X|X|}
\hline
\rowcolor{popblueL} \textbf{Structure} & \textbf{Exemple (Caractères + Pinyin)} & \textbf{Traduction} \\
\hline
A 比 B + Adj + \zh{一点 / 一些} & \zh{雅温得比杜阿拉凉快一点。} \emph{Yǎwēndé bǐ Dù'ālā liángkuai yìdiǎn.} & Yaoundé est un peu plus frais que Douala. \\
\hline
A 比 B + Adj + \zh{多了 / 得多} & \zh{杜阿拉的车比雅温得多多了。} \emph{Dù'ālā de chē bǐ Yǎwēndé duō duō le.} & Il y a beaucoup plus de voitures à Douala qu'à Yaoundé. \\
\hline
\end{tabularx}
\end{center}
\textbf{Règles clés :}
\begin{itemize}
\item \zh{一点} (yìdiǎn) ou \zh{一些} (yìxiē) : écart faible, nuance subtile. \zh{一点} s'emploie surtout avec des adjectifs monosyllabiques ou disyllabiques d'état ; \zh{一些} souvent avec des verbes ou adjectifs de quantité.
\item \zh{多了} (duō le) ou \zh{得多} (de duō) : écart important, évident. \zh{多了} est plus oral ; \zh{得多} légèrement plus formel/écrit.
\item \textbf{Ordre strict :} Le complément de degré (\zh{一点}, \zh{多了}) se place \textbf{après} l'adjectif, jamais avant.
\end{itemize}
\end{propriete}

\begin{attention}[Erreurs fréquentes]
\begin{itemize}
\item \zh{*比杜阿拉一点凉快} (ordre incorrect : le degré ne précède pas l'adjectif).
\item \zh{*雅温得比杜阿拉凉快很} (pas de \zh{很} en finale).
\item Confondre \zh{多了} (adverbe de degré) et \zh{很多} (hěn duō, « beaucoup », qui ne s'utilise pas comme complément de comparaison).
\end{itemize}
\end{attention}


\courssub{Renforcer la comparaison : A比B 更 / 还 + adjectif}

\begin{propriete}[Renforcement avec \zh{更} (gèng) et \zh{还} (hái)]
\textbf{Schéma :} \cle{Sujet A + 比 + Sujet B + 更 / 还 + Adjectif} \\
\textbf{Nuances :}
\begin{itemize}
\item \zh{更} (gèng) : « encore plus », « davantage ». Neutre, factuel. Utilisé pour affirmer une supériorité claire.
\item \zh{还} (hái) : « encore plus », « même plus ». Souvent subjectif, insiste sur un dépassement des attentes ou une surprise. Fréquent à l'oral.
\end{itemize}
\begin{center}
\begin{tabularx}{\linewidth}{|X|X|X|}
\hline
\rowcolor{popblueL} \textbf{Structure} & \textbf{Exemple} & \textbf{Traduction} \\
\hline
A 比 B \zh{更} + Adj & \zh{雅温得的饭比杜阿拉更便宜。} \emph{Yǎwēndé de fàn bǐ Dù'ālā gèng piányi.} & Les repas de Yaoundé sont encore moins chers que ceux de Douala. \\
\hline
A 比 B \zh{还} + Adj & \zh{这个问题比那个还难。} \emph{Zhège wèntí bǐ nàge hái nán.} & Ce problème est encore plus difficile que celui-là. \\
\hline
\end{tabularx}
\end{center}
\end{propriete}

\begin{attention}[Piège classique]
Ne placez \textbf{jamais} \zh{更} ou \zh{still} avant \zh{比}.
\emph{Incorrect :} \zh{*雅温得更比杜阿拉便宜。} \\
\emph{Correct :} \zh{雅温得比杜阿拉更便宜。}
\end{attention}


\courssub{La comparaison négative : A 不如 B (+ adjectif)}

\begin{propriete}[Infériorité : \zh{不如} (bùrú)]
\textbf{Schéma :} \cle{Sujet A + 不如 + Sujet B (+ Adjectif)} \\
\textbf{Emploi :} « A n'est pas aussi… que B » / « A est inférieur à B ». Structure élégante, souvent préférée à \zh{比} + négation (\zh{没有}) pour sa concision et sa politesse (voir \emph{Le saviez-vous ?}).
\textbf{Exemples :}
\begin{itemize}
\item \zh{杜阿拉不如雅温得凉快。} \emph{Dù'ālā bùrú Yǎwēndé liángkuai.} « Douala n'est pas aussi frais que Yaoundé. »
\item \zh{我的中文不如你的好。} \emph{Wǒ de Zhōngwén bùrú nǐ de hǎo.} « Mon chinois n'est pas aussi bon que le tien. »
\item \zh{这家饭馆不如那家。} \emph{Zhè jiā fànguǎn bùrú nà jiā.} « Ce restaurant n'est pas aussi bon que celui-là. » (Adjectif omis par contexte).
\end{itemize}
\end{propriete}

\begin{attention}[Distinction \zh{不如} vs \zh{没有}]
\begin{itemize}
\item \zh{A 不如 B (+ Adj)} : Comparaison de qualité/degré. Standard, neutre à formel.
\item \zh{A 没有 B (+ 这么/那么 + Adj)} : « A n'est pas aussi… que B ». Plus oral, insiste sur la non-égalité.
\item \textbf{Interdiction :} Ne jamais mettre \zh{比} dans une phrase avec \zh{不如}. \emph{Incorrect :} \zh{*杜阿拉比雅温得不如凉快。}
\end{itemize}
\end{attention}


\courssub{Exprimer le but avec 为了}

\begin{propriete}[Finalité : \zh{为了} (wèile) « pour, afin de, dans le but de »]
\textbf{Schéma :} \cle{为了 + But (Nom / Verbe / Proposition), Sujet + Verbe/Action principale} \\
\textbf{Position :} En chinois, la proposition de but (\zh{为了…}) se place \textbf{systématiquement en tête de phrase} (thème), avant le sujet. En français, « pour + infinitif » se place souvent à la fin (rhème).
\textbf{Exemples :}
\begin{itemize}
\item \zh{为了学习中文，我每天练习。} \emph{Wèile xuéxí Zhōngwén, wǒ měi tiān liànxí.} « Pour apprendre le chinois, je m'entraîne tous les jours. »
\item \zh{为了健康，我们不抽烟。} \emph{Wèile jiànkāng, wǒmen bù chōuyān.} « Pour la santé, nous ne fumons pas. »
\item \zh{为了不迟到，他早早就出发了。} \emph{Wèile bù chídào, tā zǎozǎo jiù chūfā le.} « Pour ne pas être en retard, il est parti très tôt. »
\end{itemize}
\end{propriete}

\begin{attention}[Erreurs de calque français]
\begin{itemize}
\item \emph{Incorrect :} \zh{*我每天练习，为了学习中文。} (Place finale à la française).
\item \emph{Correct :} \zh{为了学习中文，我每天练习。}
\item \zh{为了} est suivi d'un nom, d'un verbe ou d'une phrase. Pas de \zh{de} (的) après \zh{为了}.
\end{itemize}
\end{attention}


\courssub{Additionner les idées : 不但…而且…}

\begin{propriete}[Coordination additive : \zh{不但…而且…} (búdàn… érqiě…) « non seulement… mais aussi… »]
\textbf{Schéma :}
\begin{itemize}
\item \textbf{Sujet commun :} \cle{Sujet + 不但 + Verbe/Adj 1, 而且 + Verbe/Adj 2}. (Le sujet \textbf{ne se répète pas}).
\item \textbf{Sujets différents :} \cle{不但 + Sujet 1 + Verbe/Adj 1, 而且 + Sujet 2 + Verbe/Adj 2}.
\end{itemize}
\textbf{Exemples :}
\begin{itemize}
\item \zh{雅温得不但环境好，而且空气也很新鲜。} \emph{Yǎwēndé búdàn huánjìng hǎo, érqiě kōngqì yě hěn xīnxiān.} (Sujet commun \zh{雅温得}).
\item \zh{不但学生努力，而且老师也很负责。} \emph{Búdàn xuésheng nǔlì, érqiě lǎoshī yě hěn fùzé.} (Sujets différents).
\item \zh{为了考试，我不但复习了课文，而且做了很多练习。} \emph{Wèile kǎoshì, wǒ búdàn fùxí le kèwén, érqiě zuò le hěn duō liànxí.} (Combinaison \zh{为了} + \zh{不但…而且…}).
\end{itemize}
\end{propriete}

\begin{attention}[Points de vigilance]
\begin{itemize}
\item \zh{而且} (érqiě) s'écrit avec \zh{而} (ér) + \zh{且} (qiě). Ne pas confondre avec \zh{并且} (bìngqiě) « et en plus », qui ne s'emploie pas dans la corrélation \zh{不但}.
\item Ponctuation : virgule chinoise \zh{，} obligatoire entre les deux propositions.
\item \zh{也} (yě) ou \zh{还} (hái) s'ajoutent souvent dans la seconde proposition pour renforcer l'addition (\zh{而且…也…} / \zh{而且…还…}).
\end{itemize}
\end{attention}


\begin{popfigure}
\begin{tikzpicture}[
  mindmap,
  grow cyclic,
  every node/.style=concept,
  concept color=popblue,
  level 1/.append style={level distance=4.5cm, sibling angle=72},
  level 2/.append style={level distance=3cm, sibling angle=40},
  text=white,
  font=\small
]
\node[align=center]{Structures\\de la Leçon 10}
  child[concept color=popgreen] { node[align=center]{Nuancer l'écart\\ \zh{比...一点/多了}} }
  child[concept color=poporange] { node[align=center]{Renforcer\\ \zh{比...更/还}} }
  child[concept color=poppurple] { node[align=center]{Infériorité\\ \zh{不如}} }
  child[concept color=poppink] { node[align=center]{But\\ \zh{为了}} }
  child[concept color=popgold] { node[align=center]{Addition\\ \zh{不但...而且}} };
\end{tikzpicture}
\legende{Récapitulatif visuel des cinq structures grammaticales de la leçon}
\end{popfigure}

\courssub{Activité d'intégration : Débat « Ma ville contre ta ville »}

\courssub{Objectifs de l'activité}

À l'issue de cette activité d'intégration, tu seras capable de :

\begin{itemize}
\item mobiliser dans une même production les cinq structures de la leçon : \zh{不如}, \zh{A比B} + adjectif + \zh{一点/多了}, \zh{A比B更/还} + adjectif, \zh{为了} et \zh{不但…而且…} ;
\item comparer deux villes camerounaises sur des critères précis (climat, transport, nourriture, sport, études) ;
\item argumenter oralement et par écrit pour défendre un point de vue, en respectant l'ordre des mots chinois ;
\item réagir aux arguments d'un camarade en reformulant et en renforçant une comparaison ;
\item rédiger un court texte argumentatif de huit à dix phrases, sans faute de structure.
\end{itemize}

\courssub{Situation}

Le club de chinois de ton lycée organise un grand débat amical intitulé « Ma ville contre ta ville ». Deux équipes s'affrontent : l'équipe de Yaoundé et l'équipe de Douala. Chaque équipe doit convaincre le jury (le professeur et les autres élèves) que sa ville est la meilleure pour vivre, étudier et se distraire. Pour préparer le débat, le professeur distribue une fiche d'information rédigée par deux élèves chinois en échange scolaire au Cameroun.

\begin{textebox}[Fiche d'information : Yaoundé et Douala]
雅温得是喀麦隆的首都，天气不太热，比杜阿拉凉快一点。\\
Yǎwēndé shì Kāmàilóng de shǒudū, tiānqì bú tài rè, bǐ Dù'ālā liángkuai yìdiǎn.\\
« Yaoundé est la capitale du Cameroun ; le temps n'y est pas trop chaud, il y fait un peu plus frais qu'à Douala. »\\[4pt]
杜阿拉很大，人很多，车也很多，所以路上常常堵车。\\
Dù'ālā hěn dà, rén hěn duō, chē yě hěn duō, suǒyǐ lùshang chángcháng dǔchē.\\
« Douala est très grande, il y a beaucoup de gens et beaucoup de voitures, donc il y a souvent des embouteillages. »\\[4pt]
为了学习中文，雅温得的学生不但去孔子学院上课，而且常常练习说汉语。\\
Wèile xuéxí Zhōngwén, Yǎwēndé de xuésheng búdàn qù Kǒngzǐ Xuéyuàn shàngkè, érqiě chángcháng liànxí shuō Hànyǔ.\\
« Pour apprendre le chinois, les élèves de Yaoundé vont non seulement en cours à l'Institut Confucius, mais s'entraînent aussi souvent à parler chinois. »\\[4pt]
杜阿拉的海鲜很好吃，可是雅温得的饭更便宜。\\
Dù'ālā de hǎixiān hěn hǎochī, kěshì Yǎwēndé de fàn gèng piányi.\\
« Les fruits de mer de Douala sont très bons, mais les repas de Yaoundé sont encore moins chers. »
\end{textebox}

\courssub{Consignes}

\textbf{Étape 1 — Comprendre la fiche (travail individuel, 10 minutes).} Lis le document ci-dessus. Souligne les cinq structures de la leçon et recopie-les dans ton cahier avec leur schéma. Réponds ensuite en français : quel avantage chaque ville possède-t-elle selon la fiche ?

\textbf{Étape 2 — Enrichir le lexique (par deux, 10 minutes).} Avec ton voisin, complète le tableau de vocabulaire utile ci-dessous : ajoute trois adjectifs de ton choix pour décrire ta ville (par exemple \zh{干净} gānjìng « propre », \zh{安全} ānquán « sûr », \zh{有名} yǒumíng « célèbre ») et vérifie le pinyin dans le dictionnaire.

\textbf{Étape 3 — Préparer les arguments (en équipe, 15 minutes).} Chaque équipe rédige au moins cinq arguments, un par structure imposée :
\begin{enumerate}
\item une phrase avec \zh{不如} pour minimiser un point faible de la ville adverse ;
\item une phrase avec \zh{比…一点} pour un petit écart ;
\item une phrase avec \zh{比…多了} pour un grand écart ;
\item une phrase avec \zh{更} ou \zh{还} pour renforcer un avantage ;
\item une phrase avec \zh{为了} pour exprimer un but, et une phrase avec \zh{不但…而且…} pour additionner deux qualités.
\end{enumerate}

\textbf{Étape 4 — Le débat (oral, 15 minutes).} Les équipes présentent leurs arguments à tour de rôle. Règle du jeu : pour répondre à un argument adverse, tu dois commencer par reformuler la phrase de ton adversaire avec \zh{不如}, puis contre-attaquer avec \zh{更} ou \zh{还}. Le jury note chaque structure correctement employée.

\textbf{Étape 5 — Production écrite finale (individuel, 20 minutes).} Rédige un court texte de huit à dix phrases intitulé \zh{我的城市最好} (« Ma ville est la meilleure »), en utilisant au moins une fois chacune des cinq structures. Tu peux défendre ta vraie ville (Bafoussam, Buea, Garoua, Bamenda…) contre une autre ville camerounaise.

\courssub{Ressources à mobiliser}

\begin{aretenir}[Rappel synthétique des cinq structures]
\begin{itemize}
\item \cle{A 不如 B (+ adjectif)} : A est inférieur à B. \zh{杜阿拉不如雅温得凉快。} \emph{Dù'ālā bùrú Yǎwēndé liángkuai.} « Douala est moins fraîche que Yaoundé. »
\item \cle{A 比 B + adjectif + 一点/一些} : petit écart. \zh{雅温得比杜阿拉凉快一点。} \emph{Yǎwēndé bǐ Dù'ālā liángkuai yìdiǎn.} « Yaoundé est un peu plus fraîche que Douala. »
\item \cle{A 比 B + adjectif + 多了/得多} : grand écart. \zh{杜阿拉的车比雅温得多多了。} \emph{Dù'ālā de chē bǐ Yǎwēndé duō duō le.} « Il y a beaucoup plus de voitures à Douala qu'à Yaoundé. »
\item \cle{A 比 B 更/还 + adjectif} : renforcement. \zh{雅温得的饭更便宜。} \emph{Yǎwēndé de fàn gèng piányi.} « Les repas de Yaoundé sont encore moins chers. »
\item \cle{为了 + but, …} : « pour, afin de ». \zh{为了健康，我们多运动。} \emph{Wèile jiànkāng, wǒmen duō yùndòng.} « Pour la santé, nous faisons plus de sport. »
\item \cle{不但…而且…} : « non seulement… mais aussi… ». \zh{这个城市不但大，而且很干净。} \emph{Zhège chéngshì búdàn dà, érqiě hěn gānjìng.} « Cette ville est non seulement grande, mais aussi très propre. »
\end{itemize}
\end{aretenir}

\begin{center}
\begin{tabularx}{\linewidth}{|X|X|X|X|}
\hline
\rowcolor{popblueL} Caractères & Pinyin & Nature & Traduction \\
\hline
城市 & chéngshì & nom & ville \\
\hline
首都 & shǒudū & nom & capitale \\
\hline
凉快 & liángkuai & adjectif & frais (climat) \\
\hline
堵车 & dǔchē & verbe & être embouteillé \\
\hline
海鲜 & hǎixiān & nom & fruits de mer \\
\hline
便宜 & piányi & adjectif & bon marché \\
\hline
环境 & huánjìng & nom & environnement \\
\hline
空气 & kōngqì & nom & air \\
\hline
方便 & fāngbiàn & adjectif & pratique \\
\hline
生活 & shēnghuó & nom / verbe & vie ; vivre \\
\hline
适合 & shìhé & verbe & convenir, être adapté \\
\hline
新鲜 & xīnxiān & adjectif & frais (air, nourriture), pur \\
\hline
负责 & fùzé & adjectif & responsable \\
\hline
\end{tabularx}
\end{center}

\begin{attention}[Pièges à éviter pendant le débat]
\begin{itemize}
\item Ne dis jamais \zh{比…很热} : après \zh{比}, l'adjectif ne prend ni \zh{很} ni \zh{非常}. Dis \zh{雅温得比杜阿拉凉快一点}.
\item Avec \zh{更/还}, l'ordre est : A + \zh{比} + B + \zh{更/还} + adjectif. Ne place pas \zh{更} avant \zh{比}.
\item \zh{不如} remplace la comparaison entière : pas de \zh{比} dans la même phrase.
\item Dans \zh{不但…而且…}, si le sujet est le même, il se met avant \zh{不但} et ne se répète pas.
\item \zh{为了} + but se place \textbf{en début de phrase}, avant le sujet. Jamais à la fin comme en français.
\end{itemize}
\end{attention}

\courssub{Proposition de production et corrigé commenté}

\begin{exoresolu}[Texte modèle : 我的城市最好]
Énoncé : rédige huit à dix phrases pour défendre ta ville contre une ville rivale, en employant les cinq structures de la leçon.
\tcblower
\textbf{Production modèle (équipe Yaoundé) :}\\[4pt]
我的城市最好\\
Wǒ de chéngshì zuì hǎo\\
« Ma ville est la meilleure »\\[4pt]
我住在雅温得，我觉得雅温得比杜阿拉更适合生活。\\
Wǒ zhù zài Yǎwēndé, wǒ juéde Yǎwēndé bǐ Dù'ālā gèng shìhé shēnghuó.\\
« J'habite à Yaoundé et je pense que Yaoundé convient mieux à la vie que Douala. »\\[4pt]
第一，雅温得的天气比杜阿拉凉快一点，不太热。\\
Dì-yī, Yǎwēndé de tiānqì bǐ Dù'ālā liángkuai yìdiǎn, bú tài rè.\\
« Premièrement, le climat de Yaoundé est un peu plus frais que celui de Douala, il n'y fait pas trop chaud. »\\[4pt]
第二，杜阿拉的车比雅温得多多了，路上常常堵车。\\
Dì-èr, Dù'ālā de chē bǐ Yǎwēndé duō duō le, lùshang chángcháng dǔchē.\\
« Deuxièmement, il y a beaucoup plus de voitures à Douala qu'à Yaoundé, et les routes sont souvent bloquées. »\\[4pt]
杜阿拉的海鲜很好吃，可是雅温得的饭更便宜。\\
Dù'ālā de hǎixiān hěn hǎochī, kěshì Yǎwēndé de fàn gèng piányi.\\
« Les fruits de mer de Douala sont très bons, mais les repas de Yaoundé sont encore moins chers. »\\[4pt]
在交通方面，杜阿拉不如雅温得方便。\\
Zài jiāotōng fāngmiàn, Dù'ālā bùrú Yǎwēndé fāngbiàn.\\
« En matière de transport, Douala est moins pratique que Yaoundé. »\\[4pt]
雅温得不但环境好，而且空气也很新鲜。\\
Yǎwēndé búdàn huánjìng hǎo, érqiě kōngqì yě hěn xīnxiān.\\
« Yaoundé a non seulement un bon environnement, mais l'air y est aussi très pur. »\\[4pt]
为了学习中文，我不但去孔子学院上课，而且每天练习说汉语。\\
Wèile xuéxí Zhōngwén, wǒ búdàn qù Kǒngzǐ Xuéyuàn shàngkè, érqiě měi tiān liànxí shuō Hànyǔ.\\
« Pour apprendre le chinois, je vais non seulement en cours à l'Institut Confucius, mais je m'entraîne aussi à parler chinois tous les jours. »\\[4pt]
所以，我觉得我的城市最好！\\
Suǒyǐ, wǒ juéde wǒ de chéngshì zuì hǎo!\\
« C'est pourquoi je pense que ma ville est la meilleure ! »\\[6pt]
\textbf{Commentaire des choix de langue :}
\begin{itemize}
\item La phrase d'ouverture pose la thèse avec \zh{更} + adjectif (\zh{更适合}) : le renforcement annonce d'emblée la position de l'équipe.
\item \zh{凉快一点} nuance un petit écart (le climat), tandis que \zh{多多了} marque un grand écart (les voitures) : les deux degrés de comparaison sont bien distingués.
\item \zh{不如} est utilisé avec un complément de domaine \zh{在交通方面} « en matière de transport », placé en tête de phrase comme en chinois standard : le cadre (temps, lieu, domaine) précède toujours le commentaire.
\item La concession \zh{可是} (« mais ») avant \zh{更便宜} montre qu'on peut reconnaître un avantage adverse puis le dépasser : stratégie argumentative efficace.
\item \zh{不但…而且…} apparaît deux fois avec le même sujet, donc le sujet précède \zh{不但} et n'est pas répété : \zh{雅温得不但…而且…}, \zh{我不但…而且…}.
\item \zh{为了} ouvre la phrase de but ; en chinois, le but exprimé par \zh{为了} se place avant la proposition principale, contrairement au français « pour + infinitif » souvent en fin de phrase.
\item La conclusion reprend le titre avec le superlatif \zh{最}, ce qui boucle élégamment le texte.
\end{itemize}
\end{exoresolu}

\courssub{Bilan de l'activité}

Cette activité t'a permis de passer de la manipulation mécanique des structures à leur usage réel dans un échange argumenté. Vérifie que tu sais désormais :

\begin{itemize}
\item choisir entre \zh{一点} (petit écart) et \zh{多了} (grand écart) selon ce que tu veux dire ;
\item employer \zh{更/还} pour renforcer une comparaison déjà exprimée avec \zh{比} ;
\item utiliser \zh{不如} pour une comparaison négative sans ajouter \zh{比} ;
\item placer correctement \zh{为了} + but en tête de phrase ;
\item enchaîner deux qualités ou deux actions avec \zh{不但…而且…} sans répéter le sujet ;
\item respecter l'ordre des mots chinois : cadre (temps, lieu, domaine) d'abord, commentaire ensuite.
\end{itemize}

\begin{experience}[Pour aller plus loin : le débat s'étend]
Recommence le débat en changeant de sujet : « Le lycée contre la maison » (\zh{学校不如家里舒服吗？}), « Le football contre le basketball » ou « La ville contre le village ». Chaque élève doit utiliser au moins trois structures différentes et poser une question à l'équipe adverse avec \zh{为什么} ou \zh{怎么}. Le jury attribue un point par structure correcte et un point bonus pour toute réaction spontanée bien formulée.
\end{experience}

\begin{savaistu}[Comparer en Chine : une question de politesse]
En Chine, comme au Cameroun, complimenter directement en écrasant l'autre peut sembler impoli. C'est pourquoi la structure \zh{不如} est très appréciée : dire \zh{我不如你} (\emph{Wǒ bùrú nǐ}, « je suis moins bon que toi ») est une formule de modestie très courante, par exemple après un match ou un examen. À l'inverse, \zh{你比我强多了} (\emph{Nǐ bǐ wǒ qiáng duō le}, « tu es bien plus fort que moi ») est un compliment sincère et naturel entre amis. La grammaire de la comparaison est donc aussi un outil de politesse et d'harmonie sociale.
\end{savaistu}

\newpage

\coursec{Méthodes et exercices résolus}

\courssub{Construire correctement les structures étudiées}

\begin{methode}[Construire une phrase comparative en cinq étapes]
Pour construire sans faute une comparaison avec \zh{比} (bǐ), \zh{不如} (bùrú), \zh{更} (gèng) ou \zh{还} (hái), suis toujours la même démarche :
\begin{enumerate}
\item \textbf{Identifier les deux termes comparés} (A et B) et l'adjectif de comparaison (grand, cher, loin, intéressant…).
\item \textbf{Choisir la structure} selon le sens : supériorité simple \zh{A比B + adj.} ; supériorité renforcée \zh{A比B更/还 + adj.} ; infériorité \zh{A不如B (+ adj.)}.
\item \textbf{Placer le cadre comparatif AVANT l'adjectif} : \zh{比B} ou \zh{不如B} se mettent toujours devant l'adjectif, jamais après. Schéma : A + \zh{比} + B + adjectif.
\item \textbf{Nuancer l'écart si nécessaire} : petit écart → adjectif + \zh{一点儿} (yìdiǎnr) ; grand écart → adjectif + \zh{多了} (duō le). Ces deux mots se placent APRÈS l'adjectif.
\item \textbf{Vérifier les interdits} : jamais \zh{很} (hěn), \zh{非常} (fēicháng) ou \zh{太} (tài) devant l'adjectif dans une phrase avec \zh{比} ; utiliser \zh{更} ou \zh{还} à la place.
\end{enumerate}
Réflexe de traduction : « plus… que » → \zh{比} ; « beaucoup plus… que » → \zh{比…多了} ; « un peu plus… que » → \zh{比…一点儿} ; « moins… que / ne pas être aussi… que » → \zh{不如}.
\end{methode}

\begin{center}
\begin{tabularx}{\linewidth}{|X|X|X|}
\hline
\rowcolor{popblueL} Structure & Exemple (caractères + pinyin) & Traduction \\
\hline
A + \zh{比} + B + adj. & \zh{他比我高。} \emph{Tā bǐ wǒ gāo.} & Il est plus grand que moi. \\
\hline
A + \zh{比} + B + adj. + \zh{一点儿} & \zh{他比我高一点儿。} \emph{Tā bǐ wǒ gāo yìdiǎnr.} & Il est un peu plus grand que moi. \\
\hline
A + \zh{比} + B + adj. + \zh{多了} & \zh{他比我高多了。} \emph{Tā bǐ wǒ gāo duō le.} & Il est beaucoup plus grand que moi. \\
\hline
A + \zh{比} + B + \zh{更/还} + adj. & \zh{他比我更高。} \emph{Tā bǐ wǒ gèng gāo.} & Il est encore plus grand que moi. \\
\hline
A + \zh{不如} + B (+ adj.) & \zh{我不如他高。} \emph{Wǒ bùrú tā gāo.} & Je ne suis pas aussi grand que lui. \\
\hline
\end{tabularx}
\end{center}

\begin{exoresolu}[Comparer deux villes]
Énoncé — Traduis en chinois : 1. Douala est plus chaude que Buea. 2. Yaoundé est beaucoup plus grande que mon village. 3. Le train est un peu plus rapide que le bus. 4. Buea est moins chère que Douala.
\tcblower
Solution détaillée :
\begin{enumerate}
\item Étape 1 : A = Douala, B = Buea, adjectif = \zh{热} (rè, chaud). Étape 2 : supériorité simple → \zh{比}. Étape 3 : cadre avant l'adjectif. Réponse : \zh{杜阿拉比布埃亚热。} \emph{Dù'ālā bǐ Bù'āiyà rè.} « Douala est plus chaude que Buea. »
\item « beaucoup plus grande » → grand écart → \zh{多了} après l'adjectif \zh{大} (dà). Réponse : \zh{雅温得比我的村子大多了。} \emph{Yǎwēndé bǐ wǒ de cūnzi dà duō le.} « Yaoundé est beaucoup plus grande que mon village. »
\item « un peu plus rapide » → petit écart → \zh{一点儿} après \zh{快} (kuài). Réponse : \zh{火车比公共汽车快一点儿。} \emph{Huǒchē bǐ gōnggòng qìchē kuài yìdiǎnr.} « Le train est un peu plus rapide que le bus. »
\item « moins chère que » → infériorité → \zh{不如} + adjectif \zh{贵} (guì, cher). Réponse : \zh{布埃亚不如杜阿拉贵。} \emph{Bù'āiyà bùrú Dù'ālā guì.} « Buea n'est pas aussi chère que Douala. »
\end{enumerate}
Conclusion : le choix de la structure dépend du sens (supériorité, infériorité, ampleur de l'écart) ; la place des mots est fixe : \zh{比B} / \zh{不如B} avant l'adjectif, \zh{一点儿} / \zh{多了} après.
\end{exoresolu}

\begin{methode}[Nuancer ou renforcer une comparaison déjà construite]
Pour transformer une comparaison simple en comparaison nuancée ou renforcée :
\begin{enumerate}
\item Repérer la phrase de base : A + \zh{比} + B + adjectif.
\item Pour un \textbf{petit écart}, ajouter \zh{一点儿} juste après l'adjectif : \zh{他比我高一点儿。} \emph{Tā bǐ wǒ gāo yìdiǎnr.} « Il est un peu plus grand que moi. »
\item Pour un \textbf{grand écart}, ajouter \zh{多了} (ou \zh{得多}) après l'adjectif : \zh{他比我高多了。} \emph{Tā bǐ wǒ gāo duō le.} « Il est beaucoup plus grand que moi. »
\item Pour \textbf{renforcer le degré}, insérer \zh{更} (gèng) ou \zh{还} (hái) entre \zh{比B} et l'adjectif : \zh{他比我更高。} \emph{Tā bǐ wǒ gèng gāo.} « Il est encore plus grand que moi. »
\item Attention : on peut dire \zh{比B更/还 + adj.}, mais on ne combine pas \zh{更/还} avec \zh{一点儿/多了} dans la même phrase au niveau HSK 3.
\end{enumerate}
\end{methode}

\begin{exoresolu}[Transformation de phrases]
Énoncé — Réécris chaque phrase selon la consigne entre parenthèses. 1. \zh{今天比昨天热。} (petit écart) 2. \zh{汉语比英语难。} (grand écart) 3. \zh{这本书比那本书有意思。} (renforcer avec \zh{更})
\tcblower
Solution détaillée :
\begin{enumerate}
\item Phrase de base : \zh{今天比昨天热。} \emph{Jīntiān bǐ zuótiān rè.} « Aujourd'hui il fait plus chaud qu'hier. » Petit écart → ajouter \zh{一点儿} après l'adjectif \zh{热} : \zh{今天比昨天热一点儿。} \emph{Jīntiān bǐ zuótiān rè yìdiǎnr.} « Aujourd'hui il fait un peu plus chaud qu'hier. »
\item Grand écart → ajouter \zh{多了} après \zh{难} (nán, difficile) : \zh{汉语比英语难多了。} \emph{Hànyǔ bǐ Yīngyǔ nán duō le.} « Le chinois est beaucoup plus difficile que l'anglais. »
\item Renforcer → insérer \zh{更} entre \zh{比那本书} et l'adjectif \zh{有意思} (yǒu yìsi, intéressant) : \zh{这本书比那本书更有意思。} \emph{Zhè běn shū bǐ nà běn shū gèng yǒu yìsi.} « Ce livre est encore plus intéressant que celui-là. »
\end{enumerate}
Conclusion : les trois modifications se font à des places différentes — \zh{更/还} avant l'adjectif, \zh{一点儿/多了} après ; ne jamais déplacer \zh{比B}.
\end{exoresolu}

\courssub{Éviter les erreurs fréquentes (ordre des mots, mots-outils)}

\begin{methode}[Éviter les pièges de la comparaison négative 不如]
\begin{enumerate}
\item Comprendre le sens : \zh{A不如B} signifie « A n'est pas aussi… que B », c'est-à-dire A < B. Schéma : A + \zh{不如} + B (+ adjectif).
\item Si l'adjectif est exprimé, il se place après B : \zh{我不如他高。} \emph{Wǒ bùrú tā gāo.} « Je ne suis pas aussi grand que lui. »
\item Ne jamais écrire \zh{不比} pour « moins… que » : \zh{他不比我高} signifie « il n'est pas plus grand que moi » (nous sommes peut-être égaux), ce n'est pas la même chose que \zh{他不如我高}.
\item Ne pas ajouter de négation supplémentaire : \zh{不如} contient déjà la négation ; on ne dit jamais \zh{不如…不高}.
\item Pour traduire « moins… que », préférer toujours \zh{不如} : « Le bus est moins rapide que le train » → \zh{公共汽车不如火车快。} \emph{Gōnggòng qìchē bùrú huǒchē kuài.}
\end{enumerate}
\end{methode}

\begin{exoresolu}[Phrases justes ou fausses]
Énoncé — Corrige les phrases fautives et justifie. 1. \zh{我弟弟不如我一点儿高。} 2. \zh{今天比昨天很冷。} 3. \zh{我的房间不如他的房间干净。}
\tcblower
Solution détaillée :
\begin{enumerate}
\item \zh{我弟弟不如我一点儿高。} est FAUSSE. \zh{一点儿} ne se combine pas avec \zh{不如} : la nuance d'écart avec \zh{一点儿/多了} ne fonctionne que dans les phrases en \zh{比}. Correction : \zh{我弟弟不如我高。} \emph{Wǒ dìdi bùrú wǒ gāo.} « Mon petit frère n'est pas aussi grand que moi. »
\item \zh{今天比昨天很冷。} est FAUSSE. Dans une phrase comparative avec \zh{比}, l'adjectif ne peut pas être précédé de \zh{很} (ni de \zh{非常}, \zh{太}). Correction : \zh{今天比昨天冷。} \emph{Jīntiān bǐ zuótiān lěng.} « Aujourd'hui il fait plus froid qu'hier. » Ou, pour marquer le degré : \zh{今天比昨天更冷。} \emph{Jīntiān bǐ zuótiān gèng lěng.}
\item \zh{我的房间不如他的房间干净。} est JUSTE : structure A + \zh{不如} + B + adjectif (\zh{干净}, gānjìng, propre) parfaitement respectée. \emph{Wǒ de fángjiān bùrú tā de fángjiān gānjìng.} « Ma chambre n'est pas aussi propre que la sienne. »
\end{enumerate}
Conclusion : deux interdits à mémoriser — pas de \zh{很} avec \zh{比}, pas de \zh{一点儿} avec \zh{不如}.
\end{exoresolu}

\begin{methode}[Exprimer le but avec 为了]
\begin{enumerate}
\item Sens : \zh{为了} (wèile) = « pour, afin de », il introduit le BUT ou la PERSONNE au profit de qui on agit.
\item Schéma : \zh{为了} + but / personne, + proposition principale. Exemple : \zh{为了健康，他每天运动。} \emph{Wèile jiànkāng, tā měitiān yùndòng.} « Pour sa santé, il fait du sport tous les jours. »
\item Le groupe \zh{为了…} se place en tête de phrase, suivi d'une virgule ; il peut aussi se placer après le sujet : \zh{他为了学好汉语很努力。} \emph{Tā wèile xuéhǎo Hànyǔ hěn nǔlì.} « Il travaille dur pour bien apprendre le chinois. »
\item Attention aux tons : \zh{为了} se prononce wèile (quatrième ton + ton léger) ; ne pas confondre avec \zh{因为} (yīnwèi, parce que) qui exprime la CAUSE, pas le but.
\item Distinction : cause → \zh{因为…所以…} ; but → \zh{为了…}. « Il étudie parce qu'il aime la Chine » ≠ « Il étudie pour aller en Chine ».
\end{enumerate}
\end{methode}

\begin{exoresolu}[Choisir entre 因为 et 为了]
Énoncé — Complète avec \zh{因为} ou \zh{为了}, puis traduis : 1. \zh{……身体好，他每天跑步。} 2. \zh{……下雨了，我们不去市场了。} 3. \zh{……通过考试，她学习很努力。}
\tcblower
Solution détaillée :
\begin{enumerate}
\item « ……身体好，他每天跑步。 » : « avoir une bonne santé » est le BUT recherché, pas la cause → \zh{为了}. \zh{为了身体好，他每天跑步。} \emph{Wèile shēntǐ hǎo, tā měitiān pǎobù.} « Pour être en bonne santé, il court tous les jours. »
\item « ……下雨了，我们不去市场了。 » : la pluie est la CAUSE de l'annulation → \zh{因为}. \zh{因为下雨了，我们不去市场了。} \emph{Yīnwèi xià yǔ le, wǒmen bú qù shìchǎng le.} « Parce qu'il a plu, nous n'allons plus au marché. »
\item « ……通过考试，她学习很努力。 » : réussir l'examen est le BUT de ses efforts → \zh{为了}. \zh{为了通过考试，她学习很努力。} \emph{Wèile tōngguò kǎoshì, tā xuéxí hěn nǔlì.} « Pour réussir l'examen, elle étudie très dur. »
\end{enumerate}
Conclusion : poser la question « est-ce la raison (cause) ou l'objectif (but) ? » : cause → \zh{因为}, but → \zh{为了}.
\end{exoresolu}

\courssub{Utiliser les structures dans des phrases et de courts textes}

\begin{methode}[Additionner les idées avec 不但…而且… et rédiger un court texte]
\begin{enumerate}
\item Sens : \zh{不但…而且…} (búdàn… érqiě…) = « non seulement… mais aussi… ». Schéma : Sujet + \zh{不但} + verbe/adjectif 1, \zh{而且} + verbe/adjectif 2.
\item Si les deux parties ont le MÊME sujet, le sujet se met avant \zh{不但} : \zh{他不但会说汉语，而且会说英语。} \emph{Tā búdàn huì shuō Hànyǔ, érqiě huì shuō Yīngyǔ.} « Non seulement il sait parler chinois, mais il sait aussi parler anglais. »
\item Si les sujets sont DIFFÉRENTS, chaque sujet suit son connecteur : \zh{不但他会说汉语，而且他妹妹也会说。} \emph{Búdàn tā huì shuō Hànyǔ, érqiě tā mèimei yě huì shuō.} « Non seulement lui sait parler chinois, mais sa petite sœur aussi. »
\item Ne jamais remplacer \zh{而且} par \zh{也} seul, ni oublier la deuxième partie : \zh{不但} appelle toujours \zh{而且} (ou \zh{还/也}).
\item Pour rédiger un court texte (par exemple sur l'environnement ou la santé), combiner les structures de la leçon : une comparaison en \zh{比}, une nuance en \zh{一点儿/多了}, une infériorité en \zh{不如}, un but en \zh{为了}, une addition en \zh{不但…而且…}.
\end{enumerate}
\end{methode}

\begin{exoresolu}[Composition guidée : la santé à Yaoundé]
Énoncé — Rédige un court texte de 4 à 6 phrases sur le thème « le sport et la santé » en utilisant au moins une fois chacune des structures : \zh{比…多了}, \zh{不如}, \zh{更}, \zh{为了}, \zh{不但…而且…}.
\tcblower
Solution détaillée — Raisonnement : on construit d'abord une phrase par structure, puis on les enchaîne logiquement (situation → comparaison → but → conclusion).
\begin{enumerate}
\item Avec \zh{不但…而且…} : \zh{运动不但对身体好，而且让人高兴。} \emph{Yùndòng búdàn duì shēntǐ hǎo, érqiě ràng rén gāoxìng.} « Le sport est non seulement bon pour la santé, mais il rend aussi heureux. » (Même sujet \zh{运动} placé avant \zh{不但}.)
\item Avec \zh{比…多了} : \zh{现在运动的人比以前多了。} \emph{Xiànzài yùndòng de rén bǐ yǐqián duō le.} « Maintenant, il y a beaucoup plus de gens qui font du sport qu'avant. » (\zh{多了} après l'adjectif \zh{多}.)
\item Avec \zh{不如} : \zh{跑步不如游泳有意思。} \emph{Pǎobù bùrú yóuyǒng yǒu yìsi.} « La course n'est pas aussi intéressante que la natation. »
\item Avec \zh{更} : \zh{但是为了健康，跑步比游泳更方便。} \emph{Dànshì wèile jiànkāng, pǎobù bǐ yóuyǒng gèng fāngbiàn.} « Mais pour la santé, la course est encore plus pratique que la natation. » (\zh{更} entre \zh{比游泳} et \zh{方便}.)
\item Conclusion avec \zh{为了} : \zh{为了身体好，我们每天都应该运动。} \emph{Wèile shēntǐ hǎo, wǒmen měitiān dōu yīnggāi yùndòng.} « Pour être en bonne santé, nous devrions faire du sport tous les jours. »
\end{enumerate}
Texte rédigé : \zh{运动不但对身体好，而且让人高兴。现在运动的人比以前多了。跑步不如游泳有意思，但是为了健康，跑步比游泳更方便。为了身体好，我们每天都应该运动。}
Conclusion : un bon texte à l'examen aligne des structures variées mais correctes ; chaque structure est vérifiée séparément (place de \zh{比B}, de \zh{更}, de \zh{多了}, de \zh{为了}) avant l'assemblage final.
\end{exoresolu}

\begin{attention}[Les 5 erreurs qui coûtent des points à l'examen]
\begin{enumerate}
\item \textbf{Calque du français « plus… que »} : écrire \zh{他高比我} au lieu de \zh{他比我高。} En chinois, le cadre \zh{比B} précède TOUJOURS l'adjectif : A + \zh{比} + B + adjectif.
\item \textbf{Ajouter \zh{很} dans une comparaison} : \zh{今天比昨天很热。} est fautif. Avec \zh{比}, l'adjectif seul suffit ; pour insister, utiliser \zh{更} ou \zh{还} : \zh{今天比昨天更热。}
\item \textbf{Mal placer \zh{一点儿} et \zh{多了}} : ils se mettent APRÈS l'adjectif (\zh{贵一点儿}, \zh{大多了}), jamais avant, et jamais dans une phrase en \zh{不如}.
\item \textbf{Confondre \zh{为了} (but) et \zh{因为} (cause)} : \zh{为了下雨，我们不去。} est un non-sens. Cause → \zh{因为} ; objectif → \zh{为了}. Attention aussi aux tons : wèile (quatrième ton) contre yīnwèi.
\item \textbf{Casser le couple \zh{不但…而且…}} : oublier \zh{而且}, ou mal placer les sujets. Même sujet → sujet avant \zh{不但} ; sujets différents → sujet après chaque connecteur. Erreur typique : \zh{他不但会说汉语，会说英语。} (il manque \zh{而且}).
\end{enumerate}
\end{attention}

\newpage

\coursec{Exercices}

\courssub{Je vérifie mes connaissances}

\begin{exercice}[QCM : la bonne structure]
Entoure la bonne réponse.
\begin{enumerate}
\item Pour dire « A est un peu plus grand que B », on dit :
\begin{enumerate}
\item A 比 B 大一点 \hspace{1cm} b) A 比 B 大多了 \hspace{1cm} c) A 不如 B 大
\end{enumerate}
\item \zh{为了} sert à exprimer :
\begin{enumerate}
\item la cause \hspace{1cm} b) le but \hspace{1cm} c) la conséquence
\end{enumerate}
\item « Non seulement\ldots{} mais aussi\ldots{} » se traduit par :
\begin{enumerate}
\item \zh{因为\ldots{}所以\ldots{}} \hspace{1cm} b) \zh{虽然\ldots{}但是\ldots{}} \hspace{1cm} c) \zh{不但\ldots{}而且\ldots{}}
\end{enumerate}
\item \zh{他不如我高。} signifie :
\begin{enumerate}
\item Il est plus grand que moi. \hspace{1cm} b) Il est moins grand que moi. \hspace{1cm} c) Il est aussi grand que moi.
\end{enumerate}
\end{enumerate}
\end{exercice}

\begin{exercice}[Vrai ou faux ? Justifie]
Dis si la phrase est correcte ; si elle est fausse, explique pourquoi en français.
\begin{enumerate}
\item \zh{雅温得比杜阿拉更安静。} \emph{Yǎwēndé bǐ Dù'ālā gèng ānjìng.}
\item \zh{我比他很高。} \emph{Wǒ bǐ tā hěn gāo.}
\item \zh{为了保护环境，我们种树。} \emph{Wèile bǎohù huánjìng, wǒmen zhòng shù.}
\item \zh{他不但会说汉语，而且会说英语。} \emph{Tā bùdàn huì shuō Hànyǔ, érqiě huì shuō Yīngyǔ.}
\item \zh{今天比昨天热一点儿。} \emph{Jīntiān bǐ zuótiān rè yìdiǎnr.}
\end{enumerate}
\end{exercice}

\begin{exercice}[Vocabulaire : complète le tableau]
Complète le tableau avec le pinyin et la traduction française.

\begin{center}
\begin{tabularx}{\linewidth}{|X|X|X|X|}
\hline
\rowcolor{popblueL} Caractères & Pinyin & Nature & Traduction \\
\hline
不如 & & prép./v. & \\
\hline
为了 & & prép. & \\
\hline
更 & & adv. & \\
\hline
环境 & & n. & \\
\hline
健康 & & n./adj. & \\
\hline
吵 & & adj. & \\
\hline
\end{tabularx}
\end{center}
\end{exercice}

\begin{exercice}[Associe caractères et pinyin]
Relie chaque mot chinois à son pinyin.
\begin{enumerate}
\item \zh{不但} \hspace{2cm} a) \emph{érqiě}
\item \zh{而且} \hspace{2cm} b) \emph{bùdàn}
\item \zh{比较} \hspace{2cm} c) \emph{wèile}
\item \zh{为了} \hspace{2cm} d) \emph{bǐjiào}
\item \zh{干净} \hspace{2cm} e) \emph{gānjìng}
\end{enumerate}
\end{exercice}

\courssub{Je m'entraîne}

\begin{exercice}[Phrases à trous]
Complète avec : \zh{比、不如、更、还、为了、因为、不但、而且} (certains mots peuvent servir deux fois).
\begin{enumerate}
\item 杜阿拉 \underline{\hspace{1cm}} 雅温得热多了。
\item 我 \underline{\hspace{1cm}} 他高，他一米八，我一米七。
\item 这个公园 \underline{\hspace{1cm}} 那个公园 \underline{\hspace{1cm}} 漂亮。
\item \underline{\hspace{1cm}} 健康，他每天跑步。
\item 他 \underline{\hspace{1cm}} 会说汉语，\underline{\hspace{1cm}} 会写汉字。
\item \underline{\hspace{1cm}} 下雨，我们不去市场了。
\item 今天 \underline{\hspace{1cm}} 昨天冷一点儿。
\item 汉语 \underline{\hspace{1cm}} 英语难，但是很有意思。
\item \underline{\hspace{1cm}} 保护环境，我们应该少开车。
\item 这家饭店的菜 \underline{\hspace{1cm}} 那家好吃多了。
\end{enumerate}
\end{exercice}

\begin{exercice}[Transformation : de 比 à 不如]
Réécris chaque phrase avec \zh{不如} sans changer le sens. Attention à l'ordre des termes !
\begin{enumerate}
\item 我比他矮。\emph{Wǒ bǐ tā ǎi.}
\item 今天比昨天热。\emph{Jīntiān bǐ zuótiān rè.}
\item 公共汽车比出租车慢。\emph{Gōnggòng qìchē bǐ chūzūchē màn.}
\item 法语比汉语容易。\emph{Fǎyǔ bǐ Hànyǔ róngyì.}
\item 农村比城市安静。\emph{Nóngcūn bǐ chéngshì ānjìng.}
\end{enumerate}
\end{exercice}

\begin{exercice}[Correction d'erreurs]
Chaque phrase contient une faute typique des francophones. Trouve-la et réécris la phrase correctement.
\begin{enumerate}
\item 他比我很高。
\item 为了他很累，他休息。
\item 我不但学习汉语，我而且学习英语。
\item 雅温得不如杜阿拉很大。
\item 她比姐姐更很漂亮。
\end{enumerate}
\end{exercice}

\begin{exercice}[Remise en ordre et classement]
\textbf{A.} Remets les mots dans l'ordre pour faire une phrase correcte.
\begin{enumerate}
\item 比 / 一点儿 / 杜阿拉 / 雅温得 / 大
\item 为了 / 我们 / 身体 / 运动 / 每天 / 做
\item 不但 / 而且 / 他 / 聪明 / 努力 / 很 / 也
\item 不如 / 我 / 哥哥 / 高 / 的
\end{enumerate}
\textbf{B.} Voici trois villes : 北京（人口很多）、上海（人口多）、雅温得（人口不太多）。Écris trois phrases pour les comparer : une avec \zh{比\ldots{}多了}, une avec \zh{更}, une avec \zh{不如}.
\end{exercice}

\courssub{Je me prépare au Bac}

\begin{exercice}[Compréhension de texte]
Lis le texte suivant, puis réponds aux questions en français.

\begin{textebox}[为了我们的城市]
为了健康，喀麦隆的城市应该更干净。\\
\emph{Wèile jiànkāng, Kāmàilóng de chéngshì yīnggāi gèng gānjìng.}\\
杜阿拉比雅温得热多了，也比雅温得大一点儿。\\
\emph{Dù'ālā bǐ Yǎwēndé rè duō le, yě bǐ Yǎwēndé dà yìdiǎnr.}\\
但是雅温得不如杜阿拉吵，那里的空气也比较好。\\
\emph{Dànshì Yǎwēndé bùrú Dù'ālā chǎo, nàli de kōngqì yě bǐjiào hǎo.}\\
为了保护环境，我们不但应该少开车，而且应该多走路、多种树。\\
\emph{Wèile bǎohù huánjìng, wǒmen bùdàn yīnggāi shǎo kāi chē, érqiě yīnggāi duō zǒulù, duō zhòng shù.}\\
城市干净了，我们的身体就更健康了。\\
\emph{Chéngshì gānjìng le, wǒmen de shēntǐ jiù gèng jiànkāng le.}
\end{textebox}

\begin{enumerate}
\item Quelle ville est la plus chaude ? La plus grande ? Justifie avec le texte.
\item Que signifie \zh{雅温得不如杜阿拉吵} ? Quelle structure grammaticale est utilisée ?
\item Releve dans le texte une phrase avec \zh{为了} et traduis-la.
\item Releve la phrase avec \zh{不但\ldots{}而且\ldots{}} et traduis-la.
\item Langue : transforme \zh{杜阿拉比雅温得热多了} en phrase avec \zh{不如} sans changer le sens.
\end{enumerate}
\end{exercice}

\begin{exercice}[Thème et version]
\textbf{A. Thème.} Traduis en chinois (caractères) les phrases suivantes.
\begin{enumerate}
\item Pour protéger l'environnement, Douala est non seulement beaucoup plus sale que Yaoundé, mais aussi encore plus bruyante.
\item Mon frère est un peu plus grand que moi ; je suis moins grand que lui.
\item Pour la santé, nous devons manger moins de riz et plus de légumes.
\end{enumerate}
\textbf{B. Version.} Traduis en français.
\begin{enumerate}
\item 为了学习好，他不但上课认真听，而且每天复习。\emph{Wèile xuéxí hǎo, tā bùdàn shàngkè rènzhēn tīng, érqiě měitiān fùxí.}
\item 公共汽车不如出租车快，但是便宜多了。\emph{Gōnggòng qìchē bùrú chūzūchē kuài, dànshì piányi duō le.}
\end{enumerate}
\end{exercice}

\begin{exercice}[Expression écrite : 为了健康]
Sujet : \zh{为了健康} « Pour la santé ».

Rédige en chinois un court texte de \textbf{80 à 100 caractères} dans lequel tu expliques ce que tu fais (ou ce que tu devrais faire) pour être en bonne santé.

Consignes obligatoires :
\begin{itemize}
\item utilise au moins une fois \zh{为了} ;
\item utilise au moins une fois \zh{不但\ldots{}而且\ldots{}} ;
\item utilise au moins une comparaison avec \zh{比} ou \zh{不如} ;
\item rédige en caractères chinois, avec des phrases complètes et correctes.
\end{itemize}

Piste : tu peux parler du sport, de la nourriture, du sommeil, et comparer ta vie d'avant et ta vie d'aujourd'hui.
\end{exercice}

\newpage

\coursec{Corrigés des exercices}

\begin{corrige}[Exercice 1 — QCM : la bonne structure]
\begin{enumerate}
\item \textbf{a) A 比 B 大一点.} \zh{一点} (ou \zh{一点儿}) après l'adjectif marque un \cle{petit écart} (« un peu plus ») ; \zh{多了} marque un grand écart ; \zh{不如} exprime l'infériorité, pas la supériorité.
\item \textbf{b) le but.} \zh{为了} + nom/proposition introduit le \cle{but} (« pour, afin de »). La cause s'exprime avec \zh{因为}.
\item \textbf{c) 不但……而且……} « non seulement\ldots{} mais aussi\ldots{} ». \zh{因为……所以……} = cause/conséquence ; \zh{虽然……但是……} = concession (« bien que\ldots{} mais »).
\item \textbf{b) Il est moins grand que moi.} \zh{A 不如 B + adj.} signifie « A n'atteint pas B, A est moins\ldots{} que B ».
\end{enumerate}
\end{corrige}

\begin{corrige}[Exercice 2 — Vrai ou faux ? Justifie]
\begin{enumerate}
\item \textbf{Vrai.} \zh{雅温得比杜阿拉更安静。} Structure correcte : A + \zh{比} + B + \zh{更} + adjectif. \zh{更} renforce la comparaison (« encore plus »).
\item \textbf{Faux.} \zh{我比他很高。} est incorrect : dans une comparaison avec \zh{比}, l'adjectif ne peut pas être précédé de \zh{很} (ni de \zh{非常}). On dit : \zh{我比他高。} \emph{Wǒ bǐ tā gāo.} « Je suis plus grand que lui. »
\item \textbf{Vrai.} \zh{为了保护环境，我们种树。} \zh{为了} + proposition de but placée en tête de phrase, suivie d'une virgule : « Pour protéger l'environnement, nous plantons des arbres. »
\item \textbf{Vrai.} \zh{他不但会说汉语，而且会说英语。} \zh{不但} se place après le sujet (même sujet dans les deux propositions), \zh{而且} introduit le second élément : « Il sait non seulement parler chinois, mais aussi anglais. »
\item \textbf{Vrai.} \zh{今天比昨天热一点儿。} A + \zh{比} + B + adjectif + \zh{一点儿} : « Aujourd'hui, il fait un peu plus chaud qu'hier. »
\end{enumerate}
\end{corrige}

\begin{corrige}[Exercice 3 — Vocabulaire : complète le tableau]
\begin{center}
\begin{tabularx}{\linewidth}{|X|X|X|X|}
\hline
\rowcolor{popblueL} Caractères & Pinyin & Nature & Traduction \\
\hline
不如 & bùrú & prép./v. & ne pas égaler, être moins\ldots{} que \\
\hline
为了 & wèile & prép. & pour, afin de (but) \\
\hline
更 & gèng & adv. & encore plus, davantage \\
\hline
环境 & huánjìng & n. & environnement, milieu \\
\hline
健康 & jiànkāng & n./adj. & santé ; sain, en bonne santé \\
\hline
吵 & chǎo & adj. & bruyant \\
\hline
\end{tabularx}
\end{center}
\end{corrige}

\begin{corrige}[Exercice 4 — Associe caractères et pinyin]
\begin{enumerate}
\item \zh{不但} → \textbf{b)} \emph{bùdàn}
\item \zh{而且} → \textbf{a)} \emph{érqiě}
\item \zh{比较} → \textbf{d)} \emph{bǐjiào}
\item \zh{为了} → \textbf{c)} \emph{wèile}
\item \zh{干净} → \textbf{e)} \emph{gānjìng}
\end{enumerate}
\end{corrige}

\begin{corrige}[Exercice 5 — Phrases à trous]
\begin{enumerate}
\item 杜阿拉 \zh{比} 雅温得热多了。\emph{Dù'ālā bǐ Yǎwēndé rè duō le.} « Douala est beaucoup plus chaude que Yaoundé. » (\zh{多了} = grand écart.)
\item 我 \zh{不如} 他高，他一米八，我一米七。\emph{Wǒ bùrú tā gāo.} « Je suis moins grand que lui. » (1,70 m < 1,80 m : infériorité.)
\item 这个公园 \zh{比} 那个公园 \zh{更} 漂亮。\emph{Zhège gōngyuán bǐ nàge gōngyuán gèng piàoliang.} « Ce parc est encore plus beau que ce parc-là. »
\item \zh{为了} 健康，他每天跑步。\emph{Wèile jiànkāng, tā měitiān pǎobù.} « Pour sa santé, il court tous les jours. »
\item 他 \zh{不但} 会说汉语，\zh{而且} 会写汉字。\emph{Tā bùdàn huì shuō Hànyǔ, érqiě huì xiě Hànzì.} « Non seulement il sait parler chinois, mais il sait aussi écrire les caractères. »
\item \zh{因为} 下雨，我们不去市场了。\emph{Yīnwèi xià yǔ, wǒmen bú qù shìchǎng le.} « Parce qu'il pleut, nous n'allons plus au marché. » (cause, pas but.)
\item 今天 \zh{比} 昨天冷一点儿。\emph{Jīntiān bǐ zuótiān lěng yìdiǎnr.} « Aujourd'hui, il fait un peu plus froid qu'hier. »
\item 汉语 \zh{比} 英语难，但是很有意思。\emph{Hànyǔ bǐ Yīngyǔ nán, dànshì hěn yǒu yìsi.} « Le chinois est plus difficile que l'anglais, mais très intéressant. »
\item \zh{为了} 保护环境，我们应该少开车。\emph{Wèile bǎohù huánjìng, wǒmen yīnggāi shǎo kāi chē.} « Pour protéger l'environnement, nous devrions moins conduire. »
\item 这家饭店的菜 \zh{比} 那家好吃多了。\emph{Zhè jiā fàndiàn de cài bǐ nà jiā hǎochī duō le.} « Les plats de ce restaurant sont bien meilleurs que ceux de l'autre. »
\end{enumerate}
\end{corrige}

\begin{corrige}[Exercice 6 — Transformation : de 比 à 不如]
Méthode : \zh{A 比 B + adj.} devient \zh{B 不如 A + adj.} — on \cle{inverse l'ordre des termes} et on garde le même adjectif.
\begin{enumerate}
\item \zh{他不如我矮。} \emph{Tā bùrú wǒ ǎi.} « Il n'est pas aussi petit que moi. » (il est plus grand que moi ; phrase d'origine : \zh{我比他矮}.)
\item \zh{昨天不如今天热。} \emph{Zuótiān bùrú jīntiān rè.} « Hier, il faisait moins chaud qu'aujourd'hui. »
\item \zh{出租车不如公共汽车慢。} \emph{Chūzūchē bùrú gōnggòng qìchē màn.} « Le taxi est moins lent que le bus. »
\item \zh{汉语不如法语容易。} \emph{Hànyǔ bùrú Fǎyǔ róngyì.} « Le chinois est moins facile que le français. »
\item \zh{城市不如农村安静。} \emph{Chéngshì bùrú nóngcūn ānjìng.} « La ville est moins calme que la campagne. »
\end{enumerate}
\end{corrige}

\begin{corrige}[Exercice 7 — Correction d'erreurs]
\begin{enumerate}
\item Faute : \zh{很} devant l'adjectif dans une comparaison. Correction : \zh{他比我高。} \emph{Tā bǐ wǒ gāo.} « Il est plus grand que moi. » (ou \zh{他比我高多了。})
\item Faute : \zh{为了} exprime le \cle{but}, pas la cause ; ici il faut \zh{因为}. Correction : \zh{因为他很累，所以他休息。} \emph{Yīnwèi tā hěn lèi, suǒyǐ tā xiūxi.} « Parce qu'il est fatigué, il se repose. »
\item Faute : \zh{而且} ne se place pas après le sujet répété ; avec le même sujet, on dit \zh{不但……而且……} autour des prédicats. Correction : \zh{我不但学习汉语，而且学习英语。} \emph{Wǒ bùdàn xuéxí Hànyǔ, érqiě xuéxí Yīngyǔ.}
\item Faute : \zh{很} après \zh{不如}. Correction : \zh{雅温得不如杜阿拉大。} \emph{Yǎwēndé bùrú Dù'ālā dà.} « Yaoundé est moins grande que Douala. »
\item Faute : double marqueur d'intensité \zh{更} + \zh{很}. Correction : \zh{她比姐姐更漂亮。} \emph{Tā bǐ jiějie gèng piàoliang.} « Elle est encore plus jolie que sa grande sœur. »
\end{enumerate}
\end{corrige}

\begin{corrige}[Exercice 8 — Remise en ordre et classement]
\textbf{A. Remise en ordre.}
\begin{enumerate}
\item \zh{杜阿拉比雅温得大一点儿。} \emph{Dù'ālā bǐ Yǎwēndé dà yìdiǎnr.} « Douala est un peu plus grande que Yaoundé. »
\item \zh{为了身体，我们每天做运动。} \emph{Wèile shēntǐ, wǒmen měitiān zuò yùndòng.} « Pour notre corps (notre santé), nous faisons du sport tous les jours. »
\item \zh{他不但很聪明，而且也很努力。} \emph{Tā bùdàn hěn cōngming, érqiě yě hěn nǔlì.} « Non seulement il est intelligent, mais il est aussi travailleur. »
\item \zh{我不如我的哥哥高。} \emph{Wǒ bùrú wǒ de gēge gāo.} « Je suis moins grand que mon grand frère. »
\end{enumerate}
\textbf{B. Comparaison des trois villes (propositions).}
\begin{enumerate}
\item \zh{北京的人口比上海多多了。} \emph{Běijīng de rénkǒu bǐ Shànghǎi duō duō le.} « La population de Pékin est bien plus grande que celle de Shanghai. » (on accepte aussi \zh{多了})
\item \zh{上海的人口比雅温得更多。} \emph{Shànghǎi de rénkǒu bǐ Yǎwēndé gèng duō.} « La population de Shanghai est encore plus grande que celle de Yaoundé. »
\item \zh{雅温得的人口不如上海多。} \emph{Yǎwēndé de rénkǒu bùrú Shànghǎi duō.} « La population de Yaoundé est moins grande que celle de Shanghai. »
\end{enumerate}
\end{corrige}

\begin{corrige}[Exercice 9 — Compréhension de texte]
\begin{enumerate}
\item La ville la plus chaude est \textbf{Douala} : \zh{杜阿拉比雅温得热多了} « Douala est beaucoup plus chaude que Yaoundé ». La plus grande est aussi \textbf{Douala} : \zh{也比雅温得大一点儿} « elle est aussi un peu plus grande que Yaoundé ».
\item \zh{雅温得不如杜阿拉吵} signifie « Yaoundé est moins bruyante que Douala ». C'est la structure de \cle{comparaison d'infériorité} : A + \zh{不如} + B + adjectif.
\item Phrase avec \zh{为了} : \zh{为了保护环境，我们不但应该少开车，而且应该多走路、多种树。} « Pour protéger l'environnement, nous devrions non seulement moins conduire, mais aussi marcher davantage et planter plus d'arbres. » (On accepte aussi la première phrase : \zh{为了健康，喀麦隆的城市应该更干净。} « Pour la santé, les villes du Cameroun devraient être plus propres. »)
\item Phrase avec \zh{不但……而且……} : \zh{我们不但应该少开车，而且应该多走路、多种树。} « Nous devrions non seulement moins prendre la voiture, mais aussi marcher plus et planter plus d'arbres. »
\item Transformation : \zh{雅温得不如杜阿拉热。} \emph{Yǎwēndé bùrú Dù'ālā rè.} « Yaoundé est moins chaude que Douala. » (inversion des termes, adjectif conservé.)
\end{enumerate}
\textbf{Barème indicatif (10 pts)} : 2 pts par question (1 pt pour la réponse, 1 pt pour la justification ou la traduction correcte).
\textbf{Points de vigilance} : ne pas traduire \zh{不如} par « ne pas aimer » ; dans la question 5, bien inverser l'ordre des villes et ne pas ajouter \zh{很} ni \zh{多了} après l'adjectif.
\end{corrige}

\begin{corrige}[Exercice 10 — Thème et version]
\textbf{A. Thème.}
\begin{enumerate}
\item \zh{为了保护环境，杜阿拉不但比雅温得脏多了，而且还更吵。}\\ \emph{Wèile bǎohù huánjìng, Dù'ālā bùdàn bǐ Yǎwēndé zāng duō le, érqiě hái gèng chǎo.} (On accepte \zh{比雅温得更吵} pour « encore plus bruyante ».)
\item \zh{我哥哥比我高一点儿；我不如他高。}\\ \emph{Wǒ gēge bǐ wǒ gāo yìdiǎnr ; wǒ bùrú tā gāo.}
\item \zh{为了健康，我们应该少吃米饭，多吃蔬菜。}\\ \emph{Wèile jiànkāng, wǒmen yīnggāi shǎo chī mǐfàn, duō chī shūcài.}
\end{enumerate}
\textbf{B. Version.}
\begin{enumerate}
\item \zh{为了学习好，他不但上课认真听，而且每天复习。} « Pour bien réussir ses études, non seulement il écoute attentivement en classe, mais en plus il révise tous les jours. »
\item \zh{公共汽车不如出租车快，但是便宜多了。} « Le bus est moins rapide que le taxi, mais il est beaucoup moins cher. »
\end{enumerate}
\textbf{Barème indicatif (10 pts)} : 2 pts par phrase (1 pt structure correcte, 1 pt vocabulaire et caractères).
\textbf{Points de vigilance} : pas de \zh{很} dans les comparaisons ; \zh{一点儿}/\zh{多了} se placent \emph{après} l'adjectif ; \zh{为了} en tête de phrase suivi d'une virgule ; dans \zh{不但……而且……}, ne pas répéter le sujet après \zh{而且} quand c'est le même.
\end{corrige}

\begin{corrige}[Exercice 11 — Expression écrite : 为了健康]
\textbf{Proposition de rédaction modèle} (environ 90 caractères) :

\zh{为了健康，我现在每天做运动。以前我不如同学们爱运动，身体也不太好。现在我的身体比以前好多了。我不但每天跑步，而且多吃蔬菜、少吃肉。为了身体更好，我还每天早睡早起。健康很重要，我们都要注意！}

\emph{Wèile jiànkāng, wǒ xiànzài měitiān zuò yùndòng. Yǐqián wǒ bùrú tóngxuémen ài yùndòng, shēntǐ yě bú tài hǎo. Xiànzài wǒ de shēntǐ bǐ yǐqián hǎo duō le. Wǒ bùdàn měitiān pǎobù, érqiě duō chī shūcài, shǎo chī ròu. Wèile shēntǐ gèng hǎo, wǒ hái měitiān zǎo shuì zǎo qǐ. Jiànkāng hěn zhòngyào, wǒmen dōu yào zhùyì!}

« Pour ma santé, je fais maintenant du sport tous les jours. Avant, j'aimais moins le sport que mes camarades et ma santé n'était pas très bonne. Maintenant, ma santé est bien meilleure qu'avant. Non seulement je cours chaque jour, mais je mange aussi plus de légumes et moins de viande. Pour être encore en meilleure santé, je me couche et me lève tôt tous les jours. La santé est très importante, nous devons tous y faire attention ! »

\textbf{Vérification des consignes} :
\begin{itemize}
\item \zh{为了} utilisé deux fois (début et avant-dernière phrase) ;
\item \zh{不但……而且……} présent : \zh{我不但每天跑步，而且多吃蔬菜} ;
\item comparaisons : \zh{我不如同学们爱运动} (\zh{不如}) et \zh{我的身体比以前好多了} (\zh{比} + \zh{多了}) ;
\item environ 90 caractères, phrases complètes.
\end{itemize}
\textbf{Barème indicatif (10 pts)} : respect des 4 consignes obligatoires (4 pts), correction grammaticale et ordre des mots (3 pts), vocabulaire et caractères (2 pts), cohérence du texte (1 pt).
\textbf{Points de vigilance} : ne pas écrire \zh{比以前很身体好} (interdit : \zh{很} avec \zh{比}) ; placer \zh{一点儿}/\zh{多了} après l'adjectif ; éviter de répéter le sujet après \zh{而且} ; compter les caractères pour rester entre 80 et 100.
\end{corrige}

\newpage

\coursec{Fiche bilan}

\courssub{Fiche bilan — Leçon 10 : Comparaisons nuancées, but et addition}

\begin{popfigure}
\begin{tikzpicture}[
  mindmap,
  grow cyclic,
  every node/.style={concept, execute at begin node=\strut},
  concept color=popblue,
  level 1/.append style={level distance=4.5cm, sibling angle=360/7},
  level 2/.append style={level distance=3cm, sibling angle=45},
  text=white,
  font=\small\bfseries
]
\node[concept, scale=1.2, align=center]{Leçon 10\\Comparatif / But / Addition}
  child[concept color=poporange] { node[concept, align=center]{Nuancer l'écart\\A 比 B Adj \textbf{一点/多了}} }
  child[concept color=popgreen] { node[concept, align=center]{Renforcer\\A 比 B \textbf{更/还} Adj} }
  child[concept color=poppurple] { node[concept, align=center]{Infériorité\\A \textbf{不如} B (Adj)} }
  child[concept color=poppink] { node[concept, align=center]{Exprimer le but\\\textbf{为了} + Verbe / Nom} }
  child[concept color=popgold] { node[concept, align=center]{Addition\\\textbf{不但}…\textbf{而且}…} }
  child[concept color=popdark] { node[concept, align=center]{Erreurs fréquentes\\Ordre / 的/得/地 / 了} }
  child[concept color=popblue!70!poppurple] { node[concept, align=center]{Lexique clé\\Environnement / Santé} };
\end{tikzpicture}
\legende{Carte mentale de la leçon 10 — Structures comparatives, finalité et coordination}
\end{popfigure}

\begin{bilan}

\textbf{Lexique clé (HSK 2–3 — Module Environnement \& Santé)}
\begin{center}
\begin{tabularx}{\linewidth}{|X|X|X|X|}
\hline
\rowcolor{popblueL} \textbf{Caractères} & \textbf{Pinyin} & \textbf{Nature} & \textbf{Traduction} \\ \hline
比较 & bǐjiào & adv. & relativement, assez \\
程度 & chéngdù & n. & degré, niveau \\
污染 & wūrǎn & v./n. & polluer / pollution \\
环境 & huánjìng & n. & environnement \\
保护 & bǎohù & v. & protéger \\
健康 & jiànkāng & adj./n. & en bonne santé / santé \\
锻炼 & duànliàn & v. & faire de l'exercice \\
饮食 & yǐnshí & n. & alimentation, régime \\
目的 & mùdì & n. & but, objectif \\
进步 & jìnbù & v./n. & progresser / progrès \\
\hline
\end{tabularx}
\end{center}

\textbf{Règles de grammaire à connaître par cœur}
\begin{center}
\begin{tabularx}{\linewidth}{|X|X|X|}
\hline
\rowcolor{popblueL} \textbf{Structure} & \textbf{Schéma / Exemple} & \textbf{Emploi / Nuance} \\ \hline
\textbf{Écart faible} & A 比 B Adj \textbf{一点儿 / 一点} & « Un peu plus… » — oral, sujetif. \\ \hline
\textbf{Écart fort} & A 比 B Adj \textbf{多了 / 得多} & « Beaucoup plus… » — insistante. \\ \hline
\textbf{Comparatif de supériorité renforcé} & A 比 B \textbf{更 / 还} Adj & « Encore plus / encore plus… » — \textbf{更} = objectif/factuel ; \textbf{还} = subjectif/surprise. \\ \hline
\textbf{Infériorité (A < B)} & A \textbf{不如} B (Adj) & « A n'est pas aussi… que B ». Adj souvent omis si contexte clair. \\ \hline
\textbf{But (afin de / pour)} & \textbf{为了} + [Verbe / Nom] + ，Sujet + Verbe & Place \textbf{avant} le sujet (sujet commun) ou après (sujets diff.). Nég. : \textbf{为了不} / \textbf{免得}. \\ \hline
\textbf{Addition (non seulement… mais aussi)} & Sujet \textbf{不但 / 不仅} … ,\textbf{而且 / 还 / 也} … & Progression d'intensité. Sujet commun : pas de répétition. Sujets diff. : répéter structure. \\ \hline
\end{tabularx}
\end{center}

\textbf{Méthodes express}
\begin{itemize}
  \item \textbf{Comparer :} Identifier A et B $\rightarrow$ Choisir structure (écart ? force ? infériorité ?) $\rightarrow$ Placer Adj \textbf{après} la structure 比/不如 $\rightarrow$ Ajouter particules de degré (一点/多了/更/还).
  \item \textbf{Traduire « pour / afin de » :} Repérer l'action-but $\rightarrow$ Placer \textbf{为了} + Verbe en tête de phrase (si même sujet) ou avant 2e verbe $\rightarrow$ Attention : \textbf{为了} ne s'utilise pas seul devant nom comme « pour » (ex. \zh{这个礼物是为了你} $\neq$ « Ce cadeau est pour toi » $\rightarrow$ \zh{给你的}).
  \item \textbf{Relier deux qualités/actions :} Ordonner de la moins forte à la plus forte $\rightarrow$ Insérer \textbf{不但} devant 1er verbe/adj, \textbf{而且} devant 2e $\rightarrow$ Vérifier accord sujet (commun ou non).
\end{itemize}

\end{bilan}

\begin{aretenir}[Checklist avant le Bac]
$\square$ Je sais écrire l'ordre des traits de \zh{为了}、\zh{不如}、\zh{不但}、\zh{而且} (haut→bas, gauche→droite, horizontal avant vertical).\\
$\square$ Je distingue \zh{一点} (faible) de \zh{多了} (fort) et je les place \textbf{après} l'adjectif dans \zh{A 比 B Adj 一点/多了}.\\
$\square$ Je choisis \zh{更} (factuel) vs \zh{还} (subjectif/insistance) dans \zh{A 比 B 更/还 Adj}.\\
$\square$ Je construis \zh{A 不如 B} sans adjectif si le sens est clair (ex. \zh{我弟弟不如我}).\\
$\square$ Je place \zh{为了} + Verbe \textbf{avant} le sujet principal (sujet commun) ou avant le 2e verbe (sujets différents).\\
$\square$ Je forme la négation de but : \zh{为了不} + Verbe / \zh{免得} + Verbe.\\
$\square$ J'utilise \zh{不但…而且…} avec progression d'intensité (2e idée > 1re).\\
$\square$ Je ne confonds pas \zh{不但} (adverbe) et \zh{不但} (conjonction) — même forme, rôle syntaxique selon place.\\
$\square$ Je traduis « non seulement… mais encore » par \zh{不但…而且…} et non \zh{不但…也…} (incorrect).\\
$\square$ Je repère les pièges : \zh{比} disparaît avec \zh{不如} ; \zh{了} (particule aspectuelle) $\neq$ \zh{多了} (complément de degré).\\
$\square$ Je produis 3 phrases types : environnement (pollution), santé (sport), vie scolaire (devoirs) avec les 5 structures.
\end{aretenir}

\findecours

\end{document}
